% Generated mechanically from frozen input p04/ko/schemes.tex.
% Do not edit this generated file; edit the bound locale source instead.

% OK, start here.
%

\setcounter{chapter}{25}
\chapter{스킴}



\phantomsection
\label{schemes-section-phantom}


\section{서론}
\label{schemes-section-introduction}

\noindent
이 문서에서는 스킴을 정의한다.
기본 참고문헌은 \cite{EGA}이다.









\section{국소 환 달린 공간}
\label{schemes-section-locally-ringed-spaces}

\noindent
Sheaves의 절 \ref{sheaves-section-ringed-spaces}에서
환 달린 공간을 정의했음을 상기하자.
간단히 말해, 환 달린 공간은 쌍 $(X, \mathcal{O}_X)$으로서,
위상공간 $X$와 환의 층 $\mathcal{O}_X$로 이루어진다.
환 달린 공간의 사상 $f : (X, \mathcal{O}_X) \to (Y, \mathcal{O}_Y)$은
연속사상 $f : X \to Y$와 환의 층들의 $f$-사상
$f^\sharp : \mathcal{O}_Y \to \mathcal{O}_X$로 주어진다.
$f^\sharp$를 사상 $\mathcal{O}_Y \to f_*\mathcal{O}_X$로 생각할 수도 있다.
Sheaves의 정의 \ref{sheaves-definition-f-map}과
보조정리 \ref{sheaves-lemma-f-map}을 보라.

\medskip\noindent
염두에 둘 만한 좋은 기하학적 예는 $\mathcal{C}^\infty$-다양체와
$\mathcal{C}^\infty$-다양체의 사상이다. 실제로 $M$이
$\mathcal{C}^\infty$-다양체이면 매끄러운 함수들의 층
$\mathcal{C}^\infty_M$은 $M$ 위의 환의 층이다. 또한 다양체의 사상
$f : M \to N$이 매끄러울 필요충분조건은 모든 국소 단면 $h$, 즉
$\mathcal{C}^\infty_N$의 모든 국소 단면에 대해 합성 $h \circ f$가
$\mathcal{C}^\infty_M$의 국소 단면인 것이다.
따라서 매끄러운 사상 $f$는 자연스럽게 환 달린 공간의 사상
$$
f : (M , \mathcal{C}^\infty_M) \longrightarrow (N, \mathcal{C}^\infty_N)
$$
을 준다. Sheaves의 예 \ref{sheaves-example-continuous-map-ringed}을 보라.
줄기에서 무슨 일이 일어나는지 살펴보면 이해에 도움이 된다. 즉,
$m \in M$이고 그 상이 $f(m) = n \in N$이라고 하자. 줄기
$\mathcal{C}^\infty_{M, m}$은 $m$에서 매끄러운 함수들의 싹으로 이루어진 환이다.
Sheaves의 예 \ref{sheaves-example-germs-functions}을 보라.
$(M, m)$ 위 함수의 싹 대수는 국소환이며, 그 극대 아이디얼은
$m$에서 사라지는 함수들로 이루어진다.
$\mathcal{C}^\infty_{N, n}$도 마찬가지이다. 줄기 위의 사상
$f^\sharp : \mathcal{C}^\infty_{N, n} \to \mathcal{C}^\infty_{M, m}$은
극대 아이디얼을 극대 아이디얼 안으로 보낸다. 이는 단순히
$f(m) = n$이기 때문이다.

\medskip\noindent
대수기하학에서는 스킴을 연구한다. 스킴 위의 환의 층은
그 공간의 내재적 성질만으로 결정되지 않는다.
환 $R$의 스펙트럼
(Algebra의 절 \ref{algebra-section-spectrum-ring}을 보라)에
$R$로부터 구성한 환의 층(아래에서 설명한다)을 갖춘 것이
우리의 기본 구성요소가 된다. $\mathcal{O}$의 $\Spec(R)$ 위 줄기는
$R$을 그 소 아이디얼에서 국소화한 국소환임을 알게 될 것이다.
이 상황에서 국소 환 달린 공간을 도입하는 이유는 두 가지이다.
(1) 일반적으로 스펙트럼 사이의 연속사상에 대응하는 환들의 사상을
부여하는 장치가 없다. 이 때문에 사상 $f^\sharp$를 추가 자료로 둔다.
(2) 이 스펙트럼들의 사상을 환 달린 공간의 범주에서 생각하면
줄기 위의 사상들이 국소 준동형이 아닐 수도 있다. 기하학적 직관에 따르면
국소 준동형이어야 하므로 다음과 같이 국소 환 달린 공간을 도입한다.

\begin{definition}
\label{schemes-definition-locally-ringed-space}
국소 환 달린 공간에 관하여 다음을 정의한다.
\begin{enumerate}
\item {\it 국소 환 달린 공간 $(X, \mathcal{O}_X)$}은 위상공간 $X$와
환의 층 $\mathcal{O}_X$로 이루어진 쌍으로서, 그 모든 줄기가 국소환인 것이다.
\item 국소 환 달린 공간 $(X, \mathcal{O}_X)$이 주어졌을 때
$\mathcal{O}_{X, x}$를 {\it $X$의 $x$에서의 국소환}이라 한다.
$\mathfrak{m}_{X, x}$ 또는 간단히 $\mathfrak{m}_x$로
$\mathcal{O}_{X, x}$의 극대 아이디얼을 나타낸다. 또한
{\it $X$의 $x$에서의 잉여체}는 잉여체
$\kappa(x) = \mathcal{O}_{X, x}/\mathfrak{m}_x$이다.
\item {\it 국소 환 달린 공간의 사상}
$(f, f^\sharp) : (X, \mathcal{O}_X) \to (Y, \mathcal{O}_Y)$은
모든 $x \in X$에 대해 유도된 환 사상
$\mathcal{O}_{Y, f(x)} \to \mathcal{O}_{X, x}$이 국소 준동형인
환 달린 공간의 사상이다.
\end{enumerate}
\end{definition}

\noindent
국소 환 달린 공간을 논할 때에는 대개 표기에서 환의 층
$\mathcal{O}_X$를 생략한다. 단순히 “국소 환 달린 공간 $X$”라고 할 것이다.
또 표기의 남용으로 $X$를 그 바탕 위상공간으로도 생각할 것이다.
마지막으로 이에 대응하는 환의 층
$\mathcal{O}_X$를 {\it $X$의 구조층}이라 한다.
더 나아가 국소환 $\mathcal{O}_{X, x}$의 극대 아이디얼을
$\mathfrak{m}_{X, x}$ 또는 간단히 $\mathfrak{m}_x$로 나타내는 것이 관례이다.
“$f : X \to Y$를 국소 환 달린 공간의 사상이라 하자”라고 말할 때에는
구조층까지 생략한다. 이 경우에도 표기의 남용으로
$f : X\to Y$를 바탕 위상공간들의 연속사상으로 생각한다.
$f$에 대응하는 $f$-사상은 관례상 $f^\sharp$로 나타낸다.
$f$가 국소 환 달린 공간의 사상이라는 조건은 모든 $x\in X$에 대해
줄기 위의 사상
$$
f^\sharp_x : \mathcal{O}_{Y, f(x)} \longrightarrow \mathcal{O}_{X, x}
$$
이 극대 아이디얼 $\mathfrak m_{Y, f(x)}$을
$\mathfrak m_{X, x}$ 안으로 보낸다는 말로 표현할 수 있다.

\medskip\noindent
이 표기 관례를 이용하여 국소 환 달린 공간들과 그 사상들이
범주를 이룸을 보이자. 이를 위해서는 국소 환 달린 공간의 사상들의
합성이 다시 국소 환 달린 공간의 사상임을 보이면 된다. 따라서
$f : X \to Y$와 $g : Y \to Z$를 국소 환 달린 공간의 사상이라 하자.
$f$와 $g$의 합성은 Sheaves의 정의
\ref{sheaves-definition-composition-maps-ringed-spaces}에서 정의되어 있다.
$x \in X$라 하자. Sheaves의 보조정리
\ref{sheaves-lemma-compose-f-maps-stalks}에 의해 합성
$$
\mathcal{O}_{Z, g(f(x))}
\xrightarrow{g^\sharp}
\mathcal{O}_{Y, f(x)}
\xrightarrow{f^\sharp}
\mathcal{O}_{X, x}
$$
은 사상 $g \circ f$에 딸린 줄기 위의 사상이다.
국소환 준동형들의 합성은 국소환 준동형이므로 결론이 따른다.

\medskip\noindent
이 정의의 좋은 성질 하나는 함자
$$
\textit{Locally ringed spaces}
\longrightarrow
\textit{Ringed spaces}
$$
가 동형을 반영한다는 것이다(그보다 더 많은 것을 반영한다).
다음은 이를 덜 추상적으로 진술한 것이다.

\begin{lemma}
\label{schemes-lemma-isomorphism-locally-ringed}
\begin{slogan}
국소 환 달린 공간들 사이의 환 달린 공간의 동형은
국소 환 달린 공간의 동형이다.
\end{slogan}
$X$, $Y$를 국소 환 달린 공간이라 하자.
$f : X \to Y$가 환 달린 공간의 동형이면 $f$는
국소 환 달린 공간의 동형이다.
\end{lemma}

\begin{proof}
이는 대수에서의 대응하는 사실로부터 자명하게 따른다.
$A$, $B$가 국소환이라고 하자. 임의의 환 동형
$A \to B$는 국소환 동형이다.
\end{proof}


\section{국소 환 달린 공간의 열린 몰입}
\label{schemes-section-open-immersion}

\begin{definition}
\label{schemes-definition-immersion-locally-ringed-spaces}
$f : X \to Y$를 국소 환 달린 공간의 사상이라 하자.
$f$를 {
\it 열린 몰입}이라 하는 것은 $f$가 $X$를 $Y$의 열린 부분집합 위로
보내는 위상동형이고 사상 $f^{-1}\mathcal{O}_Y \to \mathcal{O}_X$가
동형일 때이다.
\end{definition}

\noindent
다음 구성은 Sheaves의 정의 \ref{sheaves-definition-restriction} (3)과
평행한 구성이다.

\begin{example}
\label{schemes-example-open-subspace}
$X$를 국소 환 달린 공간이라 하자.
$U \subset X$를 열린 부분집합이라 하자.
$\mathcal{O}_U = \mathcal{O}_X|_U$를
$\mathcal{O}_X$의 $U$로의 제한이라 하자.
$u \in U$에 대해 줄기 $\mathcal{O}_{U, u}$는
줄기 $\mathcal{O}_{X, u}$와 같으므로 국소환이다.
따라서 $(U, \mathcal{O}_U)$는 국소 환 달린 공간이고, 사상
$j : (U, \mathcal{O}_U) \to (X, \mathcal{O}_X)$는 열린 몰입이다.
\end{example}

\begin{definition}
\label{schemes-definition-open-subspace}
$X$를 국소 환 달린 공간이라 하자.
$U \subset X$를 열린 부분집합이라 하자.
위의 예 \ref{schemes-example-open-subspace}에 나오는 국소 환 달린 공간
$(U, \mathcal{O}_U)$를 {\it $X$의 $U$에 딸린 열린 부분공간}이라 한다.
\end{definition}

\begin{lemma}
\label{schemes-lemma-open-immersion}
$f : X \to Y$를 국소 환 달린 공간의 열린 몰입이라 하자.
$j : V = f(X) \to Y$를 $Y$의, $f$의 상에 딸린 열린 부분공간이라 하자.
국소 환 달린 공간의 유일한 동형 $f' : X \cong V$가 존재하여
$f = j \circ f'$을 만족한다.
\end{lemma}

\begin{proof}
$f'$을 $X$와 $V$ 사이에서 $f$가 유도하는 위상동형이라 하자. 그러면
위상공간의 사상으로서 $f = j \circ f'$이다. 층의 동형
$f^\sharp : f^{-1}(\mathcal{O}_Y) \to \mathcal{O}_X$가 있으므로, 환의 동형
$f^\sharp : \Gamma(U, f^{-1}(\mathcal{O}_Y)) \to \Gamma(U, \mathcal{O}_X)$가
각 열린 부분집합 $U \subset X$에 대해 존재한다.
$\mathcal{O}_V = j^{-1}\mathcal{O}_Y$이고 $f^{-1} = f'^{-1} j^{-1}$이므로
(Sheaves의 보조정리 \ref{sheaves-lemma-pullback-composition}),
$f^{-1}\mathcal{O}_Y = f'^{-1}\mathcal{O}_V$이다. 따라서 사상
$\Gamma(U, f'^{-1}(\mathcal{O}_V)) \to \Gamma(U, f^{-1}(\mathcal{O}_Y))$는
모든 열린 $U \subset X$에 대해 동형이다. 이들을 합성하면 각 열린 부분집합
$U \subset X$에 대해 환의 동형
$$
\Gamma(U, f'^{-1}(\mathcal{O}_V)) \to \Gamma(U, \mathcal{O}_X)
$$
을 얻고, 따라서 층의 동형
$f^{-1}(\mathcal{O}_V) \to \mathcal{O}_X$을 얻는다. 다시 말해 동형
$f'^{\sharp} : f'^{-1}(\mathcal{O}_V) \to \mathcal{O}_X$가 있고, 따라서
국소 환 달린 공간의 동형
$(f', f'^{\sharp}) : (X, \mathcal{O}_X) \to (V, \mathcal{O}_V)$가 있다
(보조정리 \ref{schemes-lemma-isomorphism-locally-ringed}을 사용한다).
구성에 의해 국소 환 달린 공간의 사상으로서 $f = j \circ f'$이다.

\medskip\noindent
다른 사상 $f'' : (X, \mathcal{O}_X) \to (V, \mathcal{O}_V)$가
$f = j \circ f''$을 만족한다고 하자.
각 점 $x \in X$에서 $j(f'(x)) = j(f''(x))$이다. 여기서
$f'(x) = f''(x)$인 것은 $j$가 포함사상이기 때문이다. 따라서
$f'$과 $f''$은 위상공간의 사상으로서 같다.
구조층 위에서는 각 열린 부분집합 $U \subset X$에 대해 가환 그림
$$
\xymatrix @R=5em{
\Gamma(U, f^{-1}(\mathcal{O}_Y)) \ar[d]_\cong\ar[r]^\cong &
\Gamma(U, \mathcal{O}_X) \\
\Gamma(U, f'^{-1}(\mathcal{O}_V)) \ar@/^/[ru]^{f'^\sharp}
\ar@/_/[ru]_{f''^\sharp} &
}
$$
을 얻는다. 이로부터 $f'^\sharp$과 $f''^\sharp$이 같은 층의 사상을
정의함을 알 수 있다.
\end{proof}

\noindent
이제부터 열린 부분집합과 그에 딸린 부분공간을 구별하지 않는다.

\begin{lemma}
\label{schemes-lemma-restrict-map-to-opens}
$f : X \to Y$를 국소 환 달린 공간의 사상이라 하자.
$U \subset X$와 $V \subset Y$를 열린 부분집합이라 하자.
$f(U) \subset V$라고 가정하자. 다음 그림이 국소 환 달린 공간들의
가환 사각형이 되도록 하는 국소 환 달린 공간의 유일한 사상
$f|_U : U \to V$가 존재한다.
$$
\xymatrix{
U \ar[d]_{f|_U} \ar[r] & X \ar[d]^f \\
V \ar[r] & Y
}
$$
\end{lemma}

\begin{proof}
생략한다.
\end{proof}

\noindent
이하에서는 위 보조정리로부터 따르는 다음 사실을 따로 언급하지 않고 사용한다.
국소 환 달린 공간의 임의의 사상 $f : Y \to X$와 임의의 열린 부분집합
$U \subset X$가 주어지고 $f(Y) \subset U$를 만족하면,
국소 환 달린 공간의 유일한 사상 $Y \to U$가 존재하여 합성
$Y \to U \to X$가 $f$와 같아진다. 실제로 혼동을 일으키는 일이 드물기 때문에
표기의 남용으로 $f : Y \to U$라고까지 쓸 것이다.


\section{국소 환 달린 공간의 닫힌 몰입}
\label{schemes-section-closed-immersion}

\noindent
Modules의 정의 \ref{modules-definition-closed-immersion}에서 도입한
관례를 따른다.

\begin{definition}
\label{schemes-definition-closed-immersion-locally-ringed-spaces}
$i : Z \to X$를 국소 환 달린 공간의 사상이라 하자.
$i$를 다음 조건들을 만족하는 경우 {
\it 닫힌 몰입}이라 한다.
\begin{enumerate}
\item 사상 $i$는 $Z$를 $X$의 닫힌 부분집합 위로 보내는 위상동형이다.
\item 사상 $\mathcal{O}_X \to i_*\mathcal{O}_Z$는 전사이다.
그 핵을 $\mathcal{I}$로 나타낸다.
\item $\mathcal{O}_X$-가군 $\mathcal{I}$는 단면들로 국소 생성된다.
\end{enumerate}
\end{definition}

\begin{lemma}
\label{schemes-lemma-closed-local-target}
$f : Z \to X$를 국소 환 달린 공간의 사상이라 하자.
$f$가 닫힌 몰입이려면 열린 덮개 $X = \bigcup U_i$가 존재하여
각 사상 $f : f^{-1}U_i \to U_i$가 닫힌 몰입이면 충분하다.
\end{lemma}

\begin{proof}
생략한다.
\end{proof}

\begin{example}
\label{schemes-example-closed-subspace}
$X$를 국소 환 달린 공간이라 하자.
$\mathcal{I} \subset \mathcal{O}_X$를 $\mathcal{O}_X$-가군의 층으로서
단면들로 국소 생성되는 아이디얼층이라 하자. $Z$를 환의 층
$\mathcal{O}_X/\mathcal{I}$의 받침이라 하자. 이는 Modules의 보조정리
\ref{modules-lemma-support-sheaf-rings-closed}에 의해 $X$의 닫힌 부분집합이다.
$i : Z \to X$로 포함사상을 나타내자.
Modules의 보조정리 \ref{modules-lemma-i-star-exact}에 의해
유일한 환의 층 $\mathcal{O}_Z$가 $Z$ 위에 존재하며
$i_*\mathcal{O}_Z = \mathcal{O}_X/\mathcal{I}$를 만족한다. 임의의 $z \in Z$에 대해 줄기
$\mathcal{O}_{Z, z}$는 국소환의 0이 아닌 몫
$\mathcal{O}_{X, i(z)}/\mathcal{I}_{i(z)}$와 같으므로 국소환이다.
따라서 $i : (Z, \mathcal{O}_Z) \to (X, \mathcal{O}_X)$는
국소 환 달린 공간의 닫힌 몰입이다.
\end{example}

\begin{definition}
\label{schemes-definition-closed-subspace}
$X$를 국소 환 달린 공간이라 하자.
$\mathcal{I}$를 $X$ 위에서 단면들로 국소 생성되는 아이디얼층이라 하자.
위의 예 \ref{schemes-example-closed-subspace}에 나오는 국소 환 달린 공간
$(Z, \mathcal{O}_Z)$를 {
\it $X$의 아이디얼층 $\mathcal{I}$에 딸린 닫힌 부분공간}이라 한다.
\end{definition}

\begin{lemma}
\label{schemes-lemma-closed-immersion}
$f : X \to Y$를 국소 환 달린 공간의 닫힌 몰입이라 하자.
$\mathcal{I}$를 사상 $\mathcal{O}_Y \to f_*\mathcal{O}_X$의 핵이라 하자.
$i : Z \to Y$를 $Y$의 아이디얼층 $\mathcal{I}$에 딸린 닫힌 부분공간이라 하자.
국소 환 달린 공간의 유일한 동형 $f' : X \cong Z$가 존재하여
$f = i \circ f'$을 만족한다.
\end{lemma}

\begin{proof}
생략한다.
\end{proof}

\begin{lemma}
\label{schemes-lemma-characterize-closed-subspace}
$X$, $Y$를 국소 환 달린 공간이라 하자. $\mathcal{I} \subset \mathcal{O}_X$를
단면들로 국소 생성되는 아이디얼층이라 하자. $i : Z \to X$를 이에 딸린
닫힌 부분공간이라 하자. 사상 $f : Y \to X$가 $Z$를 통해 인수분해될
필요충분조건은 사상
$f^*\mathcal{I} \to f^*\mathcal{O}_X = \mathcal{O}_Y$가 0인 것이다.
이 조건이 성립하면 사상 $g : Y \to Z$가 유일하게 존재하여
$f = i \circ g$를 만족한다.
\end{lemma}

\begin{proof}
$f$가 $Y \to Z \to X$로 인수분해되면 사상
$f^*\mathcal{I} \to \mathcal{O}_Y$가 0임은 분명하다. 역으로
$f^*\mathcal{I} \to \mathcal{O}_Y$가 0이라고 하자.
임의의 $y \in Y$를 택하고 환 사상
$f^\sharp_y : \mathcal{O}_{X, f(y)} \to \mathcal{O}_{Y, y}$을 생각하자.
가정에 의해 합성
$\mathcal{I}_{f(y)} \to \mathcal{O}_{X, f(y)} \to \mathcal{O}_{Y, y}$가
0이고 $f^\sharp_y(1) = 1$이므로 $1 \not \in \mathcal{I}_{f(y)}$, 즉
$\mathcal{I}_{f(y)} \not = \mathcal{O}_{X, f(y)}$이다. 따라서
$f(Y) \subset Z = \text{Supp}(\mathcal{O}_X/\mathcal{I})$임을 얻는다.
그러므로 $f = i \circ g$이고 $g : Y \to Z$는 연속이다.
사상 $f^\sharp : \mathcal{O}_X \to f_*\mathcal{O}_Y$를 생각하자.
$f^*\mathcal{I} \to \mathcal{O}_Y$가 0이라는 가정은 합성
$\mathcal{I} \to \mathcal{O}_X \to f_*\mathcal{O}_Y$가 0임을 뜻한다.
이는 $f_*$와 $f^*$의 수반성에서 따른다.
다시 말해 환의 층의 사상
$\overline{f^\sharp} : \mathcal{O}_X/\mathcal{I} \to f_*\mathcal{O}_Y$를 얻는다.
$f_*\mathcal{O}_Y = i_*g_*\mathcal{O}_Y$이고
$\mathcal{O}_X/\mathcal{I} = i_*\mathcal{O}_Z$임에 유의하자.
Sheaves의 보조정리 \ref{sheaves-lemma-equivalence-categories-closed-structures}에
의해 환의 층의 유일한 사상
$g^\sharp : \mathcal{O}_Z \to g_*\mathcal{O}_Y$를 얻는데, 이를 $i$로
직상한 것은 $\overline{f^\sharp}$이다. $(g, g^\sharp)$가 국소 환 달린 공간의
사상을 정의한다는 것과 국소 환 달린 공간의 사상으로서
$f = i \circ g$라는 것의 확인은 생략한다. $(g, g^\sharp)$의 유일성은
위에서 이미 지적하였다.
\end{proof}

\begin{lemma}
\label{schemes-lemma-restrict-map-to-closed}
$f : X \to Y$를 국소 환 달린 공간의 사상이라 하자.
$\mathcal{I} \subset \mathcal{O}_Y$를 단면들로 국소 생성되는 아이디얼층이라 하자.
$i : Z \to Y$를 아이디얼층 $\mathcal{I}$에 딸린 닫힌 부분공간이라 하자.
$\mathcal{J}$를 사상
$f^*\mathcal{I} \to f^*\mathcal{O}_Y = \mathcal{O}_X$의 상이라 하자.
그러면 이 아이디얼은 단면들로 국소 생성된다. 또한
$i' : Z' \to X$를 이에 딸린 $X$의 닫힌 부분공간이라 하자.
다음 그림이 국소 환 달린 공간들의 가환 사각형이 되도록 하는
국소 환 달린 공간의 유일한 사상 $f' : Z' \to Z$가 존재한다.
$$
\xymatrix{
Z' \ar[d]_{f'} \ar[r]_{i'} & X \ar[d]^f \\
Z \ar[r]^{i} & Y
}
$$
또한 이 그림은 국소 환 달린 공간의 범주에서 올곱 사각형이다.
\end{lemma}

\begin{proof}
아이디얼 $\mathcal{J}$는 Modules의 보조정리
\ref{modules-lemma-pullback-locally-generated}에 의해 단면들로 국소 생성된다.
보조정리의 나머지는 위의 보조정리
\ref{schemes-lemma-characterize-closed-subspace}에서 주어진, 사상이 닫힌 부분공간을
통해 인수분해된다는 것의 특성화로부터 따른다.
\end{proof}












\section{아핀 스킴}
\label{schemes-section-affine-schemes}

\noindent
환 $R$을 잡자. Algebra의 절 \ref{algebra-section-spectrum-ring}을 참고하여
위상공간 $\Spec(R)$을 생각하자. 이는 $R$에 딸린 공간이다. 이 공간에 환의 층
$\mathcal{O}_{\Spec(R)}$을 부여할 것이며, 그 결과 얻는 쌍
$(\Spec(R), \mathcal{O}_{\Spec(R)})$은 아핀 스킴이 된다.

\medskip\noindent
$\Spec(R)$은 $D(f)$, $f \in R$ 꼴의 열린집합들을 기저로 가지며,
이를 표준 열린집합이라 부른다는 것을 상기하자. Algebra의 정의
\ref{algebra-definition-Zariski-topology}를 보라. 또한 두 표준 열린집합의
교집합도 표준 열린집합이다.
$D(f) \cap D(g) = D(fg)$, $f, g\in R$.

\begin{lemma}
\label{schemes-lemma-standard-open}
환 $R$과 원소 $f \in R$을 잡자.
\begin{enumerate}
\item $g\in R$이고 $D(g) \subset D(f)$이면 다음이 성립한다.
\begin{enumerate}
\item $f$는 $R_g$에서 가역이다.
\item $g^e = af$인 어떤 $e \geq 1$과 $a \in R$가 존재한다.
\item 표준적인 환 사상 $R_f \to R_g$가 존재한다.
\item 표준적인 $R_f$-가군 사상 $M_f \to M_g$가 임의의
$R$-가군 $M$에 대하여 존재한다.
\end{enumerate}
\item $D(f)$의 모든 열린 덮개는
$D(f) = \bigcup_{i = 1}^n D(g_i)$ 꼴의 유한 열린 덮개로 세분된다.
\item $g_1, \ldots, g_n \in R$이면, $D(f) \subset \bigcup D(g_i)$일
필요충분조건은 $g_1, \ldots, g_n$이 $R_f$에서 단위 아이디얼을 생성하는 것이다.
\end{enumerate}
\end{lemma}

\begin{proof}
$D(g) = \Spec(R_g)$임을 상기하자(Algebra의 보조정리
\ref{algebra-lemma-standard-open}을 보라). 따라서 (a)는 $f$의 $R_g$에서의
상이 어떤 소 아이디얼에도 들어 있지 않아 가역이라는 사실에서 따른다.
Algebra의 보조정리 \ref{algebra-lemma-Zariski-topology}를 보라.
$f$의 $R_g$에서의 역원을 $a/g^d$로 쓰자. 이는 $g^d - af$가 $g$의 어떤
거듭제곱에 의해 소멸된다는 뜻이므로 (b)를 얻는다. (c)의 사상
$R_f \to R_g$는 (a)와 국소화의 보편 성질로부터 존재한다. 또는
$b/f^n$을 $a^nb/g^{ne}$로 보내어 정의할 수 있다. 등식
$M_f = M \otimes_R R_f$를 사용하여 가군 위의 사상을 얻을 수 있고,
또는 $M_f \to M_g$를 $x/f^n$을 $a^nx/g^{ne}$로 보내어 정의할 수 있다.

\medskip\noindent
$D(f)$가 준콤팩트임을 상기하자. Algebra의 보조정리
\ref{algebra-lemma-qc-open}을 보라. 따라서 두 번째 명제는 표준
열린집합들이 위상의 기저를 이룬다는 사실에서 바로 따른다.

\medskip\noindent
세 번째 명제는 Algebra의 보조정리
\ref{algebra-lemma-Zariski-topology}에서 바로 따른다.
\end{proof}

\noindent
Sheaves의 절 \ref{sheaves-section-bases}에서 기저 위의 층이라는 개념을
정의했고, 이것이 공간 위의 층이라는 개념과 본질적으로 동치임을 보였다.
Sheaves의 보조정리 \ref{sheaves-lemma-extend-off-basis}와
\ref{sheaves-lemma-extend-off-basis-structures}를 보라. 더 나아가 Sheaves의
보조정리 \ref{sheaves-lemma-cofinal-systems-coverings-standard-case}에서 각
표준 열린집합에 대하여 공종인 열린 덮개들의 계에 대해서만 층 조건을
확인하면 충분함을 보였다. 위의 보조정리에 의해 표준 열린집합들로 이루어진
유한 덮개에 대해서만 확인하면 충분하다.

\begin{definition}
\label{schemes-definition-standard-covering}
환 $R$을 잡자.
\begin{enumerate}
\item $\Spec(R)$의 {\it 표준 열린 덮개}란
$\Spec(R) = \bigcup_{i = 1}^n D(f_i)$ 꼴의 덮개를 말한다. 여기서
$f_1, \ldots, f_n \in R$이다.
\item $D(f) \subset \Spec(R)$이 표준 열린집합이라고 하자. $D(f)$의
{\it 표준 열린 덮개}란 $D(f) = \bigcup_{i = 1}^n D(g_i)$ 꼴의 덮개를
말한다. 여기서 $g_1, \ldots, g_n \in R$이다.
\end{enumerate}
\end{definition}

\noindent
환 $R$과 가군 $M$을 잡고, 후자가 $R$-가군이라고 하자. 표준 열린집합들의 기저 위에 준층
$\widetilde M$을 정의하겠다. $U \subset \Spec(R)$이 표준 열린집합이라고
하자. $f, g \in R$가 $D(f) = D(g)$를 만족한다고 하자. 위의 보조정리
\ref{schemes-lemma-standard-open}에 의해 서로 역인 표준 사상
$M_f \to M_g$와 $M_g \to M_f$가 존재한다. 따라서 임의의 $f$를 골라
$U = D(f)$가 되게 하고 다음과 같이 정의할 수 있다.
$$
\widetilde M(U) = M_f.
$$
$D(g) \subset D(f)$이면 위의 보조정리 \ref{schemes-lemma-standard-open}에 의해
표준 사상
$$
\widetilde M(D(f)) = M_f \longrightarrow M_g = \widetilde M(D(g)).
$$
을 얻음에 유의하자. 이것은 표준 열린집합들의 기저 위에 아벨군의 준층을
분명히 정의한다. $M = R$이면 $\widetilde R$은 표준 열린집합들의 기저
위의 환의 준층이다.

\medskip\noindent
$\widetilde M$의 점 $x \in \Spec(R)$에서의 줄기를 계산하자. $x$가 소
아이디얼 $\mathfrak p \subset R$에 대응한다고 하자. 줄기의 정의에 의해
$$
\widetilde M_x = \colim_{f\in R, f\not\in \mathfrak p} M_f
$$
임을 알 수 있다. 여기서 집합 $\{f \in R, f \not \in \mathfrak p\}$에는
$f \geq f' \Leftrightarrow D(f) \subset D(f')$라는 규칙으로 전순서를 준다.
$f_1, f_2 \in R \setminus \mathfrak p$이면 이 순서에서
$f_1f_2 \geq f_1$이다. 따라서 Algebra의 보조정리
\ref{algebra-lemma-localization-colimit}에 의해
$$
\widetilde M_x = M_{\mathfrak p}.
$$
임을 얻는다.

\medskip\noindent
이제 표준 열린 덮개에 대한 층 조건을 확인하자.
$D(f) = \bigcup_{i = 1}^n D(g_i)$이면 이 덮개에 대한 층 조건은 열
$$
0 \to M_f \to \bigoplus M_{g_i} \to \bigoplus M_{g_ig_j}.
$$
의 완전성과 동치이다. $D(g_i) = D(fg_i)$이므로 이 열을
$$
0 \to M_f \to \bigoplus M_{fg_i} \to \bigoplus M_{fg_ig_j}.
$$
로 다시 쓸 수 있음에 유의하자. 또한 위의 보조정리
\ref{schemes-lemma-standard-open}에 의해 $g_1, \ldots, g_n$은 $R_f$에서 단위
아이디얼을 생성한다. 따라서 가군 $M_f$를 환 $R_f$ 위에서 생각하고 원소
$g_1, \ldots, g_n$에 Algebra의 보조정리
\ref{algebra-lemma-cover-module}을 적용할 수 있다. 이 열이 완전함을 얻는다.
위의 논의에 의해 $\widetilde M$은 표준 열린집합들의 기저 위의 층이다.

\medskip\noindent
따라서 Sheaves의 절 \ref{sheaves-section-bases}의 내용으로부터 유일한 환의 층
$\mathcal{O}_{\Spec(R)}$이 존재하며, 표준 열린집합들 위에서
$\widetilde R$과 일치함을 얻는다. 위에서 계산한 줄기에 의해
이 환의 층의 모든 줄기가 국소환임에 유의하자.

\medskip\noindent
마찬가지로 임의의 $R$-가군 $M$에 대하여 유일한
$\mathcal{O}_{\Spec(R)}$-가군의 층 $\mathcal{F}$가 존재하며, 표준
열린집합들 위에서 $\widetilde M$과 일치한다. Sheaves의 보조정리
\ref{sheaves-lemma-extend-off-basis-module}을 보라.

\begin{definition}
\label{schemes-definition-structure-sheaf}
환 $R$을 잡자.
\begin{enumerate}
\item {\it 구조층 $\mathcal{O}_{\Spec(R)}$}은 환 $R$의 스펙트럼에서
정의되는 유일한 환의 층 $\mathcal{O}_{\Spec(R)}$이며, 표준 열린집합들의
기저 위에서 $\widetilde R$과 일치한다.
\item 국소 환 달린 공간 $(\Spec(R), \mathcal{O}_{\Spec(R)})$을 $R$의
{\it 스펙트럼}이라 부르고 $\Spec(R)$으로 나타낸다.
\item $\mathcal{O}_{\Spec(R)}$-가군의 층 중 $\widetilde M$을
$\Spec(R)$의 모든 열린집합으로 확장한 것을 $\mathcal{O}_{\Spec(R)}$-가군의
층으로서 $M$에 딸린 층이라 한다. 이 층도 $\widetilde M$으로
나타낸다.
\end{enumerate}
\end{definition}

\noindent
지금까지 얻은 결과를 요약하자.

\begin{lemma}
\label{schemes-lemma-spec-sheaves}
환 $R$과 가군 $M$을 잡고, 후자가 $R$-가군이라고 하자. $\widetilde M$을
$\mathcal{O}_{\Spec(R)}$-가군의 층 중 $M$에 딸린 층이라 하자.
\begin{enumerate}
\item $\Gamma(\Spec(R), \mathcal{O}_{\Spec(R)}) = R$이다.
\item $\Gamma(\Spec(R), \widetilde M) = M$이며, 이는 $R$-가군의 등식이다.
\item 모든 $f \in R$에 대하여
$\Gamma(D(f), \mathcal{O}_{\Spec(R)}) = R_f$이다.
\item 모든 $f\in R$에 대하여 $\Gamma(D(f), \widetilde M) = M_f$이며,
이는 $R_f$-가군의 등식이다.
\item $D(g) \subset D(f)$일 때마다 $\mathcal{O}_{\Spec(R)}$과
$\widetilde M$ 위의 제한 사상은 보조정리 \ref{schemes-lemma-standard-open}의 사상
$R_f \to R_g$와 $M_f \to M_g$이다.
\item $\mathfrak p$를 $R$의 소 아이디얼이라 하고, $x \in \Spec(R)$을
그에 대응하는 점이라 하자. $\mathcal{O}_{\Spec(R), x} = R_{\mathfrak p}$이다.
\item $\mathfrak p$를 $R$의 소 아이디얼이라 하고, $x \in \Spec(R)$을
그에 대응하는 점이라 하자. $\widetilde M_x = M_{\mathfrak p}$이며,
이는 $R_{\mathfrak p}$-가군의 등식이다.
\end{enumerate}
더 나아가 이 모든 동일시는 $R$-가군 $M$에 관하여 함자적이다. 특히
함자 $M \mapsto \widetilde M$은 $R$-가군의 범주에서
$\mathcal{O}_{\Spec(R)}$-가군의 범주로 가는 완전함자이다.
\end{lemma}

\begin{proof}
(1)--(7)은 위의 논의에서 분명하다. 함자 $M \mapsto \widetilde M$의
완전성은 함자 $M \mapsto M_{\mathfrak p}$가 완전하다는 사실과 짧은
완전열의 완전성을 줄기에서 확인할 수 있다는 사실에서 따른다.
Modules의 보조정리 \ref{modules-lemma-abelian}을 보라.
\end{proof}

\begin{definition}
\label{schemes-definition-affine-scheme}
{\it 아핀 스킴}은 $\Spec(R)$과 국소 환 달린 공간으로서 동형인
국소 환 달린 공간이다. 여기서 $R$은 어떤 환이다. {\it 아핀 스킴의 사상}은
국소 환 달린 공간의 범주에서의 사상이다.
\end{definition}

\noindent
아핀 스킴은 모든 국소 환 달린 공간 가운데 특별한 역할을 한다.
다음 절은 바로 이를 다룬다.












\section{아핀 스킴의 범주}
\label{schemes-section-category-affine-schemes}

\noindent
$Y$가 아핀 스킴이면 그 점들은 표준적인 $1-1$ 대응에 의해
$\Gamma(Y, \mathcal{O}_Y)$의 소 아이디얼들과 대응함에 유의하자.

\begin{lemma}
\label{schemes-lemma-morphism-into-affine-where-point-goes}
$X$를 국소 환 달린 공간이라 하고 $Y$를 아핀 스킴이라 하자.
$f \in \Mor(X, Y)$를 국소 환 달린 공간의 사상이라 하자. 점
$x \in X$가 주어졌을 때 환 사상들을 생각하자.
$$
\Gamma(Y, \mathcal{O}_Y) \xrightarrow{f^\sharp}
\Gamma(X, \mathcal{O}_X) \to \mathcal{O}_{X, x}
$$
$\mathfrak p \subset \Gamma(Y, \mathcal{O}_Y)$를 $\mathfrak m_x$의
역상이라 하자. $y \in Y$를 이에 대응하는 점이라 하면 $f(x) = y$이다.
\end{lemma}

\begin{proof}
가환 그림
$$
\xymatrix{
\Gamma(X, \mathcal{O}_X) \ar[r] &
\mathcal{O}_{X, x} \\
\Gamma(Y, \mathcal{O}_Y) \ar[r] \ar[u] &
\mathcal{O}_{Y, f(x)} \ar[u]
}
$$
을 생각하자(Sheaves의 정의 \ref{sheaves-definition-f-map} 아래에 나오는
$f$-사상의 논의를 보라). 오른쪽 수직 화살표는 국소 사상이므로
$\mathfrak m_{f(x)}$는 $\mathfrak m_x$의 역상이다. 결론이 따른다.
\end{proof}

\begin{lemma}
\label{schemes-lemma-f-open}
$X$를 국소 환 달린 공간이라 하자. $f \in \Gamma(X, \mathcal{O}_X)$라 하자.
집합
$$
D(f) = \{x \in X \mid \text{image }f \not\in \mathfrak m_x\}
$$
은 열려 있다. 또한 $f|_{D(f)}$는 역원을 가진다.
\end{lemma}

\begin{proof}
이는 Modules의 보조정리 \ref{modules-lemma-s-open}의 특수한 경우이지만,
직접적인 증명도 주겠다. $U \subset X$와 $V \subset X$를 두 열린집합이라
하고, $f|_U$는 역원 $g$를, $f|_V$는 역원 $h$를 가진다고 하자. 그러면
$g|_{U\cap V} = h|_{U\cap V}$임은 분명하다. 따라서 $f$가 임의의
$x \in D(f)$의 어떤 열린 근방에서 가역임을 보이면 충분하다.
$f \not \in \mathfrak m_x$이면 $f \in \mathcal{O}_{X, x}$가 역원
$g \in \mathcal{O}_{X, x}$를 가진다. 이는 어떤 열린 근방
$x \in U \subset X$가 존재하여 $g \in \mathcal{O}_X(U)$이고
$g\cdot f|_U = 1$임을 뜻하므로 결론은 분명하다.
\end{proof}

\begin{lemma}
\label{schemes-lemma-f-open-affine}
위의 보조정리 \ref{schemes-lemma-f-open}에서 $X$가 아핀 스킴이면, 열린집합
$D(f)$는 앞서 정의한 표준 열린집합 $D(f)$와 일치한다(Algebra의 정의
\ref{algebra-definition-spectrum-ring}을 보라).
\end{lemma}

\begin{proof}
생략한다.
\end{proof}

\begin{lemma}
\label{schemes-lemma-morphism-into-affine}
\begin{reference}
이 사실의 참고문헌은 \cite[II, Err 1, Prop. 1.8.1]{EGA}이며,
거기서는 J. Tate에게 귀속되어 있다.
\end{reference}
$X$를 국소 환 달린 공간이라 하고 $Y$를 아핀 스킴이라 하자. 사상
$$
\Mor(X, Y)
\longrightarrow
\Hom(\Gamma(Y, \mathcal{O}_Y), \Gamma(X, \mathcal{O}_X))
$$
은 $f$를 전역 단면 위의 $f^\sharp$로 보내며 전단사이다.
\end{lemma}

\begin{proof}
$Y$가 아핀이므로
$(Y, \mathcal{O}_Y) \cong (\Spec(R), \mathcal{O}_{\Spec(R)})$이고,
여기서 $R$은 어떤 환이다.
증명 중에는 환의 스펙트럼에 관해 아는 사실들에서 바로 따라오는
$Y$와 그 구조층의 성질들을 사용할 것이다. 예를 들어 보조정리
\ref{schemes-lemma-spec-sheaves}를 보라.

\medskip\noindent
위의 보조정리들에 따라 역함수를 구성하자. 환 사상
$\psi_Y : \Gamma(Y, \mathcal{O}_Y) \to \Gamma(X, \mathcal{O}_X)$를 잡자.
먼저 보조정리 \ref{schemes-lemma-morphism-into-affine-where-point-goes}의 규칙으로
이에 대응하는 공간의 사상
$$
\Psi : X \longrightarrow Y
$$
를 정의한다. 다시 말해 $x \in X$가 주어지면 $\Psi(x)$를 다음 소
아이디얼에 대응하는 $Y$의 점으로 정의한다. 그 소 아이디얼은
$\Gamma(Y, \mathcal{O}_Y)$ 안에서 $\mathfrak m_x$의 역상이며, 이 역상은 합성
$
\Gamma(Y, \mathcal{O}_Y) \xrightarrow{\psi_Y}
\Gamma(X, \mathcal{O}_X) \to
\mathcal{O}_{X, x}
$
에 관해 취한다.

\medskip\noindent
사상 $\Psi : X \to Y$가 연속임을 주장한다. 표준 열린집합
$D(g)$, $g \in \Gamma(Y, \mathcal{O}_Y)$들은 $Y$의 위상의 기저이다.
따라서 $\Psi^{-1}(D(g))$가 열려 있음을 보이면 충분하다. $\Psi$의 구성에
의해 역상 $\Psi^{-1}(D(g))$는 정확히 집합
$D(\psi_Y(g)) \subset X$이며, 이는 보조정리 \ref{schemes-lemma-f-open}에 의해
열려 있다. 따라서 $\Psi$는 연속이다.

\medskip\noindent
다음으로 층의 $\Psi$-사상을 $\mathcal{O}_Y$에서 $\mathcal{O}_X$로
가도록 구성하자. Sheaves의 보조정리
\ref{sheaves-lemma-f-map-basis-below-structures}에 의해 제한 사상들과
양립하는 환 사상
$\psi_{D(g)} : \Gamma(D(g), \mathcal{O}_Y) \to
\Gamma(\Psi^{-1}(D(g)), \mathcal{O}_X)$
들을 정의하면 충분하다. 표준 동형
$\Gamma(D(g), \mathcal{O}_Y) = \Gamma(Y, \mathcal{O}_Y)_g$가 존재한다.
이는 $Y$가 아핀 스킴이기 때문이다.
$\psi_Y(g)$는 $D(\psi_Y(g))$ 위에서 가역이므로 표준 사상
$$
\Gamma(Y, \mathcal{O}_Y)_g
\longrightarrow
\Gamma(\Psi^{-1}(D(g)), \mathcal{O}_X)
=
\Gamma(D(\psi_Y(g)), \mathcal{O}_X)
$$
이 존재하며, 국소화의 보편 성질에 의해 사상 $\psi_Y$를 확장한다.
여기서는 이 표준 사상을 택하는 것 외에 선택의 여지가 없음에 유의하자.
이를 표준 동일시
$\Gamma(D(g), \mathcal{O}_Y) = \Gamma(Y, \mathcal{O}_Y)_g$와 결합하여
$\psi_{D(g)}$로 삼는다. 아핀 스킴 위의 제한 사상들도 국소화의 보편
성질로 정의되므로 이는 국소화와 양립한다. 보조정리
\ref{schemes-lemma-spec-sheaves}와 \ref{schemes-lemma-standard-open}을 보라.

\medskip\noindent
이로써 환 달린 공간의 사상
$(\Psi, \psi) : (X, \mathcal{O}_X) \to (Y, \mathcal{O}_Y)$를 정의했고,
전역 단면 위에서 $\psi_Y$를 복원한다. 이것이 국소 환 달린 공간의
사상임을 보이려면 국소환 위에 유도되는 사상
$$
\psi_x : \mathcal{O}_{Y, \Psi(x)} \longrightarrow \mathcal{O}_{X, x}
$$
들이 국소임을 보여야 한다. 이는 보조정리
\ref{schemes-lemma-morphism-into-affine-where-point-goes}의 증명에 나오는 가환
그림과 $\Psi$의 정의에서 바로 따른다.

\medskip\noindent
마지막으로 구성 $(\Psi, \psi) \mapsto \psi_Y$와 구성
$\psi_Y \mapsto (\Psi, \psi)$가 서로 역임을 보여야 한다. 분명히
$\psi_Y \mapsto (\Psi, \psi) \mapsto \psi_Y$이다. 따라서 보일 것은
$\psi_Y$가 주어졌을 때 그것을 낳는 쌍 $(\Psi, \psi)$가 많아야 하나라는
것뿐이다. $\Psi$의 유일성은 보조정리
\ref{schemes-lemma-morphism-into-affine-where-point-goes}에서 보였고, $\Psi$의
유일성이 주어졌을 때 사상 $\psi$의 유일성은 위의 증명 과정에서
지적하였다.
\end{proof}

\begin{lemma}
\label{schemes-lemma-category-affine-schemes}
아핀 스킴의 범주는 환의 범주의 반대범주와 동치이다. 이 동치는 아핀
스킴에 그 구조층의 전역 단면을 대응시키는 함자로 주어진다.
\end{lemma}

\begin{proof}
이는 이제 정의 \ref{schemes-definition-affine-scheme}와 보조정리
\ref{schemes-lemma-morphism-into-affine}에서 분명하다.
\end{proof}

\begin{lemma}
\label{schemes-lemma-standard-open-affine}
$Y$를 아핀 스킴이라 하고 $f \in \Gamma(Y, \mathcal{O}_Y)$라 하자.
열린 부분공간 $D(f)$는 아핀 스킴이다.
\end{lemma}

\begin{proof}
$Y = \Spec(R)$이고 $f \in R$이라고 가정해도 된다. 아핀 스킴의 사상
$\phi : U = \Spec(R_f) \to \Spec(R) = Y$를 생각하자. 이는 환 사상
$R \to R_f$에 의해 유도된다. Algebra의 보조정리
\ref{algebra-lemma-standard-open}에 의해 이것은 $D(f)$ 위로의
위상동형이다. 한편 사상 $\phi^{-1}\mathcal{O}_Y \to \mathcal{O}_U$는
줄기 위에서 동형이므로 동형이다. 따라서 $\phi$가 열린 몰입임을 알 수 있다.
보조정리 \ref{schemes-lemma-open-immersion}에 의해 $D(f)$는 $U$와 동형이다.
\end{proof}

\begin{lemma}
\label{schemes-lemma-fibre-product-affine-schemes}
아핀 스킴의 범주는 유한 곱과 올곱을 가진다. 다시 말해 유한 극한을
가진다. 또한 아핀 스킴의 범주에서의 곱과 올곱은 국소 환 달린 공간의
범주에서의 것과 같다. 국소 환 달린 공간의 범주에서 식으로 쓰면
$$
\Spec(R) \times \Spec(S) =
\Spec(R \otimes_{\mathbf{Z}} S)
$$
이고, 환 사상 $R \to A$, $R \to B$가 주어지면
$$
\Spec(A) \times_{\Spec(R)} \Spec(B)
=
\Spec(A \otimes_R B).
$$
이다.
\end{lemma}

\begin{proof}
이는 보조정리 \ref{schemes-lemma-morphism-into-affine}의 적용일 뿐이다. 먼저 그
보조정리에 의해 아핀 스킴 $\Spec(\mathbf{Z})$은 국소 환 달린 공간의
범주의 끝 대상이다. 따라서 첫 번째 표시식은 두 번째 것에서 따른다.
두 번째 것을 증명하기 위해, 임의의 국소 환 달린 공간 $X$에 대하여
\begin{eqnarray*}
\Mor(X, \Spec(A \otimes_R B))
& = &
\Hom(A \otimes_R B, \mathcal{O}_X(X)) \\
& = &
\Hom(A, \mathcal{O}_X(X))
\times_{\Hom(R, \mathcal{O}_X(X))}
\Hom(B, \mathcal{O}_X(X)) \\
& = &
\Mor(X, \Spec(A))
\times_{\Mor(X, \Spec(R))}
\Mor(X, \Spec(B))
\end{eqnarray*}
임에 유의하자. 이것으로 식이 증명된다. 관련 정의는 Categories의 절
\ref{categories-section-fibre-products}을 보라.
\end{proof}

\begin{lemma}
\label{schemes-lemma-disjoint-union-affines}
$X$를 국소 환 달린 공간이라 하자. $X = U \amalg V$라고 가정하자.
여기서 $U$와 $V$는 열려 있고, $U$, $V$는 아핀 스킴이라고 하자.
그러면 $X$는 아핀 스킴이다.
\end{lemma}

\begin{proof}
$R = \Gamma(X, \mathcal{O}_X)$로 놓자. 층의 성질에 의해
$R = \mathcal{O}_X(U) \times \mathcal{O}_X(V)$임에 유의하자. 보조정리
\ref{schemes-lemma-morphism-into-affine}에 의해 국소 환 달린 공간의 표준 사상
$X \to \Spec(R)$이 존재한다. Algebra의 보조정리
\ref{algebra-lemma-spec-product}에 의해 위상공간으로서
$\Spec(\mathcal{O}_X(U)) \amalg \Spec(\mathcal{O}_X(V)) =
\Spec(R)$
이고, 그 사상들은 환 준동형 $R \to \mathcal{O}_X(U)$와
$R \to \mathcal{O}_X(V)$에서 온다. 물론 이는 $\Spec(R)$이 국소 환 달린
공간의 범주에서도 쌍대곱임을 뜻한다. 가정에 의해 사상
$X \to \Spec(R)$은 $\Spec(\mathcal{O}_X(U))$에서 $U$로 가는 동형을
유도하며 $V$에 대해서도 마찬가지이다. 따라서 $X \to \Spec(R)$은
동형이다.
\end{proof}














\section{아핀 위의 준연접층}
\label{schemes-section-quasi-coherent-affine}

\noindent
Modules의 정의 \ref{modules-definition-quasi-coherent}에서 준연접층이라는
추상적인 개념을 정의했음을 상기하자. 이 절에서는 아핀 스킴 $\Spec(R)$
위의 모든 준연접층이 층 $\widetilde M$에 대응함을 보인다. 이 층은
$R$-가군 $M$에 딸린다.

\begin{lemma}
\label{schemes-lemma-compare-constructions}
$(X, \mathcal{O}_X) = (\Spec(R), \mathcal{O}_{\Spec(R)})$을 아핀 스킴이라
하자. $M$을 $R$-가군이라 하자. 층 $\widetilde M$, 즉 $R$-가군 $M$에
딸린 층(정의 \ref{schemes-definition-structure-sheaf})과 층 $\mathcal{F}_M$, 즉
$R$-가군 $M$에 딸린 층(Modules의 정의
\ref{modules-definition-sheaf-associated}) 사이에는 표준 동형이 존재한다.
이 동형은 $M$에 관하여 함자적이다. 특히 층들 $\widetilde M$은
준연접이다. 더 나아가 이들은 다음 사상 성질로 특징지어진다.
$$
\Hom_{\mathcal{O}_X}(\widetilde M, \mathcal{F})
=
\Hom_R(M, \Gamma(X, \mathcal{F}))
$$
이는 임의의 $\mathcal{O}_X$-가군의 층 $\mathcal{F}$에 대하여 성립한다.
여기서 사상 $\alpha : \widetilde M \to \mathcal{F}$는 전역 단면 위의
작용에 대응한다.
\end{lemma}

\begin{proof}
Modules의 보조정리 \ref{modules-lemma-construct-quasi-coherent-sheaves}에
의해 사상 $\mathcal{F}_M \to \widetilde M$이 존재하며, 이는 사상
$M \to \Gamma(X, \widetilde M) = M$에 대응한다. $x \in X$가 소 아이디얼
$\mathfrak p \subset R$에 대응한다고 하자. 줄기 위에 유도되는 사상들은
$\mathcal{O}_{X, x} \otimes_R M \to M_{\mathfrak p}$이고, 등식
$R_{\mathfrak p} \otimes_R M = M_{\mathfrak p}$에 의해 동형이다. 따라서
사상 $\mathcal{F}_M \to \widetilde M$은 동형이다. 사상 성질은 층
$\mathcal{F}_M$의 사상 성질에서 따른다.
\end{proof}

\begin{lemma}
\label{schemes-lemma-widetilde-constructions}
$(X, \mathcal{O}_X) = (\Spec(R), \mathcal{O}_{\Spec(R)})$을 아핀 스킴이라
하자. 다음 표준 동형들이 존재한다.
\begin{enumerate}
\item
$
\widetilde{M \otimes_R N}
\cong
\widetilde M \otimes_{\mathcal{O}_X} \widetilde N
$,
Modules의 절 \ref{modules-section-tensor-product}을 보라.
\item
$
\widetilde{\text{T}^n(M)}
\cong
\text{T}^n(\widetilde M)
$,
$
\widetilde{\text{Sym}^n(M)}
\cong
\text{Sym}^n(\widetilde M)
$, 그리고
$
\widetilde{\wedge^n(M)}
\cong
\wedge^n(\widetilde M)
$,
Modules의 절 \ref{modules-section-symmetric-exterior}를 보라.
\item $M$이 유한 표시 $R$-가군이면
$
\SheafHom_{\mathcal{O}_X}(\widetilde M, \widetilde N)
\cong
\widetilde{\Hom_R(M,  N)}
$,
Modules의 절 \ref{modules-section-internal-hom}을 보라.
\end{enumerate}
\end{lemma}

\begin{proof}[첫 번째 증명]
보조정리 \ref{schemes-lemma-compare-constructions}와 Modules의 보조정리
\ref{modules-lemma-construct-quasi-coherent-sheaves}에 의해 함자
$M \mapsto \widetilde M$을 $\pi^*$로 볼 수 있으며, 여기서 $\pi$는 환
달린 공간의 사상이다. 가군을 당기는 것은 Modules의 보조정리
\ref{modules-lemma-tensor-product-pullback}과
\ref{modules-lemma-pullback-tensor-algebra}에 의해 텐서 구성과 가환한다.
사상 $\pi : (X, \mathcal{O}_X) \to (\{*\}, R)$는 평탄하다. 예를 들어
$\mathcal{O}_X$의 줄기들이 $R$의 국소화이고(보조정리
\ref{schemes-lemma-spec-sheaves}) 따라서 $R$ 위에서 평탄하기 때문이다. 그러므로
첫 번째 가군이 유한 표시이면 $\pi$에 의한 당김은 Modules의 보조정리
\ref{modules-lemma-pullback-internal-hom}에 의해 내부 Hom과 가환한다.
\end{proof}

\begin{proof}[두 번째 증명]
(1)의 증명. 보조정리 \ref{schemes-lemma-compare-constructions}에 의해 사상
$\widetilde{M \otimes_R N}$에서
$\widetilde M \otimes_{\mathcal{O}_X} \widetilde N$으로 가는 사상을 주려면
전역 단면 위의 사상
$M \otimes_R N \to
\Gamma(X, \widetilde M \otimes_{\mathcal{O}_X} \widetilde N)$을 주어야 한다.
이는 가군의 층의 텐서곱의 정의에 의해 존재한다. 이 사상이 동형임을
보이려면 줄기 위에서 동형임을 확인하면 충분하다. 이는
$\widetilde{M}$의 줄기에 대한 기술(보조정리 \ref{schemes-lemma-spec-sheaves} 또는
Modules의 보조정리 \ref{modules-lemma-construct-quasi-coherent-sheaves}),
텐서곱이 국소화와 가환한다는 사실(Algebra의 보조정리
\ref{algebra-lemma-tensor-product-localization}), 그리고 Modules의 보조정리
\ref{modules-lemma-stalk-tensor-product}에서 따른다.

\medskip\noindent
(2)의 증명. 이는 (1)의 증명과 비슷하며, Algebra의 보조정리
\ref{algebra-lemma-tensor-algebra-localization}와 Modules의 보조정리
\ref{modules-lemma-stalk-tensor-algebra}를 사용한다.

\medskip\noindent
(3)의 증명. 구성 $M \mapsto \widetilde{M}$이 함자적이므로 $R$-선형 사상
$\Hom_R(M, N) \to \Hom_{\mathcal{O}_X}(\widetilde{M}, \widetilde{N})$이
존재한다. 이 사상의 공역은
$\SheafHom_{\mathcal{O}_X}(\widetilde M, \widetilde N)$의 전역 단면이다.
따라서 보조정리 \ref{schemes-lemma-compare-constructions}에 의해
$\mathcal{O}_X$-가군의 사상 $\widetilde{\Hom_R(M,  N)} \to
\SheafHom_{\mathcal{O}_X}(\widetilde M, \widetilde N)$을 얻는다. 줄기들을
비교하여 이것이 동형인지 확인한다. $M$이 $R$-가군으로서 유한 표시이면
$\widetilde M$은 $\mathcal{O}_X$-가군으로서 전역 유한 표시를 가진다.
따라서 Algebra의 보조정리 \ref{algebra-lemma-hom-from-finitely-presented}와
Modules의 보조정리 \ref{modules-lemma-stalk-internal-hom}을 사용하여
결론을 얻는다.
\end{proof}

\begin{proof}[(1)의 세 번째 증명]
임의의 $\mathcal{O}_X$-가군 $\mathcal{F}$에 대하여 $M$, $N$, 그리고
$\mathcal{F}$에 관해 함자적인 다음 동형들이 있다.
\begin{align*}
\Hom_{\mathcal{O}_X}(\widetilde{M} \otimes _{\mathcal{O} _X} \widetilde{N},
\mathcal{F})
& =
\Hom_{\mathcal{O}_X}(\widetilde{M},
\SheafHom_{\mathcal{O} _X} (\widetilde{N}, \mathcal{F})) \\
& =
\Hom_R(M, \Gamma(X,
\SheafHom_{\mathcal{O}_X}(\widetilde{N}, \mathcal{F}))) \\
& =
\Hom_R(M, \Hom_{\mathcal{O}_X}(\widetilde{N}, \mathcal{F})) \\
& =
\Hom_R(M, \Hom_R(N, \Gamma(X,\mathcal{F}))) \\
& =
\Hom_R(M \otimes_R N, \Gamma(X, \mathcal{F})) \\
& =
\Hom_{\mathcal{O}_X}(\widetilde{M \otimes_R N}, \mathcal{F})
\end{align*}
첫 번째 등식은 Modules의 보조정리 \ref{modules-lemma-internal-hom}이다.
두 번째 등식은 보조정리 \ref{schemes-lemma-compare-constructions}에 나오는
$\widetilde{M}$의 보편 성질이다. 세 번째 등식은 $\SheafHom$의 정의에
의해 성립한다. 네 번째 등식은 $\widetilde{N}$의 보편 성질이다. 다섯 번째
등식은 Algebra의 보조정리 \ref{algebra-lemma-hom-from-tensor-product}이다.
마지막 등식은 $\widetilde{M \otimes_R N}$의 보편 성질이다. 요네다
보조정리(Categories의 보조정리 \ref{categories-lemma-yoneda})에 의해
(1)을 얻는다.
\end{proof}

\begin{lemma}
\label{schemes-lemma-widetilde-pullback}
아핀 스킴
$(X, \mathcal{O}_X) = (\Spec(S), \mathcal{O}_{\Spec(S)})$,
$(Y, \mathcal{O}_Y) = (\Spec(R), \mathcal{O}_{\Spec(R)})$
을 잡자. $\psi : (X, \mathcal{O}_X) \to (Y, \mathcal{O}_Y)$를 환 사상
$\psi^\sharp : R \to S$에 대응하는 아핀 스킴의 사상이라 하자(보조정리
\ref{schemes-lemma-category-affine-schemes}를 보라).
\begin{enumerate}
\item $\psi^* \widetilde M = \widetilde{S \otimes_R M}$이며, 이는
$R$-가군 $M$에 관하여 함자적이다.
\item $\psi_* \widetilde N = \widetilde{N_R}$이며, 이는
$S$-가군 $N$에 관하여 함자적이다.
\end{enumerate}
\end{lemma}

\begin{proof}
첫 번째 명제는 보조정리 \ref{schemes-lemma-compare-constructions}의 동일시와
Modules의 보조정리 \ref{modules-lemma-restrict-quasi-coherent}의 결과에서
따른다. 두 번째 명제는 등식
$\psi^{-1}(D(f)) = D(\psi^\sharp(f))$과 그에 따른 등식
$$
\psi_* \widetilde N(D(f)) = \widetilde N(D(\psi^\sharp(f))) =
N_{\psi^\sharp(f)} = (N_R)_f = \widetilde{N_R}(D(f))
$$
에서 원하는 대로 따른다.
\end{proof}














\noindent
위의 보조정리 \ref{schemes-lemma-widetilde-pullback}은 특히 층 $\widetilde M$을
표준 아핀 열린 부분공간 $D(f)$에 제한하면 $\widetilde{M_f}$를 얻는다는
것을 말한다. 이제부터는 이를 더 언급하지 않고 사용할 것이다.

\begin{lemma}
\label{schemes-lemma-quasi-coherent-affine}
$(X, \mathcal{O}_X) = (\Spec(R), \mathcal{O}_{\Spec(R)})$을 아핀 스킴이라
하자. $\mathcal{F}$를 준연접 $\mathcal{O}_X$-가군이라 하자. 그러면
$\mathcal{F}$는 $R$-가군 $\Gamma(X, \mathcal{F})$에 딸린 층과 동형이다.
\end{lemma}

\begin{proof}
$\mathcal{F}$를 준연접 $\mathcal{O}_X$-가군이라 하자. 모든 표준
열린집합 $D(f)$가 준콤팩트이므로 $X$는 국소 준콤팩트하다. 즉, 모든
점은 준콤팩트 근방들의 기본계를 가진다. Topology의 정의
\ref{topology-definition-locally-quasi-compact}를 보라. 따라서 Modules의
보조정리 \ref{modules-lemma-quasi-coherent-module}에 의해 모든 소 아이디얼
$\mathfrak p \subset R$에 대하여, 이에 대응하는 점 $x \in X$와
열린 근방 $x \in U \subset X$가 존재하여 $\mathcal{F}|_U$는 어떤
$\mathcal{O}_X(U)$-가군 $M$에 딸린 준연접층과 동형이다. 다시 말해 이
성질을 가지는 $U$들로 된 열린 덮개를 얻는다. 예를 들어 보조정리
\ref{schemes-lemma-standard-open}에 의해 이 덮개를 표준 열린 덮개로 세분할 수
있다. 따라서 덮개 $\Spec(R) = \bigcup D(f_i)$, $R_{f_i}$-가군 $M_i$,
그리고 동형 $\varphi_i : \mathcal{F}|_{D(f_i)} \to \mathcal{F}_{M_i}$를
얻는다. 여기서 후자는 어떤 $R_{f_i}$-가군 $M_i$에 대한 동형이다.
겹침 위에서는 동형
$$
\xymatrix{
\mathcal{F}_{M_i}|_{D(f_if_j)}
\ar[rr]^{\varphi_i^{-1}|_{D(f_if_j)}}
& &
\mathcal{F}|_{D(f_if_j)}
\ar[rr]^{\varphi_j|_{D(f_if_j)}}
& &
\mathcal{F}_{M_j}|_{D(f_if_j)}.
}
$$
을 얻는다. 이를 $\psi_{ij}$로 나타내자. 삼중 겹침 위에서 코사이클 조건
$$
\psi_{jk}|_{D(f_if_jf_k)}
\circ
\psi_{ij}|_{D(f_if_jf_k)}
=
\psi_{ik}|_{D(f_if_jf_k)}
$$
이 성립함은 분명하다.

\medskip\noindent
열린 부분공간 $D(f_i)$, $D(f_if_j)$, $D(f_if_jf_k)$가 각각 아핀
스킴임을 상기하자. 따라서 보조정리 \ref{schemes-lemma-compare-constructions}에
의해 층 $\mathcal{F}_{M_i}$는 층 $\widetilde M_i$와 동형이다. 특히
$\mathcal{F}_{M_i}(D(f_if_j)) = (M_i)_{f_j}$ 등이다. 또한 위의 보조정리
\ref{schemes-lemma-compare-constructions}에 의해 $\psi_{ij}$는 유일한
$R_{f_if_j}$-가군 동형
$$
\psi_{ij} : (M_i)_{f_j} \longrightarrow (M_j)_{f_i}
$$
에 대응함을 알 수 있다. 즉, 이것은 $\psi_{ij}$의 $D(f_if_j)$ 위 단면에
대한 작용이다. 더 나아가 이들은 코사이클 조건, 즉
$$
\xymatrix{
(M_i)_{f_jf_k}
\ar[rd]_{\psi_{ij}}
\ar[rr]^{\psi_{ik}}
& &
(M_k)_{f_if_j} \\
&
(M_j)_{f_if_k} \ar[ru]_{\psi_{jk}}
}
$$
이 가환한다는 조건을 만족한다(임의의 삼중항 $i, j, k$에 대하여).

\medskip\noindent
이제 Algebra의 보조정리 \ref{algebra-lemma-glue-modules}에 의해
$R$-가군 $M$이 존재하여 $M_i = M_{f_i}$이고 사상들 $\psi_{ij}$와
양립한다. $\mathcal{F}_M = \widetilde M$을 생각하자. 이제
$\widetilde M$이 처음의 준연접층 $\mathcal{F}$와 동형임을 보이는 것은
형식적인 일이다. 실제로 층들 $\mathcal{F}$와 $\widetilde M$은 서로 동형인
$\mathcal{O}_X$-가군 층의 붙이기 자료를 덮개 $X = \bigcup D(f_i)$에
관하여 낳는다. Sheaves의 절 \ref{sheaves-section-glueing-sheaves},
특히 보조정리 \ref{sheaves-lemma-mapping-property-glue}를 보라. 지금
상황에서는 이를 다음과 같이 명시할 수 있다. $R$-가군 사상
$$
M \longrightarrow \Gamma(X, \mathcal{F}).
$$
을 구성하자. $m \in M$이 주어지면
$m_i = m/1 \in M_{f_i} = M_i$를 얻는다. 이는 $M$의 구성에서 따르며,
$M_i$의 구성에 의해 이 원소는 단면
$s_i \in \mathcal{F}(U_i)$에 대응한다(즉, $\varphi^{-1}_i(m_i)$이다).
$s_i|_{D(f_if_j)} = s_j|_{D(f_if_j)}$임을 주장한다. 이는 $M$의 구성에
의해 $\psi_{ij}(m_i) = m_j$이고, $\psi_{ij}$의 구성에 의해서이다. 층
$\mathcal{F}$의 조건에 의해 이 단면들의 모음은 유일한 단면 $s$, 즉
$\mathcal{F}$의 $X$ 위 단면을 낳는다. $m \mapsto s$가 $R$-가군 사상임을 보이는
것은 독자에게 맡긴다. 보조정리 \ref{schemes-lemma-compare-constructions}에 의해
이에 딸린 $\mathcal{O}_X$-가군 사상
$$
\widetilde M \longrightarrow \mathcal{F}.
$$
을 얻는다. 구성에 의해 이 사상은 동형 $\varphi_i^{-1}$로 제한되며,
그 제한은 각 $D(f_i)$ 위에서 이루어지므로 이 사상은 동형이다.
\end{proof}

\begin{lemma}
\label{schemes-lemma-equivalence-quasi-coherent}
$(X, \mathcal{O}_X) = (\Spec(R), \mathcal{O}_{\Spec(R)})$을 아핀 스킴이라
하자. 함자들 $M \mapsto \widetilde M$과
$\mathcal{F} \mapsto \Gamma(X, \mathcal{F})$는 범주 동치
$$
\xymatrix{
\QCoh(\mathcal{O}_X)
\ar@<1ex>[r]
&
\text{Mod}_R
\ar@<1ex>[l]
}
$$
를 정의하며 서로 준역이다. 이 동치는 준연접 $\mathcal{O}_X$-가군의
범주와 $R$-가군의 범주 사이의 동치이다.
\end{lemma}

\begin{proof}
위의 보조정리 \ref{schemes-lemma-compare-constructions}와
\ref{schemes-lemma-quasi-coherent-affine}를 보라.
\end{proof}

\noindent
이제부터 아핀 스킴 위의 준연접층과 $\widetilde M$ 꼴의 층을 구별하지
않을 것이다.

\begin{lemma}
\label{schemes-lemma-kernel-cokernel-quasi-coherent}
$X = \Spec(R)$을 아핀 스킴이라 하자. 준연접 $\mathcal{O}_X$-가군의
사상의 핵과 여핵은 준연접이다.
\end{lemma}

\begin{proof}
이는 함자 $\widetilde{\ }$의 완전성에서 따른다. 실제로 보조정리
\ref{schemes-lemma-compare-constructions}에 의해 모든 사상
$\psi : \widetilde{M} \to \widetilde{N}$는 어떤 $R$-가군 사상
$\varphi : M \to N$에서 온다. 따라서
$\Ker(\psi) = \widetilde{\Ker(\varphi)}$이고
$\Coker(\psi) = \widetilde{\Coker(\varphi)}$이다.
\end{proof}

\begin{lemma}
\label{schemes-lemma-colimit-quasi-coherent}
$X = \Spec(R)$을 아핀 스킴이라 하자. $X$ 위의 준연접층들의 임의의
모음의 직합은 준연접이다. 쌍극한에 대해서도 마찬가지이다.
\end{lemma}

\begin{proof}
$\mathcal{F}_i$, $i \in I$가 $X$ 위의 준연접층들의 모음이라고 하자.
위의 보조정리 \ref{schemes-lemma-equivalence-quasi-coherent}에 의해
$\mathcal{F}_i = \widetilde{M_i}$로 쓸 수 있다. 여기서 이는 어떤
$R$-가군 $M_i$에 대한 것이다.
$M = \bigoplus M_i$로 놓고 층 $\widetilde{M}$을 생각하자. 각 표준
열린집합 $D(f)$에 대하여
$$
\widetilde{M}(D(f)) = M_f =
\left(\bigoplus M_i\right)_f =
\bigoplus M_{i, f}.
$$
이다. 따라서 준연접 $\mathcal{O}_X$-가군 $\widetilde{M}$은 층들
$\mathcal{F}_i$의 직합이다. 일반 쌍극한에도 비슷한 논증이 적용된다.
\end{proof}

\begin{lemma}
\label{schemes-lemma-extension-quasi-coherent}
$(X, \mathcal{O}_X) = (\Spec(R), \mathcal{O}_{\Spec(R)})$을 아핀 스킴이라
하자. $\mathcal{O}_X$-가군 층들의 짧은 완전열
$$
0 \to
\mathcal{F}_1 \to
\mathcal{F}_2 \to
\mathcal{F}_3 \to
0
$$
이 주어졌다고 하자. 셋 중 둘이 준연접이면 나머지 하나도 준연접이다.
\end{lemma}

\begin{proof}
$\mathcal{F}_1$과 $\mathcal{F}_2$가 모두 준연접인 경우에는 함자
$M \mapsto \widetilde M$이 완전하므로 분명하다. 보조정리
\ref{schemes-lemma-spec-sheaves}를 보라. $\mathcal{F}_2$와 $\mathcal{F}_3$가
모두 준연접인 경우도 마찬가지이다. 이제
$\mathcal{F}_1 = \widetilde M_1$과 $\mathcal{F}_3 = \widetilde M_3$가
준연접이라고 하자. $M_2 = \Gamma(X, \mathcal{F}_2)$로 놓자. 열
$$
0 \to M_1 \to M_2 \to M_3 \to 0
$$
이 완전함을 보이면 충분하다고 주장한다. 실제로 그렇다면 보조정리
\ref{schemes-lemma-compare-constructions}의 사상 성질을 사용하여 가환 그림
$$
\xymatrix{
0 \ar[r] &
\widetilde M_1 \ar[r] \ar[d] &
\widetilde M_2 \ar[r] \ar[d] &
\widetilde M_3 \ar[r] \ar[d] &
0 \\
0 \ar[r] &
\mathcal{F}_1 \ar[r] &
\mathcal{F}_2 \ar[r] &
\mathcal{F}_3 \ar[r] &
0
}
$$
을 얻고 뱀 보조정리에 의해 결론을 얻는다.

\medskip\noindent
여기서 ``올바른'' 논증은 먼저 $H^1(X, \mathcal{F}) = 0$임을 임의의
준연접층 $\mathcal{F}$에 대하여 보이는 것이다. 이는 실제로 그리 어렵지
않지만 나중으로 미루는 편이 나을 것이다. 대신 작은 요령을 사용한다.

\medskip\noindent
$m \in M_3 = \Gamma(X, \mathcal{F}_3)$을 택하자. 집합
$$
I = \{ f \in R \mid \text{the element }fm\text{ comes from }M_2\}.
$$
을 생각하자. 이는 분명 아이디얼이다. $1 \in I$임을 보이면 충분하다.
따라서 임의의 소 아이디얼 $\mathfrak p$에 대하여 $f \in I$이고
$f \not\in \mathfrak p$인 원소가 존재함을 보이면 충분하다. $x \in X$를
$\mathfrak p$에 대응하는 점이라 하자. 전사성은 줄기에서 확인할 수
있으므로 열린 근방 $U$가 점 $x$에 대하여 존재하여 $m|_U$는 국소 단면
$s \in \mathcal{F}_2(U)$에서 온다. 실제로 $U = D(f)$가 표준
열린집합이라고 가정할 수 있다. 즉, $f \in R$, $f \not \in \mathfrak p$이다.
어떤 $N \gg 0$에 대하여 $f^N \in I$임을 보일 것이며, 그러면 증명이
끝난다.

\medskip\noindent
임의의 점 $z \in V(f)$를 택하자. 이 점은 소 아이디얼
$\mathfrak q \subset R$에 대응한다고 할 수 있다. 또한 $g \in R$,
$g \not \in \mathfrak q$를 찾아
$m|_{D(g)}$가 어떤 $s' \in \mathcal{F}_2(D(g))$로 올라가게 할 수 있다.
차 $s|_{D(fg)} - s'|_{D(fg)}$를 생각하자. 이는 원소 $m'$, 즉
$\mathcal{F}_1(D(fg)) = (M_1)_{fg}$의 원소이다. 어떤 정수
$n = n(z)$에 대하여 원소 $f^n m'$는 어떤 $m'_1 \in (M_1)_g$에서 온다.
$f^n s$는 단면 $\sigma$, 즉 $\mathcal{F}_2$의 $D(f) \cup D(g)$ 위 단면으로
확장됨을 알 수 있다. 실제로 이는 $f^n s' + m'_1$의 제한과
$D(f) \cap D(g) = D(fg)$ 위에서 일치한다. 더 나아가 $\sigma$는
$f^n m$의 $D(f) \cup D(g)$에 대한 제한으로 보내진다.

\medskip\noindent
$V(f)$가 준콤팩트이므로 원소들의 유한 목록 $g_1, \ldots, g_m \in R$,
즉 $V(f) \subset \bigcup D(g_j)$를 만족하는 목록과, 정수 $n > 0$ 및 단면
$\sigma_j \in \mathcal{F}_2(D(f) \cup D(g_j))$들을 택할 수 있다. 이들은
$\sigma_j|_{D(f)} = f^n s$를 만족하고, $\sigma_j$는
$f^nm|_{D(f) \cup D(g_j)}$로 보내진다. 이는 $\mathcal{F}_3$의 단면이다. 차들
$$
\sigma_j|_{D(f) \cup D(g_jg_k)}
-
\sigma_k|_{D(f) \cup D(g_jg_k)}.
$$
을 생각하자. 이들은 $\mathcal{F}_1$의 $D(f) \cup D(g_jg_k)$ 위
단면들에 대응하며 $D(f)$ 위에서는 0이다. 특히
$\mathcal{F}_1(D(g_jg_k)) = (M_1)_{g_jg_k}$에서의 상들은
$(M_1)_{g_jg_kf}$에서 0이다. 따라서 $f$의 어떤 높은 거듭제곱은 이들을
하나도 빠짐없이 소멸시킨다. 다시 말해 원소들 $f^N \sigma_j$는 어떤
$N \gg 0$에 대하여 층 성질의 붙이기 조건을 만족하고, 원하는 단면
$\sigma $, 즉 $\mathcal{F}_2$의 단면을 낳는다. 그 정의역은
$\bigcup (D(f) \cup D(g_j)) = X$이다.
\end{proof}







\section{아핀 스킴의 닫힌 부분공간}
\label{schemes-section-closed-in-affine}

\begin{example}
\label{schemes-example-closed-immersion-affines}
$R$을 환이라 하고 $I \subset R$을 아이디얼이라 하자. 아핀 스킴의 사상
$i : Z = \Spec(R/I) \to \Spec(R) = X$를 생각하자. Algebra의 보조정리
\ref{algebra-lemma-spec-closed}에 의해 이는 $Z$에서 $X$의 닫힌 부분집합
위로 가는 위상동형이다. 더 나아가 $I \subset \mathfrak p \subset R$이
점 $x = i(z)$, $x \in X$, $z \in Z$에 대응하는 소 아이디얼이면 줄기
위에서 사상
$$
\mathcal{O}_{X, x} = R_{\mathfrak p}
\longrightarrow
R_{\mathfrak p}/IR_{\mathfrak p} = \mathcal{O}_{Z, z}
$$
을 얻는다. 따라서 $i$가 국소 환 달린 공간의 닫힌 몰입임을 알 수 있다.
정의 \ref{schemes-definition-closed-immersion-locally-ringed-spaces}를 보라. 이는
예 \ref{schemes-example-closed-subspace}에서처럼 준연접 아이디얼층
$\widetilde I$에 딸린 닫힌 부분공간과 (동형 아래에서) 같음이 분명하다.
\end{example}

\begin{lemma}
\label{schemes-lemma-closed-immersion-affine-case}
\begin{slogan}
아핀 스킴에서 닫힌 몰입은 아이디얼에 대응한다.
\end{slogan}
$(X, \mathcal{O}_X) = (\Spec(R), \mathcal{O}_{\Spec(R)})$을 아핀 스킴이라
하자. $i : Z \to X$를 국소 환 달린 공간의 임의의 닫힌 몰입이라 하자.
그러면 유일한 아이디얼 $I \subset R$이 존재하여 사상 $i : Z \to X$를
위의 예 \ref{schemes-example-closed-immersion-affines}에서 구성한 닫힌 몰입
$\Spec(R/I) \to \Spec(R)$과 동일시할 수 있다.
\end{lemma}

\begin{proof}
이는 다소 우스운 일이다! 실제로 보조정리 \ref{schemes-lemma-closed-immersion}에
의해 $Z \to X$를 정의 \ref{schemes-definition-closed-subspace}와 예
\ref{schemes-example-closed-subspace}에서처럼 아이디얼층
$\mathcal{I} \subset \mathcal{O}_X$에 딸린 닫힌 부분공간과 동일시할 수
있다. 우리의 관례에 따라 이 아이디얼층은 $\mathcal{O}_X$-가군의
층으로서 단면들로 국소 생성된다. 따라서 몫층
$\mathcal{O}_X / \mathcal{I}$는 $X$ 위에서 국소적으로 사상
$\bigoplus_{j \in J} \mathcal{O}_U \to \mathcal{O}_U$의 여핵이다. 그러므로
정의에 의해 $\mathcal{O}_X / \mathcal{I}$는 준연접이다. 절
\ref{schemes-section-quasi-coherent-affine}의 결과에 의해 이는 $\widetilde S$
꼴이며, 어떤 $R$-가군 $S$에 대한 것이다. 더 나아가
$\mathcal{O}_X = \widetilde R \to \widetilde S$가 전사이므로 보조정리
\ref{schemes-lemma-extension-quasi-coherent}에 의해 $\mathcal{I}$도 준연접이며,
$\mathcal{I} = \widetilde I$라고 쓸 수 있다. 물론 $I \subset R$이고
$S = R/I$이므로 모든 것이 분명하다.
\end{proof}













\section{스킴}
\label{schemes-section-schemes}

\begin{definition}
\label{schemes-definition-scheme}
\begin{history}
\cite{EGA1}에서는 우리가 스킴이라 부르는 것을 ``pre-sch\'ema''라
불렀고, ``sch\'ema''라는 이름은 Stacks project에서 분리 스킴이라 하는
것을 위해 남겨 두었다. 제2판 \cite{EGA1-second}에서는 용어를 오늘날의
표준 용어로 바꾸었다. 그러나 때때로 ``prescheme''이라는 용어를 만날 수
있으며, 예를 들면 \cite{Murre-lectures}에서 그러하다.
\end{history}
{\it 스킴}은 모든 점이 아핀 스킴인 열린 근방을 가지는 국소 환 달린
공간이다. {\it 스킴의 사상}은 국소 환 달린 공간의 사상이다. 스킴의
범주는 $\Sch$로 나타낼 것이다.
\end{definition}

\noindent
$X$를 스킴이라 하자. 다음과 같은 (아주 가벼운) 언어의 남용을 사용할
것이다. $U \subset X$이고 열린 부분공간 $U$가 아핀 스킴이면 이를
{\it 아핀 열린집합} 또는 {\it 열린 아핀}이라 할 것이다. 흔히
$U = \Spec(R)$이라고 써서 $U$가 $\Spec(R)$과 동형임을 나타내고, 더
나아가 $U$와 $\Spec(R)$을 (일시적으로) 동일시할 것이다.

\begin{lemma}
\label{schemes-lemma-open-subspace-scheme}
$X$를 스킴이라 하자. $j : U \to X$를 국소 환 달린 공간의 열린
몰입이라 하자. 그러면 $U$는 스킴이다. 특히 $X$의 모든 열린
부분공간은 스킴이다.
\end{lemma}

\begin{proof}
$U \subset X$라 하고 $u \in U$라 하자. 아핀 열린 근방
$u \in V \subset X$를 택하자. $V$의 표준 열린집합들이 $V$의 위상의
기저를 이루므로 어떤 $f\in \mathcal{O}_V(V)$가 존재하여
$u \in D(f) \subset U$이다. 보조정리 \ref{schemes-lemma-standard-open-affine}에
의해 $D(f)$는 아핀 스킴이다. 이것으로 $U$의 모든 점이 아핀인 열린
근방을 가짐을 증명한다.
\end{proof}

\noindent
분명히 이 보조정리(또는 그 증명)는 모든 스킴 $X$의 위상이 아핀
열린집합들로 이루어진 기저를 가짐을 보인다(Topology의 절
\ref{topology-section-bases}를 보라).

\begin{example}
\label{schemes-example-not-affine}
$k$를 체라 하자. 아핀이 아닌 스킴의 한 예는 열린 부분공간
$U = \Spec(k[x, y]) \setminus \{ (x, y)\}$로 주어진다. 이는 아핀 스킴
$X =\Spec(k[x, y])$의 열린 부분공간이다. 이는 두 아핀
$D(x) = \Spec(k[x, y, 1/x])$와 $D(y) = \Spec(k[x, y, 1/y])$로 덮이며,
그 교집합은 $D(xy) = \Spec(k[x, y, 1/xy])$이다. $\mathcal{O}_U$의 층
성질에 의해 완전열
$$
0 \to
\Gamma(U, \mathcal{O}_U) \to
k[x, y, 1/x] \times k[x, y, 1/y] \to
k[x, y, 1/xy]
$$
이 존재한다. 사상 $k[x, y] \to \Gamma(U, \mathcal{O}_U)$이 동형임을
얻는다. 이 사상은 $U \to X$에서 온다. 따라서
$U$는 아핀일 수 없다. 실제로 아핀이면 보조정리
\ref{schemes-lemma-category-affine-schemes}에 의해 $U \cong X$이어야 한다.
\end{example}


\section{스킴의 몰입}
\label{schemes-section-immersions}

\noindent
보조정리 \ref{schemes-lemma-open-subspace-scheme}에서 스킴의 모든 열린
부분공간이 스킴임을 보았다. 아래에서는 스킴의 닫힌 부분공간에도 같은
사실이 성립함을 증명할 것이다.

\medskip\noindent
$\mathcal{O}_X$-가군의 준연접층이라는 개념은 모든 환 달린 공간 $X$에
대하여, 특히 $X$가 스킴일 때 정의됨에 유의하자. 절
\ref{schemes-section-quasi-coherent-affine}의 작업에 의해 그러한 층은 모든 아핀
열린집합 $U \subset X$ 위에서 $\widetilde M$ 꼴임을 안다. 이는 어떤
$\mathcal{O}_X(U)$-가군 $M$에 대한 것이다.

\begin{lemma}
\label{schemes-lemma-closed-subspace-scheme}
$X$를 스킴이라 하자. $i : Z \to X$를 국소 환 달린 공간의 닫힌
몰입이라 하자.
\begin{enumerate}
\item 국소 환 달린 공간 $Z$는 스킴이다.
\item 핵 $\mathcal{I}$, 즉 사상 $\mathcal{O}_X \to i_*\mathcal{O}_Z$의
핵은 준연접 아이디얼층이다.
\item 임의의 아핀 열린집합 $U = \Spec(R)$에 대하여, 이 열린집합은
$X$에 속하며 사상 $i^{-1}(U) \to U$는 사상
$\Spec(R/I) \to \Spec(R)$과 동일시할 수 있다. 여기서 $I \subset R$은
어떤 아이디얼이다.
\item $\mathcal{I}|_U = \widetilde I$이다.
\end{enumerate}
특히 단면들로 국소 생성되는 모든 아이디얼층은 준연접
아이디얼층이고(그 역도 성립하며), $X$의 모든 닫힌 부분공간은 스킴이다.
\end{lemma}

\begin{proof}
$i : Z \to X$를 닫힌 몰입이라 하자. 점 $z \in Z$를 잡고 임의의 아핀
열린 근방 $i(z) \in U \subset X$를 택하자. $U = \Spec(R)$이라고 하자.
보조정리 \ref{schemes-lemma-closed-immersion-affine-case}에 의해
$i^{-1}(U) \to U$를 아핀 스킴의 사상
$\Spec(R/I) \to \Spec(R)$과 동일시할 수 있다. 먼저 이는
$z \in i^{-1}(U) \subset Z$가 $z$의 아핀 근방임을 뜻한다. 따라서 $Z$는
스킴이다. 둘째, 이는 $\mathcal{I}|_U$가 $\widetilde I$임을 뜻한다. 다시
말해 모든 점 $x \in i(Z)$에 대하여 어떤 열린 근방이 존재하여
$\mathcal{I}$는 그 근방에서 준연접이다. 또한
$\mathcal{I}|_{X \setminus i(Z)}
\cong \mathcal{O}_{X \setminus i(Z)}$임에 유의하자. 따라서 아이디얼층의
제한은 $X \setminus i(Z)$ 위에서도 준연접이다. 그러므로
$\mathcal{I}$는 준연접이다.
\end{proof}

\begin{definition}
\label{schemes-definition-immersion}
$X$를 스킴이라 하자.
\begin{enumerate}
\item 스킴의 사상이 국소 환 달린 공간의 열린 몰입이면 이를 {\it 열린
몰입}이라 한다(정의 \ref{schemes-definition-immersion-locally-ringed-spaces}를 보라).
\item $X$의 {\it 열린 부분스킴}은 정의 \ref{schemes-definition-open-subspace}의
의미에서 $X$의 열린 부분공간이다. 보조정리
\ref{schemes-lemma-open-subspace-scheme}에 의해 $X$의 열린 부분스킴은 스킴이다.
\item 스킴의 사상이 국소 환 달린 공간의 닫힌 몰입이면 이를 {\it 닫힌
몰입}이라 한다(정의
\ref{schemes-definition-closed-immersion-locally-ringed-spaces}를 보라).
\item $X$의 {\it 닫힌 부분스킴}은 정의 \ref{schemes-definition-closed-subspace}의
의미에서 $X$의 닫힌 부분공간이다. 보조정리
\ref{schemes-lemma-closed-subspace-scheme}에 의해 닫힌 부분스킴은 스킴이다.
\item 스킴의 사상 $f : X \to Y$가 $j \circ i$로 인수분해되고, 여기서
$i$가 닫힌 몰입이며 $j$가 열린 몰입이면 이를 {\it 몰입} 또는
{\it 국소 닫힌 몰입}이라 한다.
\end{enumerate}
\end{definition}

\noindent
절 \ref{schemes-section-open-immersion}과 \ref{schemes-section-closed-immersion}의
보조정리들로부터, 스킴의 모든 열린 몰입(각각 닫힌 몰입)은 공역의 열린
부분스킴(각각 닫힌 부분스킴)의 포함사상과 동형임이 따른다.

\medskip\noindent
우리의 닫힌 몰입 정의는 Hartshorne과 EGA의 정의 사이에 있다.
Hartshorne은 스킴의 사상 $f : X \to Y$ 중에서 $X$를 $Y$의 닫힌
부분집합 위로 위상동형으로 보내고
$f^\# : \mathcal{O}_Y \to f_*\mathcal{O}_X$가 전사인 것을 닫힌
몰입으로 정의한다. \cite[Page 85]{H}를 보라. 보조정리
\ref{schemes-lemma-characterize-closed-immersions}에서 이것이 우리의 개념과
동치임을 보일 것이다. \cite{EGA}에서 Grothendieck과 Dieudonn\'e는 먼저
준연접 아이디얼층을 사용하여 예 \ref{schemes-example-closed-subspace}의 구성으로
닫힌 부분스킴을 정의하고, 이어서 스킴의 사상 $f : X \to Y$ 중 닫힌
부분스킴과의 동형을 유도하는 것을 닫힌 몰입으로 정의한다. 보조정리
\ref{schemes-lemma-closed-subspace-scheme}에서 이것이 우리의 개념과 일치함이
따른다.

\medskip\noindent
교육적인 관점에서 위의 정의는 재앙이자 악몽이다. 이 내용을 학생들에게
가르칠 때에는, 스킴의 아핀 사상 $f : X \to Y$ 중
$f^\# : \mathcal{O}_Y \to f_*\mathcal{O}_X$가 전사인 것을 닫힌 몰입으로
정의하는 편이 흔히 편리했다. 실제로 아핀 사상이라는 개념(Morphisms의
절 \ref{morphisms-section-affine})은 매우 자연스럽고 이해하기 쉽다.

\medskip\noindent
닫힌 몰입에 대한 자세한 내용은 Morphisms의 절
\ref{morphisms-section-closed-immersions}와
\ref{morphisms-section-closed-immersions-quasi-coherent}를 권한다.

\medskip\noindent
국소 닫힌 부분스킴과 몰입은 이 절의 끝에서 논의할 것이다.

\begin{remark}
\label{schemes-remark-not-reverse-open-closed}
$f : X \to Y$가 스킴의 몰입이면 일반적으로 $f$를 열린 몰입 뒤에 닫힌
몰입이 이어지는 합성으로 인수분해할 수 없다. Morphisms의 예
\ref{morphisms-example-thibaut}를 보라.
\end{remark}

\begin{lemma}
\label{schemes-lemma-immersion-when-closed}
$f : Y \to X$를 스킴의 몰입이라 하자. $f$가 닫힌 몰입일
필요충분조건은 $f(Y) \subset X$가 닫힌 부분집합인 것이다.
\end{lemma}

\begin{proof}
$f$가 닫힌 몰입이면 정의에 의해 $f(Y)$는 닫혀 있다. 역으로
$f(Y)$가 닫혀 있다고 하자. 정의에 의해 열린 부분스킴 $U \subset X$가
존재하여 $f$는 닫힌 몰입 $i : Y \to U$와 열린 몰입 $j : U \to X$의
합성이다. $\mathcal{I} \subset \mathcal{O}_U$를 닫힌 몰입 $i$에 딸린
준연접 아이디얼층이라 하자. 등식
$\mathcal{I}|_{U \setminus i(Y)}
= \mathcal{O}_{U \setminus i(Y)}
= \mathcal{O}_{X \setminus i(Y)}|_{U \setminus i(Y)}$에 유의하자. 따라서
$\mathcal{I}$와 $\mathcal{O}_{X \setminus i(Y)}$를 붙여(Sheaves의 절
\ref{sheaves-section-glueing-sheaves}를 보라) 아이디얼층
$\mathcal{J} \subset \mathcal{O}_X$를 얻을 수 있다. $X$의 모든 점은
$\mathcal{J}$가 준연접인 근방을 가지므로 $\mathcal{J}$가 준연접임을 알
수 있다(특히 단면들로 국소 생성된다). 구성에 의해
$\mathcal{O}_X/\mathcal{J}$는 $U$ 위에 받침을 가지며, 그곳으로 제한하면
$\mathcal{O}_U/\mathcal{I}$와 같다. 따라서 $\mathcal{I}$와 $\mathcal{J}$에
딸린 닫힌 부분공간들이 표준적으로 동형임을 알 수 있다. 예
\ref{schemes-example-closed-subspace}를 보라. 특히 $U$의, $\mathcal{I}$에 딸린
닫힌 부분공간은 $X$의 닫힌 부분공간과 동형이다. 보조정리
\ref{schemes-lemma-closed-immersion}에 의해 $Y \to U$는 $\mathcal{I}$에
딸린 닫힌 부분공간과 동일시되므로 $Y \to U \to X$가 닫힌 몰입임을 얻는다.
\end{proof}

\noindent
$f : Y \to X$를 몰입이라 하자. $Z = \overline{f(Y)} \setminus f(Y)$를
$X$의 닫힌 부분집합이라 하고 $U = X \setminus Z$라 하자. 보조정리에
의해 $U$는 $X$의 열린 부분공간 중 가장 큰 것으로서, 사상
$f : Y \to X$가 $U$로의 닫힌 몰입을 통해 인수분해된다. $X$의
{\it 국소 닫힌 부분스킴}을 쌍 $(Z, U)$로 정의한다. 이때 $Z$는 열린
부분스킴 $U$의 닫힌 부분스킴이고, 그 열린 부분스킴은 $X$의 것이며, 추가로
$\overline{Z} \cup U = X$이다. 보통은 단지 ``$Z$를 $X$의 국소 닫힌
부분스킴이라 하자''라고 말한다. 왜냐하면 $U$는 사상 $Z \to X$로부터
복원할 수 있기 때문이다.
위의 논의는 모든 몰입 $f : Y \to X$가
$Y \to Z \to X$로 유일하게 인수분해됨을 보인다. 여기서 $Z$는 $X$의
국소 닫힌 부분공간이고 $Y \to Z$는 동형이다.

\medskip\noindent
이것이 유용한 이유는 $X$의 국소 닫힌 부분스킴들의 모음이 집합을
이룬다는 점이다. 명백한 이유로 포함이라 부르는 부분순서를 이 집합에
정의할 수 있다. 명시적으로, $Z \to X$와 $Z' \to X$가 $X$의 두 국소
닫힌 부분스킴이면, {\it $Z$가 $Z'$에 포함된다}고 하는 것은 단순히 사상
$Z \to X$가 $Z'$을 통해 인수분해된다는 뜻이다. 그러한 경우 물론 $Z$는 $Z'$의
유일한 국소 닫힌 부분스킴과 동일시되며, 계속 이와 같이 할 수 있다.

\medskip\noindent
몰입에 대한 자세한 내용은 Morphisms의 절
\ref{morphisms-section-immersions}을 보라.









\section{스킴의 자리스키 위상}
\label{schemes-section-topology}

\noindent
스킴의 자리스키 위상에 맞춘 위상수학의 몇 가지 기본 내용은 Topology의
절 \ref{topology-section-introduction}을 보라.

\begin{lemma}
\label{schemes-lemma-scheme-sober}
$X$를 스킴이라 하자. $X$의 모든 기약 닫힌 부분집합은 유일한 일반점을
가진다. 다시 말해 $X$는 소버 위상공간이다. Topology의 정의
\ref{topology-definition-generic-point}를 보라.
\end{lemma}

\begin{proof}
$Z \subset X$를 기약 닫힌 부분집합이라 하자. 모든 아핀 열린집합
$U \subset X$, $U = \Spec(R)$에 대하여, $Z \cap U = V(I)$가 되는
유일한 근기 아이디얼 $I \subset R$이 존재함을 안다. $Z \cap U$는
공집합이거나 기약이다. 두 번째 경우에는(적어도 하나의 $U$에 대해
일어난다) $I = \mathfrak p$가 소 아이디얼이고, $\xi$는 이 소 아이디얼이
주는 $Z \cap U$의 일반점이다. 따라서 $Z = \overline{\{\xi\}}$이다. 다시 말해
$\xi$는 $Z$의 일반점이다. $\xi'$가 두 번째 일반점이라면
$\xi' \in Z \cap U$이고, 여기서 바로 $\xi' = \xi$가 따른다.
\end{proof}

\begin{lemma}
\label{schemes-lemma-basis-affine-opens}
$X$를 스킴이라 하자. $X$의 아핀 열린집합들의 모음은 $X$의 위상의
기저를 이룬다.
\end{lemma}

\begin{proof}
이는 절 \ref{schemes-section-schemes}에서 열린 부분스킴에 관해 논의한 내용에서
따른다.
\end{proof}

\begin{remark}
\label{schemes-remark-intersection-affine-opens}
일반적으로 $X$ 안의 두 아핀 열린집합의 교집합은 아핀 열린집합이 아니다.
예 \ref{schemes-example-affine-space-zero-doubled}를 보라.
\end{remark}

\begin{lemma}
\label{schemes-lemma-locally-quasi-compact}
모든 스킴의 바탕 위상공간은 국소 준콤팩트이다. Topology의 정의
\ref{topology-definition-locally-quasi-compact}를 보라.
\end{lemma}

\begin{proof}
이는 위의 보조정리 \ref{schemes-lemma-basis-affine-opens}와 환의 스펙트럼이
준콤팩트라는 사실에서 따른다. Algebra의 보조정리
\ref{algebra-lemma-quasi-compact}를 보라.
\end{proof}

\begin{lemma}
\label{schemes-lemma-standard-open-two-affines}
$X$를 스킴이라 하자. $U, V$를 $X$의 아핀 열린집합이라 하고
$x \in U \cap V$라 하자. 아핀 열린 근방 $W$가 존재하여 $x$를 포함하고 $W$는
$U$와 $V$ 양쪽의 표준 열린집합이다.
\end{lemma}

\begin{proof}
$U = \Spec(A)$와 $V = \Spec(B)$로 쓰자. $x$가 소 아이디얼
$\mathfrak p \subset A$와 소 아이디얼 $\mathfrak q \subset B$에
대응한다고 하자. $f \in A$, $f \not \in \mathfrak p$를 택하여
$D(f) \subset U \cap V$가 되게 할 수 있다. $D(f)$의 모든 표준
열린집합은 $\Spec(A) = U$의 표준 열린집합임에 유의하자. 따라서
$U \subset V$라고 가정할 수 있다. 다시 말해 이제 $U$를 $V$의 아핀
열린집합으로 생각할 수 있다. 다음으로 $g \in B$,
$g \not \in \mathfrak q$를 택하여 $D(g) \subset U$가 되게 한다. 이때
$D(g) = D(g_A)$임을 알 수 있다. 여기서 $g_A \in A$는 $g$의 사상 $B \to A$
아래에서의 상을 나타낸다. 이로써 보조정리가 증명된다.
\end{proof}

\begin{lemma}
\label{schemes-lemma-good-subcover}
$X$를 스킴이라 하자. $X = \bigcup_i U_i$를 아핀 열린 덮개라 하고
$V \subset X$를 아핀 열린집합이라 하자. 표준 열린 덮개
$V = \bigcup_{j = 1, \ldots, m} V_j$(정의
\ref{schemes-definition-standard-covering}를 보라)가 존재하여 각 $V_j$는
$U_i$들 중 하나의 표준 열린집합이다.
\end{lemma}

\begin{proof}
$v \in V$를 택하자. 그러면 $v \in U_i$인 어떤 $i$가 존재한다. 위의
보조정리 \ref{schemes-lemma-standard-open-two-affines}에 의해 열린집합
$v \in W_v \subset V \cap U_i$가 존재하여 $W_v$는 $V$와 $U_i$
양쪽의 표준 열린집합이다. $V$가 준콤팩트이므로 보조정리가 따른다.
\end{proof}

\begin{lemma}
\label{schemes-lemma-sheaf-on-affines}
$X$를 스킴이라 하자. $\mathcal{B}$를 $X$의 아핀 열린집합들의 집합이라
하자. $\mathcal{F}$를 $\mathcal{B}$ 위의 집합의 준층이라 하자. Sheaves의
정의 \ref{sheaves-definition-presheaf-basis}를 보라. 다음 조건들은
동치이다.
\begin{enumerate}
\item $\mathcal{F}$는 $X$ 위의 한 층을 $\mathcal{B}$로 제한한 것이다.
\item $\mathcal{F}$는 $\mathcal{B}$ 위의 층이다.
\item $\mathcal{F}(\emptyset)$은 한원소 집합이고, $U = V \cup W$이며
$U, V, W \in \mathcal{B}$이고 $V, W \subset U$인 표준 열린집합들에 대하여
(Algebra의 정의 \ref{algebra-definition-Zariski-topology}) 사상
$$
\mathcal{F}(U) \longrightarrow \mathcal{F}(V) \times \mathcal{F}(W)
$$
은 단사이고 그 상은 쌍 $(s, t)$들 가운데
$s|_{V \cap W} = t|_{V \cap W}$를 만족하는 것들의 집합이다.
\end{enumerate}
\end{lemma}

\begin{proof}
(1)과 (2)의 동치는 Sheaves의 보조정리
\ref{sheaves-lemma-restrict-basis-equivalence}이다. (2)가 (3)을 함의함은
분명하다. 따라서 (3)이 (2)를 함의함을 증명하면 충분하다. Sheaves의
보조정리 \ref{sheaves-lemma-cofinal-systems-coverings-standard-case}와
보조정리 \ref{schemes-lemma-standard-open}에 의해 $\mathcal{B}$의 원소들의
표준 열린 덮개(정의 \ref{schemes-definition-standard-covering})에 대해 층 조건이
성립함을 증명하면 충분하다. $U = U_1 \cup \ldots \cup U_n$을 표준 열린
덮개라 하고 $U \subset X$를 아핀 열린집합이라 하자. 이 덮개의 층
조건을 $n$에 대한 귀납법으로 증명할 것이다. $n = 0$이면 $U$는
공집합이고 가정에서 층 조건을 얻는다. $n = 1$이면 증명할 것이 없다.
$n = 2$이면 이는 가정 (3)이다. $n > 2$이면
$U_i = D(f_i)$로 쓰되 $f_i \in A = \mathcal{O}_X(U)$라 하자. 단면
$s_i \in \mathcal{F}(U_i)$들이
$s_i|_{U_i \cap U_j} = s_j|_{U_i \cap U_j}$를 모든 $1 \leq i < j \leq n$에
대하여 만족한다고 하자. $U = U_1 \cup \ldots \cup U_n$이므로
$1 = \sum_{i = 1, \ldots, n} a_i f_i$가 $A$에서 성립하며, 여기서
$a_i \in A$이다. Algebra의 보조정리
\ref{algebra-lemma-Zariski-topology}를 보라.
$g = \sum_{i = 1, \ldots, n - 1} a_if_i$로 놓자. 그러면
$U = D(g) \cup D(f_n)$이다. 또한
$D(g) = D(gf_1) \cup \ldots \cup D(gf_{n - 1})$은 표준 열린 덮개이다.
귀납법에 의해 유일한 단면 $s' \in \mathcal{F}(D(g))$가 존재하여
$s_i|_{D(gfi)}$와 일치하며, 여기서 $i = 1, \ldots, n - 1$이다. $s'$와 $s_n$은
$D(gf_n)$으로의 제한이 같다고 주장한다. 이는 귀납법과 덮개
$D(gf_n) = D(gf_nf_1) \cup \ldots \cup D(gf_nf_{n - 1})$에 의해
성립한다. 따라서 유일한 단면 $s \in \mathcal{F}(U)$가 존재하여
$D(g)$으로의 제한은 $s'$이고 $D(f_n)$으로의 제한은 $s_n$이다. $s$가
$s_i$로 $D(f_i)$ 위에서 제한된다는, $i = 1, \ldots, n - 1$에 대한
확인은 생략하며, $s$가 유일하다는 확인도 생략한다.
\end{proof}

\begin{lemma}
\label{schemes-lemma-scheme-finite-discrete-affine}
바탕 위상공간이 유한 이산 집합인 스킴 $X$를 잡자. 그러면 $X$는
아핀이다.
\end{lemma}

\begin{proof}
$X = \{x_1, \ldots, x_n\}$이라고 하자. 그러면 $U_i = \{x_i\}$는
$x_i$의 열린 근방이다. 보조정리 \ref{schemes-lemma-basis-affine-opens}에 의해
이는 아핀이다. 따라서 $X$는 아핀 스킴들의 유한 서로소 합이고,
보조정리 \ref{schemes-lemma-disjoint-union-affines}에 의해 아핀이다.
\end{proof}

\begin{example}
\label{schemes-example-scheme-without-closed-points}
닫힌점이 없는 스킴이 존재한다. 실제로 국소 정역 $R$을 잡되 그 스펙트럼이
$(0) = \mathfrak p_0 \subset \mathfrak p_1 \subset \mathfrak p_2
\subset \ldots \subset \mathfrak m$처럼 생겼다고 하자.
그러면 열린 부분스킴 $\Spec(R) \setminus \{\mathfrak m\}$은 닫힌점을
가지지 않는다. 이러한 환 $R$이 존재함을 보기 위해, 임의의 전순서군
$(\Gamma, \geq)$가 주어지면 값매김환 $A$가 존재하고 그 값군이
$(\Gamma, \geq)$이라는 사실을 사용한다. \cite{Krull}을 보라. 표기는 Algebra의 절
\ref{algebra-section-valuation-rings}을 보라. 다음과 같이 잡는다.
$\Gamma = \mathbf{Z}x_1 \oplus \mathbf{Z}x_2 \oplus \mathbf{Z}x_3 \oplus
\ldots$ 그리고 $\sum_i a_i x_i \geq 0$을 다음과 같이 정의한다. 첫 번째
0이 아닌 $a_i$가 $> 0$이거나 모든 $a_i = 0$일 때에만 이 부등식이 성립한다. 그러면
$x_1 \geq x_2 \geq x_3 \geq  \ldots \geq 0$이다. 부분집합
$x_i + \Gamma_{\geq 0}$들은 $(\Gamma, \geq)$의 소 아이디얼이다.
Algebra의 보조정리 \ref{algebra-lemma-ideals-valuation-ring} 앞의 표기를
보라. 이들과 $\emptyset$, $\Gamma_{\geq 0}$이 전부인 소 아이디얼들이다.
따라서 Algebra의 보조정리 \ref{algebra-lemma-ideals-valuation-ring}에 의해
$A$는 스펙트럼이 주어진 구조를 가지는 환의 한 예이다.
\end{example}

\section{축약 스킴}
\label{schemes-section-reduced}

\begin{definition}
\label{schemes-definition-reduced}
$X$를 스킴이라 하자. $X$가 {\it 축약}이라고 함은 모든 국소환
$\mathcal{O}_{X, x}$가 축약환인 경우이다.
\end{definition}

\begin{lemma}
\label{schemes-lemma-reduced}
스킴 $X$가 축약일 필요충분조건은 $\mathcal{O}_X(U)$가 축약환이라는 것이
모든 열린집합 $U \subset X$에 대하여 성립하는 것이다.
\end{lemma}

\begin{proof}
$X$가 축약이라고 가정하자. $f \in \mathcal{O}_X(U)$가 $f^n = 0$을
만족하는 단면이라 하자. 그러면 $f$의 $\mathcal{O}_{U, u}$에서의 상은
모든 $u \in U$에 대하여 0이다. 따라서 $f$는 0이다. Sheaves의 보조정리
\ref{sheaves-lemma-sheaf-subset-stalks}를 보라. 역으로,
$\mathcal{O}_X(U)$가 모든 열린집합 $U$에 대하여 축약이라고 가정하자.
0이 아닌 임의의 원소 $f \in \mathcal{O}_{X, x}$를 택하자. 임의의 대표
$(U, f \in \mathcal{O}(U))$는 $f$의 대표로서 0이 아니므로 멱영이 아니다. 따라서
$f$는 $\mathcal{O}_{X, x}$에서 멱영이 아니다.
\end{proof}

\begin{lemma}
\label{schemes-lemma-affine-reduced}
아핀 스킴 $\Spec(R)$가 축약일 필요충분조건은 $R$이 축약환인 것이다.
\end{lemma}

\begin{proof}
정방향은 위의 보조정리 \ref{schemes-lemma-reduced}에서 바로 따른다. 역방향은
축약환의 모든 국소화가 축약이고, 특히 축약환의 국소환들이 축약이므로
따른다.
\end{proof}

\begin{lemma}
\label{schemes-lemma-reduced-closed-subscheme}
$X$를 스킴이라 하고 $T \subset X$를 닫힌 부분집합이라 하자. 다음 성질을
가지는 유일한 닫힌 부분스킴 $Z \subset X$가 존재한다. (a) $Z$의 바탕
위상공간은 $T$와 같고, (b) $Z$는 축약이다.
\end{lemma}

\begin{proof}
다음 규칙으로 정의되는 부분 준층 $\mathcal{I} \subset \mathcal{O}_X$를
생각하자.
$$
\mathcal{I}(U) = \{f \in \mathcal{O}_X(U) \mid
f(t) = 0\text{ for all }t \in T\cap U\}
$$
여기서 $f(t)$는 $f$의 잉여체 $\kappa(t)$ 안의 상을 나타내며, 이 잉여체는
$X$의 $t$에서의 잉여체이다. 이 조건은 국소적이므로 $\mathcal{I}$가 아이디얼의 층임은
분명하다. 또한 $U = \Spec(R)$를 아핀 열린집합이라 하자.
$T \cap U = V(I)$로 쓸 수 있으며, 여기서 $I \subset R$은 유일한 근기
아이디얼이다. 소 아이디얼 $\mathfrak p \in V(I)$가 $t \in T \cap U$에 대응하고 원소
$f \in R$이 주어지면
$f(t) = 0 \Leftrightarrow f \in \mathfrak p$이다. 따라서 Algebra의 보조정리
\ref{algebra-lemma-Zariski-topology}에 의해
$\mathcal{I}(U) = \bigcap_{\mathfrak p \in V(I)} \mathfrak p
= I$이다. 또한 같은 논리에 의해 임의의 표준 열린집합
$D(g) \subset \Spec(R) = U$에 대하여 $\mathcal{I}(D(g)) = I_g$이다.
따라서 $\widetilde I$와 $\mathcal{I}|_U$는 열린집합의 한 기저 위에서
(아이디얼로서) 일치하므로 같다. 그러므로 $\mathcal{I}$는 아이디얼의
준연접층이다.

\medskip\noindent
이제 닫힌 부분공간 $Z$를 아이디얼의 층 $\mathcal{I}$에 결부된 것으로
정의할 수 있다. 아핀 열린집합 $U = \Spec(R)$를 $X$ 안에서 임의로 잡으면
$Z \cap U = \Spec(R/I)$이다. 여기서 $I$는 근기 아이디얼이므로 $Z$는
축약이다(위의 보조정리 \ref{schemes-lemma-affine-reduced}). 구성에 의해 $Z$의
바탕 닫힌 부분집합은 $T$이다. 따라서 성질 (a)와 (b)를 가지는 닫힌
부분스킴을 찾았다.

\medskip\noindent
$Z' \subset X$를 성질 (a)와 (b)를 가지는 두 번째 닫힌 부분스킴이라 하자.
아핀 열린집합 $U = \Spec(R)$를 $X$ 안에서 임의로 잡으면
$Z' \cap U = \Spec(R/I')$이며, 여기서 $I' \subset R$은 어떤 아이디얼이다. 보조정리
\ref{schemes-lemma-affine-reduced}에 의해 환 $R/I'$은 축약이므로 $I'$은 근기
아이디얼이다. $V(I') = T \cap U = V(I)$이므로 Algebra의 보조정리
\ref{algebra-lemma-Zariski-topology}에 의해 $I = I'$임을 얻는다. 따라서
$Z'$와 $Z$는 같은 아이디얼의 층으로 정의되므로 서로 같다.
\end{proof}

\begin{definition}
\label{schemes-definition-reduced-induced-scheme}
$X$를 스킴이라 하고 $Z \subset X$를 닫힌 부분집합이라 하자. {\it $Z$ 위의
스킴 구조}란 닫힌 부분스킴 $Z'$를 $X$ 안에 주되 그 바탕 집합이 $Z$와
같도록 하는 것이다. 이를 나타내기 위해 흔히 ``$(Z, \mathcal{O}_Z)$를 $Z$ 위의
스킴 구조라 하자''라고 말한다. $Z$ 위의 {\it 유도된 축약 스킴 구조}란
보조정리 \ref{schemes-lemma-reduced-closed-subscheme}에서 구성한 구조이다.
{\it 축약 $X_{red}$}이란 $X$의, $X$ 자체 위의 유도된 축약 스킴 구조이다.
\end{definition}

\noindent
``$Z \subset X$를 $X$의 한 기약 성분이라 하자''라고 말할 때에는 흔히
유도된 축약 스킴 구조를 사용하여 $Z$를 $X$의 축약 닫힌 부분스킴으로
생각한다.

\begin{remark}
\label{schemes-remark-reduced-induced-locally-closed}
$X$를 스킴이라 하고 $T \subset X$를 국소 닫힌 부분집합이라 하자. 이
상황에서도 때때로 ``$T$ 위의 유도된 축약 스킴 구조''라는 표현을 쓴다.
이는 $T$를 열린 부분스킴 $X \setminus \partial T$의 닫힌 부분집합으로
보되 이 열린 부분스킴을 $X$ 안에서 취할 때
정의 \ref{schemes-definition-reduced-induced-scheme}에서 얻는 유도된 축약 스킴
구조를 뜻한다. 여기서 $\partial T = \overline{T} \setminus T$는 위상공간
안에서 $T$의 ``경계''이며, 이 위상공간은 $X$의 바탕 공간이다.
\end{remark}

\begin{lemma}
\label{schemes-lemma-map-into-reduction}
$X$를 스킴이라 하고 $Z \subset X$를 닫힌 부분스킴이라 하며 $Y$를 축약
스킴이라 하자. 사상 $f : Y \to X$가 $Z$를 통해 인수분해될 필요충분조건은
(집합론적으로) $f(Y) \subset Z$인 것이다. 특히 임의의 사상 $Y \to X$는
$Y \to X_{red} \to X$로 인수분해된다.
\end{lemma}

\begin{proof}
(집합론적으로) $f(Y) \subset Z$라고 가정하자. $\mathcal{I} \subset \mathcal{O}_X$를
$Z$의 아이디얼 층이라 하자. 임의의 아핀 열린집합
$U \subset X$, $\Spec(B) = V \subset Y$가 $f(V) \subset U$를 만족하고,
임의의 $g \in \mathcal{I}(U)$가 주어졌다고 하자. 당김
$b = f^\sharp(g) \in \Gamma(V, \mathcal{O}_Y) = B$는 모든 $y \in V$의
잉여체에서 0으로 간다. 다시 말해
$b \in \bigcap_{\mathfrak p \subset B} \mathfrak p$이다. 이는 $b = 0$을
함의한다. 실제로 $B$는 축약이다(보조정리 \ref{schemes-lemma-reduced}와 Algebra의
보조정리 \ref{algebra-lemma-Zariski-topology}). 따라서
보조정리 \ref{schemes-lemma-characterize-closed-subspace}에 의해 $f$는 $Z$를 통해
인수분해된다.
\end{proof}

\section{스킴의 점}
\label{schemes-section-points}

\noindent
스킴 $X$가 주어지면 함자
$$
h_X : \Sch^{opp}
\longrightarrow
\textit{Sets}, \quad
T \longmapsto \Mor(T, X).
$$
를 정의할 수 있다. Categories의 예 \ref{categories-example-hom-functor}를
보라. 이를 {\it $X$의 점 함자}라고 한다. 스킴 이론의 흥미로운 부분 중
하나는 $X$의 내부 기하를 함자 $h_X$로 기술하는 것이다. 이
절에서는 $X$의 점을 기술하는 간단한 방법을 찾는다.

\medskip\noindent
$X$를 스킴이라 하자. $R$을 극대 아이디얼 $\mathfrak m \subset R$을 가지는
국소환이라 하자. $f : \Spec(R) \to X$가 스킴의 사상이라고 가정하자.
$x \in X$를 닫힌점 $\mathfrak m \in \Spec(R)$의 상이라 하자. 그러면
국소환들의 국소 준동형
$$
f^\sharp :
\mathcal{O}_{X, x}
\longrightarrow
\mathcal{O}_{\Spec(R), \mathfrak m} = R.
$$
을 얻는다.

\begin{lemma}
\label{schemes-lemma-morphism-from-spec-local-ring}
$X$를 스킴이라 하고 $R$을 국소환이라 하자. 위의 구성은 사상
$\Spec(R) \to X$들과 쌍 $(x, \varphi)$들 사이의 전단사 대응을 준다.
여기서 쌍은 점 $x \in X$와 국소환들의 국소 준동형
$\varphi : \mathcal{O}_{X, x} \to R$로 이루어진다.
\end{lemma}

\begin{proof}
$A$를 환이라 하자. 임의의 환 준동형 $\psi : A \to R$에 대하여 유일한
소 아이디얼 $\mathfrak p \subset A$와 인수분해
$A \to A_{\mathfrak p} \to R$가 존재하고, 마지막 사상은 국소환들의 국소
준동형이다. 실제로 $\mathfrak p = \psi^{-1}(\mathfrak m)$이다. 보조정리
\ref{schemes-lemma-morphism-into-affine}를 적용하면 $X$가 아핀 스킴일 때 보조정리가
성립함을 알 수 있다.

\medskip\noindent
$X$를 일반적인 스킴이라 하자. 임의의 $x \in X$는 아핀 열린집합
$U \subset X$에 포함된다. 아핀 경우에 의해 모든 쌍 $(x, \varphi)$가
보조정리 앞의 구성 결과로 나타남을 알 수 있다.

\medskip\noindent
증명을 마치려면 임의의 사상 $f : \Spec(R) \to X$의 상이 닫힌점의 상
$x$를 포함하는 모든 아핀 열린집합에 포함됨을 보이면 충분하다. 여기서
닫힌점은 $\Spec(R)$의 닫힌점이다.
실제로 $x \in V \subset X$를 $x$를 포함하는 임의의 열린
근방이라 하자. 그러면 $f^{-1}(V) \subset \Spec(R)$은 유일한 닫힌점을
포함하는 열린집합이므로 $\Spec(R)$과 같다.
\end{proof}

\noindent
위 보조정리의 특수한 경우로, 각 점 $x$에 대하여 스킴 $X$에는 표준 사상
\begin{equation}
\label{schemes-equation-canonical-morphism}
\Spec(\mathcal{O}_{X, x}) \longrightarrow X
\end{equation}
을 얻으며, 이는 $X$의 $x$에서의 국소환의 항등사상에 대응한다. 위의
보조정리는 임의의 사상 $f : \Spec(R) \to X$에 대하여 유일한 점
$x \in X$가 존재하여 $f$가
$\Spec(R) \to \Spec(\mathcal{O}_{X, x}) \to X$로 인수분해되고 첫 번째
사상은 국소 준동형 $\mathcal{O}_{X, x} \to R$에서 온다는 말로 다시
표현할 수 있다.

\medskip\noindent
스킴의 사상 $f : X \to S$와 점 $x$가 있고 그 상이 점 $s \in S$인 경우에는
가환 도식
$$
\xymatrix{
\Spec(\mathcal{O}_{X, x}) \ar[r] \ar[d] & X \ar[d] \\
\Spec(\mathcal{O}_{S, s}) \ar[r] & S
}
$$
을 얻는다. 여기서 왼쪽 세로 사상은 국소환 사상
$f^\sharp_x : \mathcal{O}_{S, s} \to \mathcal{O}_{X, x}$에 대응한다.

\begin{lemma}
\label{schemes-lemma-specialize-points}
$X$를 스킴이라 하자. $x, x' \in X$를 $X$의 점들이라 하자. 그러면
$x' \in X$가 $x$의 일반화일 필요충분조건은 $x'$가 표준 사상
$\Spec(\mathcal{O}_{X, x}) \to X$의 상에 속하는 것이다.
\end{lemma}

\begin{proof}
연속사상은 특수화/일반화 관계를 보존한다.
$\Spec(\mathcal{O}_{X, x})$의 모든 점은 닫힌점의 일반화이므로
$\Spec(\mathcal{O}_{X, x}) \to X$의 상의 모든 점은 $x$의 일반화이다.
역으로 $x'$가 $x$의 일반화라고 가정하자. 아핀 열린 근방
$U = \Spec(R)$를 택하되 $x$의 근방이 되게 하자. 그러면 $x' \in U$이다.
$\mathfrak p \subset R$과 $\mathfrak p' \subset R$을 각각 $x$와 $x'$에
대응하는 소 아이디얼들이라 하자. $x'$가 $x$의 일반화이므로
$\mathfrak p' \subset \mathfrak p$이다.
이는 원하는 대로 $\mathfrak p'$가 사상
$\Spec(\mathcal{O}_{X, x}) = \Spec(R_{\mathfrak p})
\to \Spec(R) = U \subset X$의 상에 속한다는 뜻이다.
\end{proof}

\noindent
이제 체의 스펙트럼에서 오는 사상들을 논하자. $(R, \mathfrak m, \kappa)$를
극대 아이디얼 $\mathfrak m$과 잉여체 $\kappa$를 가지는 국소환이라 하자.
$K$를 체라 하자. 정의에 의해 국소 준동형 $R \to K$는
$R \to \kappa \to K$로 인수분해된다. 즉, 이는 준동형 $\kappa \to K$와
같은 데이터이다. 따라서 사상
$$
\Spec(K) \longrightarrow X
$$
들은 쌍 $(x, \kappa(x) \to K)$들과 대응한다. 체의 스펙트럼에서 $X$로
가는 사상들 위에 다음과 같이 준순서를 정의할 수 있다. 사상
$\Spec(K) \to X$가 $\Spec(L) \to X$를 지배한다고 함은
$\Spec(K) \to X$가 $\Spec(L) \to X$를 통해 인수분해되는 경우이다.
이 정의는 다음 개념을 시사한다. 당분간 두 사상
$p : \Spec(K) \to X$와 $q : \Spec(L) \to X$에 대하여, 세 번째 체
$\Omega$와 가환 도식
$$
\xymatrix{
\Spec(\Omega) \ar[r] \ar[d] &
\Spec(L) \ar[d]^q \\
\Spec(K) \ar[r]^p &
X
}
$$
이 존재하면 이 두 사상이 {\it 동치}라고 하자. 물론 이는 세 스킴
$\Spec(K)$, $\Spec(L)$, $\Spec(\Omega)$ 각각의 유일한 점이 모두 같은
$x \in X$로 간다는 것을 즉시 함의한다. 따라서 (위의 논의에 의해) 이러한
도식은 점 $x \in X$와 체들의 가환 도식
$$
\xymatrix{
\Omega &
L \ar[l] \\
K \ar[u] &
\kappa(x) \ar[l] \ar[u]
}
$$
에 대응한다. 이것은 동치관계를 정의한다. 실제로 체 확대들의 임의의 집합
$K_i/\kappa$가 주어지면 어떤 체 확대 $\Omega/\kappa$가 존재하여 모든 체
확대 $K_i$가 $\Omega$에 포함된다.

\begin{lemma}
\label{schemes-lemma-characterize-points}
$X$를 스킴이라 하자. $X$의 점들은 체의 스펙트럼에서 $X$로 가는 사상들의
동치류들과 전단사로 대응한다. 또한 각 동치류는 (유일한 동형을 제외하고
유일한) 최소 원소 $\Spec(\kappa(x)) \to X$를 포함한다.
\end{lemma}

\begin{proof}
위의 논의에서 따른다.
\end{proof}

\noindent
물론 사상 $\Spec(\kappa(x)) \to X$들은 표준 사상
$\Spec(\mathcal{O}_{X, x}) \to X$를 통해 인수분해된다. 이 맥락에서 보조정리
\ref{schemes-lemma-specialize-points}의 내용은 사상 $\Spec(\kappa(x')) \to X$가
$\Spec(\kappa(x')) \to \Spec(\mathcal{O}_{X, x}) \to X$로 인수분해된다는
것이며, 이는 $x'$가 $x$의 일반화일 때마다 성립한다. 스킴의 사상
$f : X \to S$와 점 $x$가 있고 그 상이 점
$s \in S$인 경우에는 가환 도식
$$
\xymatrix{
\Spec(\kappa(x)) \ar[r] \ar[d] &
\Spec(\mathcal{O}_{X, x}) \ar[r] \ar[d] &
X \ar[d] \\
\Spec(\kappa(s)) \ar[r] &
\Spec(\mathcal{O}_{S, s}) \ar[r] &
S.
}
$$

\section{스킴 붙이기}
\label{schemes-section-glueing-schemes}

\noindent
$I$를 집합이라 하자. 각 $i \in I$에 대하여 $(X_i, \mathcal{O}_i)$를 국소 환
달린 공간이라 하자. (사실 뒤따르는 구성은 환 달린 공간에도 똑같이
적용된다.) 각 쌍 $i, j \in I$에 대하여 $U_{ij} \subset X_i$를 열린
부분공간이라 하자. 각 쌍 $i, j \in I$에 대하여
$$
\varphi_{ij} : U_{ij} \to U_{ji}
$$
를 국소 환 달린 공간의 동형이라 하자. 편의를 위해 $U_{ii} = X_i$라고
가정한다. 각 삼중항 $i, j, k \in I$에 대하여 다음을 가정한다.
\begin{enumerate}
\item
$\varphi_{ij}^{-1}(U_{ji} \cap U_{jk}) =  U_{ij} \cap U_{ik}$이고,
\item 도식
$$
\xymatrix{
U_{ij} \cap U_{ik} \ar[rr]_{\varphi_{ik}} \ar[rd]_{\varphi_{ij}} & &
U_{ki} \cap U_{kj} \\
& U_{ji} \cap U_{jk} \ar[ru]_{\varphi_{jk}}
}
$$
이 가환한다.
\end{enumerate}
(특히 $\varphi_{ii} = \text{id}_{X_i}$가 따른다.) 위 조건들을 만족하는 모임
$(I, (X_i)_{i\in I}, (U_{ij})_{i, j\in I}, (\varphi_{ij})_{i, j\in I})$
을 붙이기 자료라고 부르자.

\begin{lemma}
\label{schemes-lemma-glue}
\begin{slogan}
두 국소 환 달린 공간이 있고 첫 번째 공간의 한 부분공간이 다른 공간의 한
부분공간과 동형이면, 두 공간을 붙여 하나의 큰 국소 환 달린 공간을 만들 수
있다.
\end{slogan}
국소 환 달린 공간의 임의의 붙이기 자료가 주어지면 국소 환 달린 공간 $X$와
열린 부분공간들 $U_i \subset X$, 그리고 국소 환 달린 공간의 동형들
$\varphi_i : X_i \to U_i$가 존재하여 다음을 만족한다.
\begin{enumerate}
\item $X=\bigcup_{i\in I} U_i$,
\item $\varphi_i(U_{ij}) = U_i \cap U_j$이고,
\item $\varphi_{ij} =
\varphi_j^{-1}|_{U_i \cap U_j} \circ \varphi_i|_{U_{ij}}$이다.
\end{enumerate}
국소 환 달린 공간 $X$는 다음 사상 성질들로 특징지어진다. 국소 환 달린
공간 $Y$가 주어지면 다음이 성립한다.
\begin{eqnarray*}
\Mor(X, Y) & = & \{ (f_i)_{i\in I} \mid
f_i : X_i \to Y, \ f_j \circ \varphi_{ij} = f_i|_{U_{ij}}\} \\
f & \mapsto & (f|_{U_i} \circ \varphi_i)_{i \in I} \\
\Mor(Y, X) & = &
\left\{
\begin{matrix}
\text{open covering }Y = \bigcup\nolimits_{i \in I} V_i\text{ and }
(g_i : V_i \to X_i)_{i \in I}
\text{ such that}\\
g_i^{-1}(U_{ij}) = V_i \cap V_j
\text{ and }
g_j|_{V_i \cap V_j} = \varphi_{ij} \circ g_i|_{V_i \cap V_j}
\end{matrix}
\right\} \\
g & \mapsto &
V_i = g^{-1}(U_i), \ g_i = \varphi_i^{-1} \circ g|_{V_i}
\end{eqnarray*}
\end{lemma}

\begin{proof}
$X$를 단계별로 구성한다. 집합으로는
$$
X = (\coprod X_i) / \sim.
$$
를 취한다. 여기서 $x \in X_i$와 $x' \in X_j$가 주어졌을 때,
$x \sim x'$라고 함은 $x \in U_{ij}$, $x' \in U_{ji}$이고
$\varphi_{ij}(x) = x'$인 경우이다. 이는 동치관계이다. 실제로 $x \in X_i$,
$x' \in X_j$, $x'' \in X_k$이고 $x \sim x'$ 및 $x' \sim x''$이면
$x' \in U_{ji} \cap U_{jk}$이다. 따라서 붙이기 자료의 조건 (1)에 의해
$x \in U_{ij} \cap U_{ik}$이고 $x'' \in U_{ki} \cap U_{kj}$이며, 조건 (2)에
의해 $\varphi_{ik}(x) = x''$이다. (반사성과 대칭성은 가정
$U_{ii} = X_i$와 $\varphi_{ii} = \text{id}_{X_i}$에서 따른다.) 자연스러운
사상들을 $\varphi_i : X_i \to X$로 나타내자. 또한
$U_i = \varphi_i(X_i) \subset X$로 나타내자. 사상
$\varphi_i : X_i \to U_i$는 전단사임에 유의하자.

\medskip\noindent
$X$ 위의 위상은 다음 규칙으로 정의한다. $U \subset X$가 열린집합일
필요충분조건은 $\varphi_i^{-1}(U)$가 모든 $i$에 대하여 열린집합인 것이다.
이 규칙이 실제로 위상을 정의한다는 확인은 독자에게 맡긴다. 특히 $U_i$는
열린집합이다. 실제로
$\varphi_j^{-1}(U_i)
= U_{ji}$는 $X_j$에서 모든 $j$에 대하여 열려 있다. 또한 임의의 열린집합
$W \subset X_i$에 대하여 상
$\varphi_i(W) \subset U_i$는 열려 있다. 실제로
$\varphi_j^{-1}(\varphi_i(W)) = \varphi_{ji}^{-1}(W \cap U_{ij})$이다. 따라서
$\varphi_i : X_i \to U_i$는 위상동형이다.

\medskip\noindent
국소 환 달린 공간을 얻으려면 환의 층 $\mathcal{O}_X$를 구성해야 한다.
환의 층들 $\mathcal{O}_{U_i} := \varphi_{i, *} \mathcal{O}_i$를 붙여 이를
구성한다. 즉, 가환 도식
$$
\xymatrix{
U_{ij} \ar[rr]_{\varphi_{ij}} \ar[rd]_{\varphi_i|_{U_{ij}}}
& &
U_{ji} \ar[ld]^{\varphi_j|_{U_{ji}}} \\
& U_i \cap U_j &
}
$$
에서 위쪽 화살표는 환 달린 공간의 동형이므로 환의 층들의 유일한 동형
$$
\mathcal{O}_{U_i}|_{U_i \cap U_j}
\longrightarrow
\mathcal{O}_{U_j}|_{U_i \cap U_j}.
$$
을 얻는다. 이들은 Sheaves의 절 \ref{sheaves-section-glueing-sheaves}에서와
같은 코사이클 조건을 만족한다. 그 절의 결과에 의해 환의 층
$\mathcal{O}_X$를 $X$ 위에 얻으며, $\mathcal{O}_X|_{U_i}$는 위에 표시한 붙이기
사상들과 양립하도록 $\mathcal{O}_{U_i}$와 동형이다. 특히
$(X, \mathcal{O}_X)$는 국소 환 달린 공간이다. 실제로 $\mathcal{O}_X$의
줄기들은 대응하는 점들에서 $\mathcal{O}_i$의 줄기들과 같다.

\medskip\noindent
사상 성질들의 증명은 생략한다.
\end{proof}

\begin{lemma}
\label{schemes-lemma-glue-schemes}
\begin{slogan}
스킴들을 붙여 새로운 스킴을 얻을 수 있다.
\end{slogan}
위의 보조정리 \ref{schemes-lemma-glue}에서 모든 $X_i$가 스킴이라고 가정하자.
그러면 얻어진 국소 환 달린 공간 $X$는 스킴이다.
\end{lemma}

\begin{proof}
각 $U_i$가 스킴이므로 모든 $x \in X$가 아핀 근방을 가진다. 따라서 이는
분명하다.
\end{proof}

\noindent
각 $X_i$를 $X$의 열린 부분공간으로 동형들 $\varphi_i$를 통해 생각하는 것이
관례이다. 다음 두 예에서도 그렇게 하겠다.

\begin{example}[원점이 이중화된 아핀 공간]
\label{schemes-example-affine-space-zero-doubled}
$k$를 체라 하고 $n \geq 1$이라 하자.
$X_1 = \Spec(k[x_1, \ldots, x_n])$,
$X_2 = \Spec(k[y_1, \ldots, y_n])$로 놓자.
$0_1 \in X_1$을 극대 아이디얼
$(x_1, \ldots, x_n) \subset k[x_1, \ldots, x_n]$에 대응하는 점이라 하자.
$0_2 \in X_2$를 극대 아이디얼
$(y_1, \ldots, y_n) \subset k[y_1, \ldots, y_n]$에 대응하는 점이라 하자.
$U_{12} = X_1 \setminus \{0_1\}$로 놓고
$U_{21} = X_2 \setminus \{0_2\}$로 놓자. $\varphi_{12} : U_{12} \to U_{21}$를
$k$-대수의 동형 $k[y_1, \ldots, y_n] \to k[x_1, \ldots, x_n]$에서 오는
동형이라 하자. 이 대수 동형은 $y_i$를 $x_i$로 보내며(따라서 유도된
$X_1 \cong X_2$는 $0_1$을 $0_2$로 보낸다), $X$를 붙이기 자료
$(X_1, X_2, U_{12}, U_{21}, \varphi_{12},
\varphi_{21} = \varphi_{12}^{-1})$에서 얻는 스킴이라 하자. 예 앞에서 도입한
약간의 표기 남용을 사용하여 $X_1, X_2 \subset X$를 열린 부분스킴으로
생각한다. 사상 $f : X \to \Spec(k[t_1, \ldots, t_n])$가 존재하며, 이는
$X_1$(각각 $X_2$) 위에서 $k$-대수 사상
$k[t_1, \ldots, t_n] \to k[x_1, \ldots, x_n]$
(각각 $k[t_1, \ldots, t_n] \to k[y_1, \ldots, y_n]$)에 대응한다. 이들은
$t_i$를 $x_i$(각각 $t_i$를 $y_i$)로 보낸다. 이 사상이
$k[t_1, \ldots, t_n]$을 $\Gamma(X, \mathcal{O}_X)$와 동일시함은 쉽게
알 수 있다. $f(0_1) = f(0_2)$이므로 $X$는 아핀이 아니다.

\medskip\noindent
$X_1$과 $X_2$는 $X$의 아핀 열린집합임에 유의하자. 그러나 $n = 2$이면
$X_1 \cap X_2$는 예 \ref{schemes-example-not-affine}에서 기술한 스킴이므로 아핀이
아니다. 따라서 일반적으로 스킴의 아핀 열린집합 두 개의 교집합은 아핀이
아니다. (이 사실은 더 일반적으로 임의의 $n > 1$에 대해서도 성립한다.)

\medskip\noindent
이 예에는 다음과 같은 또 다른 흥미로운 특징이 있다. $n > 1$이면 기약
닫힌 부분집합 $T \subset X$가 많이 존재한다(예를 들어 $X_1$의 닫히지
않은 임의의 점의 폐포를 취하라). 그러나 $T = \{0_1\}$ 또는
$T = \{0_2\}$가 아니면 $0_1 \in T \Leftrightarrow 0_2 \in T$이다. 증명은
생략한다.
\end{example}

\begin{example}[사영직선]
\label{schemes-example-projective-line}
$k$를 체라 하자. $X_1 = \Spec(k[x])$, $X_2 = \Spec(k[y])$로 놓자.
$0 \in X_1$을 극대 아이디얼 $(x) \subset k[x]$에 대응하는 점이라 하자.
$\infty \in X_2$를 극대 아이디얼 $(y) \subset k[y]$에 대응하는 점이라 하자.
$U_{12} = X_1 \setminus \{0\} = D(x) = \Spec(k[x, 1/x])$로 놓고
$U_{21} = X_2 \setminus \{\infty\} = D(y) = \Spec(k[y, 1/y])$로 놓자.
$\varphi_{12} : U_{12} \to U_{21}$를 $k$-대수의 동형
$k[y, 1/y] \to k[x, 1/x]$에서 오는 동형이라 하자. 이 대수 동형은 $y$를
$1/x$로 보낸다. $\mathbf{P}^1_k$를 붙이기 자료
$(X_1, X_2, U_{12}, U_{21}, \varphi_{12},
\varphi_{21} = \varphi_{12}^{-1})$에서 얻는 스킴이라 하자. 예 앞에서 도입한
약간의 표기 남용을 사용하여 $X_i \subset \mathbf{P}^1_k$를 열린
부분스킴으로 생각한다. 이 경우
$\Gamma(\mathbf{P}^1_k, \mathcal{O}) = k$임을 알 수 있다. 실제로 다항식
$g(x)$ 가운데 변수 $x$에 대한 다항식이면서 $g(1/y)$가 $y$에 대한
다항식이기도 한 것은 상수
다항식뿐이다. $\mathbf{P}^1_k$가 무한이므로 $\mathbf{P}^1_k$는 아핀이 아니다.

\medskip\noindent
아핀 열린집합 $U \subset \mathbf{P}^1_k$가 존재하여 $0$과 $\infty$를 모두
포함한다고 주장한다. 실제로
$U = \mathbf{P}^1_k \setminus \{1\}$로 놓자. 여기서 $1$은 $X_1$에서
극대 아이디얼 $(x - 1)$에 대응하는 점이고, $X_2$에서 극대 아이디얼
$(y - 1)$에 대응하는 점이기도 하다. 그러면
$s = 1/(x - 1) = y/(1 - y) \in \Gamma(U, \mathcal{O}_U)$임을 쉽게 알 수 있다.
실제로 $\Gamma(U, \mathcal{O}_U)$가 다항식환 $k[s]$와 같고 대응하는 사상
$U \to \Spec(k[s])$가 스킴의 동형임을 보일 수 있다. 세부사항은 생략한다.
\end{example}

\section{표현가능성 판정법}
\label{schemes-section-representable}

\noindent
이 절에서는 절 \ref{schemes-section-glueing-schemes}의 붙이기 보조정리를 함자의
언어로 다시 표현한다. Categories의 절 \ref{categories-section-opposite}에
있는 내용 일부를 상기하자. 스킴 $X$가 주어지면 함자
$$
h_X : \Sch^{opp}
\longrightarrow
\textit{Sets}, \quad
T \longmapsto \Mor(T, X).
$$
를 정의할 수 있음을 상기하자. 이를 {\it $X$의 점 함자}라고 한다.

\medskip\noindent
$F$를 스킴의 범주에서 집합의 범주로 가는 반변함자라 하자. 식으로 쓰면
$$
F : \Sch^{opp}
\longrightarrow
\textit{Sets}.
$$
이다. Sites의 절 \ref{sites-section-presheaves}와 같은 용어를 사용하겠다.
즉, 스킴 $T$, 원소 $\xi \in F(T)$, 사상 $f : T' \to T$가 주어지면
$f^*\xi$로 원소 $F(f)(\xi)$를 나타내고, 때로는 표기 $\xi|_{T'}$도 사용한다.

\begin{definition}
\label{schemes-definition-representable-functor}
(Categories의 정의 \ref{categories-definition-representable-functor}를 보라.)
$F$를 위와 같이 스킴의 범주에서 집합의 범주로 가는 반변함자라 하자.
$F$가 {\it 스킴으로 표현가능하다} 또는 간단히 {\it 표현가능하다}고 함은
스킴 $X$가 존재하여 $h_X \cong F$인 경우이다.
\end{definition}

\noindent
$F$가 스킴 $X$로 표현되고 $s : h_X \to F$가 동형이라고 가정하자.
Categories의 요네다 보조정리 \ref{categories-lemma-yoneda}에 의해 쌍
$(X, s : h_X \to F)$는 존재한다면 유일한 동형을 제외하고 유일하다. 또한
요네다 보조정리에 따르면 위와 같은 임의의 반변함자 $F$와 임의의 스킴
$Y$에 대하여 전단사
$$
\Mor_{\text{Fun}(\Sch^{opp}, \textit{Sets})} (h_Y, F)
\longrightarrow
F(Y), \quad
s \longmapsto s(\text{id}_Y).
$$
가 존재한다. 역구성은 다음과 같다. 임의의 $\xi \in F(Y)$가 주어지면 함자
사이의 변환 $s_\xi : h_Y \to F$는 임의의 사상 $f : T \to Y$에 원소
$f^*\xi \in F(T)$를 대응시킨다.

\medskip\noindent
특히 $F$가 표현가능한 경우에는 스킴 $X$와 원소 $\xi \in F(X)$가 존재하여
대응하는 사상 $h_X \to F$가 동형이다. 이 경우에는 {\it 쌍 $(X, \xi)$가
$F$를 표현한다}고도 한다. 원소 $\xi \in F(X)$는, 대수적 스택을 논할 때
더 분명해질 이유로 흔히 {\it ``보편족''}이라고 부른다(여기에 향후 참조를
삽입). 지금은 쌍 $(X, \xi)$가 $F$를 표현하면 모든 원소
$\xi' \in F(T)$가 임의의 $T$에 대하여 $\xi' = f^*\xi$의 꼴이고, 이때 사상
$f : T \to X$가 유일함을 관찰하는 것으로 충분하다.

\begin{example}
\label{schemes-example-global-sections}
모든 스킴 $T$에 집합 $F(T) = \Gamma(T, \mathcal{O}_T)$를 대응시키는 규칙을
생각하자. 사상 $f : T' \to T$에 당김 사상
$f^\sharp : \Gamma(T, \mathcal{O}_T) \to \Gamma(T', \mathcal{O}_{T'})$을
사용하면 이를 반변함자로 만들 수 있다. 환 $R$과 원소 $t \in R$가 주어지면
유일한 환 준동형 $\mathbf{Z}[x] \to R$가 존재하며, 이는 $x$를 $t$로 보낸다.
따라서 보조정리 \ref{schemes-lemma-morphism-into-affine}를 사용하면
$$
\Mor(T, \Spec(\mathbf{Z}[x])) =
\Hom(\mathbf{Z}[x], \Gamma(T, \mathcal{O}_T)) =
\Gamma(T, \mathcal{O}_T).
$$
이며, 이는 실제로 동형 $h_{\Spec(\mathbf{Z}[x])} \to F$를 준다.
``보편족'' $\xi$는 무엇인가? 이를 얻으려면 위의 동일시들을
$\text{id}_{\Spec(\mathbf{Z}[x])}$에 적용해야 한다. 위의 동일시 아래에서
이는 분명히 예상한 대로
$\xi = x \in \Gamma(\Spec(\mathbf{Z}[x]),
\mathcal{O}_{\Spec(\mathbf{Z}[x])}) = \mathbf{Z}[x]$
를 준다.
\end{example}

\begin{definition}
\label{schemes-definition-representable-by-open-immersions}
$F$를 스킴의 범주 위에서 집합에 값을 가지는 반변함자라 하자.
\begin{enumerate}
\item $F$가 {\it 자리스키 위상에 대한 층 성질을 만족한다}고 함은, 모든
스킴 $T$, 모든 열린 덮개 $T = \bigcup_{i \in I} U_i$, 그리고 임의의 원소 모임
$\xi_i \in F(U_i)$에 대하여
$\xi_i|_{U_i \cap U_j} =
\xi_j|_{U_i \cap U_j}$이면 유일한 원소 $\xi \in F(T)$가 존재하여
$\xi_i = \xi|_{U_i}$가 $F(U_i)$에서 성립하는 경우이다.
\item {\it 부분함자 $H \subset F$}란 각 스킴 $T$에 부분집합
$H(T) \subset F(T)$를 대응시키고, 사상 $F(f) : F(T) \to F(T')$가 $H(T)$를
$H(T')$ 안으로 보내도록 하는 규칙이다. 여기서 이 조건은 스킴의 모든 사상
$f : T' \to T$에 대하여 성립한다.
\item $H \subset F$를 부분함자라 하자. $H \subset F$가 {\it 열린 몰입으로
표현가능하다}고 함은 모든 쌍 $(T, \xi)$, 즉 스킴 $T$와 원소
$\xi \in F(T)$의 쌍에 대하여 다음 성질을 가지는 열린 부분스킴
$U_\xi \subset T$가 존재하는 경우이다.
\begin{itemize}
\item[(*)] 사상 $f : T' \to T$가 $U_\xi$를 통해 인수분해될 필요충분조건은
$f^*\xi \in H(T')$인 것이다.
\end{itemize}
\item $I$를 집합이라 하자. 각 $i \in I$에 대하여 $H_i \subset F$를
부분함자라 하자. 모임 $(H_i)_{i \in I}$가 {\it $F$를 덮는다}고 함은 모든
$\xi \in F(T)$에 대하여 열린 덮개 $T = \bigcup U_i$가 존재하여
$\xi|_{U_i} \in H_i(U_i)$인 경우이다.
\end{enumerate}
\end{definition}

\noindent
조건 (4)에서 $H_i \subset F$가 모든 $i$에 대하여 열린 몰입으로
표현가능하다면, $(H_i)_{i \in I}$가 $F$를 덮는지 확인하기 위해서는
$F(T) = \bigcup H_i(T)$인지, $T$가 체의 스펙트럼일 때마다 확인하면 충분하다.

\begin{lemma}
\label{schemes-lemma-glue-functors}
$F$를 스킴의 범주 위에서 집합의 범주에 값을 가지는 반변함자라 하자. 다음을
가정하자.
\begin{enumerate}
\item $F$는 자리스키 위상에 대한 층 성질을 만족한다.
\item 집합 $I$와 부분함자들의 모임 $F_i \subset F$가 존재하여 다음을
만족한다.
\begin{enumerate}
\item 각 $F_i$는 표현가능하다.
\item 각 $F_i \subset F$는 열린 몰입으로 표현가능하다.
\item 모임 $(F_i)_{i \in I}$는 $F$를 덮는다.
\end{enumerate}
\end{enumerate}
그러면 $F$는 표현가능하다.
\end{lemma}

\begin{proof}
$X_i$를 $F_i$를 표현하는 스킴이라 하고, $\xi_i \in F_i(X_i) \subset F(X_i)$를
``보편족''이라 하자. $F_j \subset F$가 열린 몰입으로 표현가능하므로 열린집합
$U_{ij} \subset X_i$가 존재하여 $T \to X_i$가 $U_{ij}$를 통해 인수분해될
필요충분조건은 $\xi_i|_T \in F_j(T)$인 것이다. 특히
$\xi_i|_{U_{ij}} \in F_j(U_{ij})$이므로 표준 사상
$\varphi_{ij} : U_{ij} \to X_j$를 얻으며
$\varphi_{ij}^*\xi_j = \xi_i|_{U_{ij}}$이다. $U_{ji}$의 정의에 의해 이는
$\varphi_{ij}$가 $U_{ji}$를 통해 인수분해됨을 함의한다.
또한
$(\varphi_{ij} \circ \varphi_{ji})^*\xi_j
=\varphi_{ji}^*(\varphi_{ij}^*\xi_j) =
\varphi_{ji}^*\xi_i = \xi_j$이므로
$\varphi_{ij} \circ \varphi_{ji} = \text{id}_{U_{ji}}$이다. 실제로 쌍
$(X_j, \xi_j)$는 $F_j$를 표현한다. 특히 사상들
$\varphi_{ij} : U_{ij} \to U_{ji}$는 스킴의 동형이다. 다음으로
$\varphi_{ij}^{-1}(U_{ji} \cap U_{jk}) = U_{ij} \cap U_{ik}$임을 보여야 한다.
이는 다음 이유로 참이다. (a) $U_{ji} \cap U_{jk}$는 $U_{ji}$의 가장 큰
열린집합이고 그 위에서 $\xi_j$의 제한이 $F_k$의 원소가 되며, (b)
$U_{ij} \cap U_{ik}$는 $U_{ij}$의 가장 큰 열린집합이고 그 위에서
$\xi_i$의 제한이 $F_k$의 원소가 되며,
가장 큰 열린집합이며, (c) $\varphi_{ij}^*\xi_j = \xi_i$이다. 또한 절
\ref{schemes-section-glueing-schemes}의 코사이클 조건은 두 사상
$\varphi_{jk}|_{U_{ji} \cap U_{jk}} \circ
\varphi_{ij}|_{U_{ij} \cap U_{ik}}$와
$\varphi_{ik}|_{U_{ij} \cap U_{ik}}$가 모두 $\xi_k$를 원소 $\xi_i$로
당기므로 성립한다. 따라서 보조정리 \ref{schemes-lemma-glue-schemes}를 적용하여
스킴 $X$를 얻는데, 이는 열린 덮개 $X = \bigcup U_i$와 동형들
$\varphi_i : X_i \to U_i$를 가지며 보조정리 \ref{schemes-lemma-glue}의 성질들을
만족한다.
$\xi_i' = (\varphi_i^{-1})^* \xi_i$로 놓자. 보조정리 \ref{schemes-lemma-glue}의
조건들에 의해
$\xi_i'|_{U_i \cap U_j} = \xi_j'|_{U_i \cap U_j}$이다. $F$가 자리스키
위상의 층 조건을 만족하므로 원소 $\xi' \in F(X)$가 존재하여
$\xi_i = \varphi_i^*\xi'|_{U_i}$가 모든 $i$에 대하여 성립한다.
$\varphi_i$가 동형이므로
$(U_i, \xi'|_{U_i})$가 함자 $F_i$를 표현한다는 것도 얻는다.

\medskip\noindent
쌍 $(X, \xi')$가 함자 $F$를 표현한다고 주장한다. 이를 보이기 위해
$T$를 스킴이라 하고 $\xi \in F(T)$라 하자. 유일한 사상 $g : T \to X$를
구성하여 $g^*\xi' = \xi$가 되게 하겠다. 실제로 부분함자들 $F_i$가 $F$를
덮는다는 조건에 의해 열린 덮개 $T = \bigcup V_i$가 존재하여 각 $i$에
대해 제한 $\xi|_{V_i} \in F_i(V_i)$이다. 또한 각 포함 $F_i \subset F$가
열린 몰입으로 표현가능하므로 각 $V_i \subset T$가 이 성질을 가지는 극대
열린집합이라고 가정할 수 있다. $(U_i, \xi'|_{U_i})$가 함자 $F_i$를
표현하므로 유일한 사상 $g_i : V_i \to U_i$를 얻어
$g_i^*\xi'|_{U_i} = \xi|_{V_i}$이다. 겹침 $V_i \cap V_j$ 위에서 사상들
$g_i$와 $g_j$는 일치한다. 예를 들어 둘 다
$\xi'|_{U_i \cap U_j} \in F_i(U_i \cap U_j)$를 같은 원소로 당기기 때문이다.
따라서 사상들 $g_i$는 원하는 유일한 사상 $T \to X$로 붙는다.
\end{proof}

\begin{remark}
\label{schemes-remark-representable-locally-ringed}
함자 $F$가 모든 국소 환 달린 공간 위에 정의되어 있고, 보조정리
\ref{schemes-lemma-glue-functors}의 조건들을 다음 조건들로 바꾼다고 가정하자.
\begin{enumerate}
\item $F$는 국소 환 달린 공간의 범주 위에서 층 성질을 만족한다.
\item 집합 $I$와 부분함자들의 모임 $F_i \subset F$가 존재하여 다음을
만족한다.
\begin{enumerate}
\item 각 $F_i$는 스킴으로 표현가능하다.
\item 각 $F_i \subset F$는 국소 환 달린 공간의 범주 위에서 열린 몰입으로
표현가능하다.
\item 모임 $(F_i)_{i \in I}$는 국소 환 달린 공간의 범주 위의 함자로서
$F$를 덮는다.
\end{enumerate}
\end{enumerate}
이를 더 상세히 쓰는 일은 독자에게 맡긴다. 그러면 최종적으로 함자 $F$는
국소 환 달린 공간의 범주에서 표현가능하고, 그 표현 대상은 스킴이다.
\end{remark}

\section{스킴의 섬유곱의 존재}
\label{schemes-section-existence-fibre-products}

\noindent
스킴의 범주에서 곱과 섬유곱이 존재하는지는 매우 기본적인 문제이다. 먼저
곱과 섬유곱이 존재함을 추상적으로 증명하고, 다음 절에서 스킴의 섬유곱을
어떻게 합리적으로 생각할 수 있는지 설명한다.

\begin{lemma}
\label{schemes-lemma-fibre-products}
스킴의 범주는 끝 대상, 곱, 섬유곱을 가진다. 다시 말해 스킴의 범주는 유한
극한을 가진다. Categories의 보조정리
\ref{categories-lemma-finite-limits-exist}를 보라.
\end{lemma}

\begin{proof}
이 증명은 건너뛰어도 좋다. 다음 절에서 설명할 섬유곱을 다루는 방법을
배우는 편이 더 중요하다.

\medskip\noindent
보조정리 \ref{schemes-lemma-morphism-into-affine}에 의해 스킴
$\Spec(\mathbf{Z})$는 국소 환 달린 공간의 범주에서 끝 대상이다. 따라서
섬유곱이 존재함을 보이면 충분하다.

\medskip\noindent
$f : X \to S$와 $g : Y \to S$를 스킴의 사상이라 하자. 함자
\begin{eqnarray*}
F : \Sch^{opp} & \longrightarrow & \textit{Sets} \\
T & \longmapsto &
\Mor(T, X) \times_{\Mor(T, S)} \Mor(T, Y)
\end{eqnarray*}
가 표현가능함을 보여야 한다. 보조정리 \ref{schemes-lemma-glue-functors}가 함자
$F$에 적용된다고 주장한다. 이를 증명하면 보조정리가 증명된다.

\medskip\noindent
먼저 $F$가 자리스키 위상에 대한 층 성질을 만족함을 보인다. 실제로 $T$가
스킴이고, $T = \bigcup_{i \in I} U_i$가 열린 덮개이며, 원소들
$\xi_i \in F(U_i)$가
$\xi_i|_{U_i \cap U_j} =  \xi_j|_{U_i \cap U_j}$를 모든 쌍 $i, j$에 대하여
만족한다고 하자. 정의에
의해 $\xi_i$는 쌍 $(a_i, b_i)$에 대응한다. 여기서
$a_i : U_i \to X$와 $b_i : U_i \to Y$는 $f \circ a_i = g \circ b_i$를
만족하는 스킴의 사상들이다. 붙이기 조건은
$a_i|_{U_i \cap U_j} =  a_j|_{U_i \cap U_j}$와
$b_i|_{U_i \cap U_j} =  b_j|_{U_i \cap U_j}$가 성립한다는 것이다. 따라서
사상들 $a_i$를 붙여 국소 환 달린 공간의 사상(즉, 스킴의 사상)
$a : T \to X$를 얻고, 마찬가지로 $b : T \to Y$를 얻는다(예를 들어 보조정리
\ref{schemes-lemma-glue}의 사상 성질을 보라). 또한 열린 덮개의 각 원소 위에서
합성들 $f \circ a$와 $g \circ b$가 일치한다. 따라서
$f \circ a = g \circ b$이고 쌍 $(a, b)$는 $F(T)$의 원소를 정의하며, 이는
쌍 $(a_i, b_i)$로 각 $U_i$ 위에서 제한된다. 층 조건을 확인하였다.

\medskip\noindent
다음으로 부분함자들의 족을 구성한다. 아핀 열린집합들로 이루어진 열린 덮개
$S = \bigcup\nolimits_{i \in I} U_i$를 택하자. 모든 $i \in I$에 대하여 아핀
열린집합들로 이루어진 열린 덮개
$f^{-1}(U_i) = \bigcup\nolimits_{j \in J_i} V_j$와
$g^{-1}(U_i) = \bigcup\nolimits_{k \in K_i} W_k$를 택하자.
$X = \bigcup_{i \in I} \bigcup_{j \in J_i} V_j$는 열린 덮개이고 $Y$도
마찬가지임에 유의하자. 임의의 $i \in I$와 각 쌍
$(j, k) \in J_i \times K_i$에 대하여 가환 도식
$$
\xymatrix{
    & W_k \ar[d] \ar[rd] &   \\
V_j \ar[rd] \ar[r] & U_i \ar[rd] & Y \ar[d] \\
    & X \ar[r]  & S
}
$$
을 가지며, 여기서 모든 빗금 화살표는 열린 몰입이다. 그러한 삼중항에
대하여 함자
\begin{eqnarray*}
F_{i, j, k} : \Sch^{opp} & \longrightarrow & \textit{Sets} \\
T & \longmapsto &
\Mor(T, V_j) \times_{\Mor(T, U_i)} \Mor(T, W_k).
\end{eqnarray*}
를 얻는다. (위의 큰 가환 도식에서 오는) 명백한 함자 변환
$F_{i, j, k} \to F$가 존재하고 이는 단사이므로, $F_{i, j, k}$를 $F$의
부분함자로 생각할 수 있다.

\medskip\noindent
보조정리 \ref{schemes-lemma-glue-functors}의 조건 (2)(a)를 확인하자. 이는 보조정리
\ref{schemes-lemma-fibre-product-affine-schemes}에서 바로 따른다. (여기서는 아핀
스킴의 범주에서의 섬유곱이 국소 환 달린 공간의 전체 범주에서도 섬유곱임을
사용한다.)

\medskip\noindent
보조정리 \ref{schemes-lemma-glue-functors}의 조건 (2)(b)를 확인하자. $T$를 스킴이라
하고 $\xi \in F(T)$라 하자. 다시 말해 $\xi = (a, b)$이고, 여기서
$a : T \to X$와 $b : T \to Y$는 $f \circ a = g \circ b$를 만족하는 스킴의
사상들이다. $V_{i, j, k} = a^{-1}(V_j) \cap b^{-1}(W_k)$로 놓자. 임의의
추가 사상 $h : T' \to T$에 대하여 $h^*\xi = (a \circ h, b \circ h)$이다.
따라서 $h^*\xi \in F_{i, j, k}(T')$일 필요충분조건은
$a(h(T')) \subset V_j$ 및 $b(h(T')) \subset W_k$인 것이다. 다시 말해
$h(T') \subset V_{i, j, k}$일 필요충분이다. 이로써 조건 (2)(b)가
증명되었다.

\medskip\noindent
보조정리 \ref{schemes-lemma-glue-functors}의 조건 (2)(c)를 확인하자. $T$를 스킴이라
하고 위와 같이 $\xi = (a, b) \in F(T)$라 하자. 위와 같이
$V_{i, j, k} = a^{-1}(V_j) \cap b^{-1}(W_k)$로 놓자. 조건 (2)(c)는 바로
$T = \bigcup V_{i, j, k}$라는 뜻이고, 이는 명백하다. 따라서 보조정리가
증명되었고 섬유곱이 존재한다.
\end{proof}

\begin{remark}
\label{schemes-remark-fibre-product-schemes-locally-ringed}
주의 \ref{schemes-remark-representable-locally-ringed}를 사용하면 스킴의 사상들의
섬유곱이 국소 환 달린 공간의 범주에서 존재하고 스킴임을 보일 수 있다.
\end{remark}

\section{스킴의 섬유곱}
\label{schemes-section-fibre-products}

\noindent
스킴의 섬유곱이 존재함을 이미 보였지만, 여기서는 일반적인 정의를
복습한다.

\begin{definition}
\label{schemes-definition-fibre-product}
스킴의 사상 $f : X \to S$와 $g : Y \to S$가 주어졌을 때 {\it 섬유곱}은
스킴 $X \times_S Y$와 사영 사상들 $p : X \times_S Y \to X$ 및
$q : X \times_S Y \to Y$로 이루어지며 다음 가환 도식 안에 놓인다.
$$
\xymatrix{
X \times_S Y \ar[r]_q \ar[d]_p & Y \ar[d]^g \\
X \ar[r]^f & S
}
$$
그리고 이러한 종류의 모든 도식 가운데 보편이다. Categories의 정의
\ref{categories-definition-fibre-products}를 보라.
\end{definition}

\noindent
다시 말해 스킴의 사상들의 임의의 실선 가환 도식
$$
\xymatrix{
T \ar[rrrd] \ar@{-->}[rrd] \ar[rrdd]
&
&
\\
&
&
X \times_S Y \ar[d] \ar[r]
&
Y \ar[d]
\\
&
&
X \ar[r]
&
S
}
$$
이 주어지면 도식을 가환하게 하는 유일한 점선 화살표가 존재한다. 이제
섬유곱을 어떻게 생각해야 하는지 알려 주는 몇 가지 보조정리를 증명하겠다.

\begin{lemma}
\label{schemes-lemma-fibre-product-affines}
$f : X \to S$와 $g : Y \to S$를 같은 공역을 가지는 스킴의 사상이라 하자.
$X, Y, S$가 모두 아핀이면 $X \times_S Y$도 아핀이다.
\end{lemma}

\begin{proof}
$X = \Spec(A)$, $Y = \Spec(B)$, $S = \Spec(R)$이라고 가정하자. 보조정리
\ref{schemes-lemma-fibre-product-affine-schemes}에 의해 아핀 스킴
$\Spec(A \otimes_R B)$는 국소 환 달린 공간의 범주에서 섬유곱
$X \times_S Y$이다. 따라서 이는 당연히 스킴의 범주에서도 섬유곱이다.
\end{proof}

\begin{lemma}
\label{schemes-lemma-open-fibre-product}
$f : X \to S$와 $g : Y \to S$를 같은 공역을 가지는 스킴의 사상이라 하자.
$X \times_S Y$, $p$, $q$를 섬유곱이라 하자. $U \subset S$,
$V \subset X$, $W \subset Y$가 열린 부분스킴들이고
$f(V) \subset U$ 및 $g(W) \subset U$라고 가정하자. 그러면 표준 사상
$V \times_U W \to X \times_S Y$는 열린 몰입이며, $V \times_U W$를
$p^{-1}(V) \cap q^{-1}(W)$와 동일시한다.
\end{lemma}

\begin{proof}
$T$를 스킴이라 하자. $a : T \to V$와 $b : T \to W$가 스킴의 사상들이고,
$f \circ a = g \circ b$가 $U$로 가는 사상으로서 성립한다고 가정하자. 그러면 이들은
$S$로 가는 사상으로서도 일치한다. 섬유곱의 보편 성질에 의해 유일한 사상
$T \to X \times_S Y$를 얻는다. 물론 이 사상의 상은 열린집합
$p^{-1}(V) \cap q^{-1}(W)$에 포함된다. 따라서
$p^{-1}(V) \cap q^{-1}(W)$는 $V$와 $W$의, $U$ 위의 섬유곱이다. 섬유곱의
유일성에서 결과가 따른다. Categories의 절
\ref{categories-section-fibre-products}를 보라.
\end{proof}

\noindent
특히 보조정리의 상황에서 $V \times_U W = V \times_S W$임을 알 수 있다.
또한 $U, V, W$가 모두 아핀이면 $V \times_U W$가 아핀임을 안다. 물론
$X \times_S Y$를 이러한 아핀 열린집합들 $V \times_U W$로 덮을 수 있다.
이를 보조정리로 표현하자.

\begin{lemma}
\label{schemes-lemma-affine-covering-fibre-product}
\begin{slogan}
섬유곱의 직접 구성: 스킴의 섬유곱의 아핀 열린 덮개는 각 조각의 양립하는
아핀 열린 덮개들로부터 조립할 수 있다.
\end{slogan}
$f : X \to S$와 $g : Y \to S$를 같은 공역을 가지는 스킴의 사상이라 하자.
$S = \bigcup U_i$를 $S$의 임의의 아핀 열린 덮개라 하자. 각 $i \in I$에
대하여 $f^{-1}(U_i) = \bigcup_{j \in J_i} V_j$를 $f^{-1}(U_i)$의 아핀 열린
덮개라 하고, $g^{-1}(U_i) = \bigcup_{k \in K_i} W_k$를 $g^{-1}(U_i)$의 아핀
열린 덮개라 하자. 그러면
$$
X \times_S Y =
\bigcup\nolimits_{i \in I}
\bigcup\nolimits_{j \in J_i, \ k \in K_i}
V_j \times_{U_i} W_k
$$
는 $X \times_S Y$의 아핀 열린 덮개이다.
\end{lemma}

\begin{proof}
보조정리 앞의 논의를 보라.
\end{proof}

\noindent
다시 말해 앞의 보조정리를 사용하여 아핀 스킴들을 붙임으로써 섬유곱을 직접
구성할 수도 있었다. (물론 이는 어쨌든 보조정리
\ref{schemes-lemma-fibre-products}의 증명에서 우리가 한 일과 정확히 같다.) 다음은
스킴의 섬유곱의 점들의 집합을 기술하는 방법이다.

\begin{lemma}
\label{schemes-lemma-points-fibre-product}
$f : X \to S$와 $g : Y \to S$를 같은 공역을 가지는 스킴의 사상이라 하자.
점 $z$들, 즉 $X \times_S Y$의 점들은 사중항
$$
(x, y, s, \mathfrak p)
$$
들과 전단사로 대응한다. 여기서 $x \in X$, $y \in Y$, $s \in S$는
$f(x) = s$, $g(y) = s$를 만족하는 점들이고, $\mathfrak p$는 환
$\kappa(x) \otimes_{\kappa(s)} \kappa(y)$의 소 아이디얼이다. $z$의
잉여체는 소 아이디얼 $\mathfrak p$의 잉여체에 대응한다.
\end{lemma}

\begin{proof}
$z$를 $X \times_S Y$의 점이라 하고 위와 같은 사중항을 구성하자. 보조정리
\ref{schemes-lemma-characterize-points}에 의해 $z$를 사상
$\Spec(\kappa(z)) \to X \times_S Y$로 생각할 수 있음을 상기하자. 이 사상은
사상들 $a : \Spec(\kappa(z)) \to X$와 $b : \Spec(\kappa(z)) \to Y$에
대응하며, 이들은 $f \circ a = g \circ b$를 만족한다. 같은 보조정리를 다시
적용하면 점들 $x \in X$, $y \in Y$를 얻고 이들은 같은 점 $s \in S$ 위에
놓이며, 체 사상들 $\kappa(x) \to \kappa(z)$,
$\kappa(y) \to \kappa(z)$를 얻으며, 합성들
$\kappa(s) \to \kappa(x) \to \kappa(z)$와
$\kappa(s) \to \kappa(y) \to \kappa(z)$는 같다. 다시 말해 환 사상
$\kappa(x) \otimes_{\kappa(s)} \kappa(y) \to \kappa(z)$를 얻는다.
$\mathfrak p$를 이 사상의 핵이라 하자.

\medskip\noindent
역으로 사중항 $(x, y, s, \mathfrak p)$가 주어지면 실선 가환 도식
$$
\xymatrix{
X \times_S Y
\ar@/_/[dddr] \ar@/^/[rrrd]
& & & \\
&
\Spec(\kappa(x) \otimes_{\kappa(s)} \kappa(y)/\mathfrak p)
\ar[r] \ar[d] \ar@{-->}[lu]
&
\Spec(\kappa(y)) \ar[d] \ar[r] &
Y \ar[dd] \\
&
\Spec(\kappa(x)) \ar[r] \ar[d] &
\Spec(\kappa(s)) \ar[rd] &
\\
&
X \ar[rr] &
&
S
}
$$
을 얻는다. 절 \ref{schemes-section-points}의 논의를 보라. 따라서 점선 화살표를
얻는다. 이에 대응하는 점 $z$, 즉 $X \times_S Y$의 이 점은
$\Spec(\kappa(x) \otimes_{\kappa(s)} \kappa(y)/\mathfrak p)$의 일반점의
상이다. 두 구성이 서로 역이라는 확인은 생략한다.
\end{proof}

\begin{lemma}
\label{schemes-lemma-fibre-product-immersion}
$f : X \to S$와 $g : Y \to S$를 같은 공역을 가지는 스킴의 사상이라 하자.
\begin{enumerate}
\item $f : X \to S$가 닫힌 몰입이면 $X \times_S Y \to Y$는 닫힌
몰입이다. 또한 $X \to S$가 아이디얼의 준연접층
$\mathcal{I} \subset \mathcal{O}_S$에 대응하면 $X \times_S Y \to Y$는
아이디얼의 층 $\Im(g^*\mathcal{I} \to \mathcal{O}_Y)$에 대응한다.
\item $f : X \to S$가 열린 몰입이면 $X \times_S Y \to Y$는 열린
몰입이다.
\item $f : X \to S$가 몰입이면 $X \times_S Y \to Y$는 몰입이다.
\end{enumerate}
\end{lemma}

\begin{proof}
$X \to S$가 아이디얼의 준연접층 $\mathcal{I} \subset \mathcal{O}_S$에
대응하는 닫힌 몰입이라고 가정하자. 보조정리
\ref{schemes-lemma-restrict-map-to-closed}에 의해 닫힌 부분공간 $Z \subset Y$, 즉
아이디얼의 층 $\Im(g^*\mathcal{I} \to \mathcal{O}_Y)$로 정의되는 공간은
국소 환 달린 공간의 범주에서 섬유곱이다. 보조정리
\ref{schemes-lemma-closed-subspace-scheme}에 의해 $Z$는 스킴이다. 따라서
$Z = X \times_S Y$이고 첫 번째 주장이 따른다. 두 번째 주장은 예를 들어
보조정리 \ref{schemes-lemma-open-fibre-product}에서 따른다. 세 번째 주장은 앞의
두 주장을 결합한 것이다.
\end{proof}

\begin{definition}
\label{schemes-definition-inverse-image-closed-subscheme}
$f : X \to Y$를 스킴의 사상이라 하고 $Z \subset Y$를 $Y$의 닫힌
부분스킴이라 하자. {\it 역상 $f^{-1}(Z)$}, 즉 닫힌 부분스킴 $Z$의 역상은
닫힌 부분스킴 $Z \times_Y X$이며, 이는 $X$의 부분스킴이다. 위의 보조정리
\ref{schemes-lemma-fibre-product-immersion}을 보라.
\end{definition}

\noindent
때때로 국소 닫힌 부분스킴과 열린 부분스킴에 대해서도 이 용어를 사용할 수
있다.

\section{대수기하학에서의 밑변환}
\label{schemes-section-base-change}

\noindent
스킴의 언어를 도입하는 한 가지 동기는 특정한 체 위에서 대수다양체를
정의한다는 말에 매우 정확한 의미를 부여하기 때문이다. 예를 들어 대수다양체
$X$가 $\mathbf{Q}$ 위에 있다는 것은 유한형이고 분리이며 기약이고 축약인
$X \to \Spec(\mathbf{Q})$와 같은 말이다(Varieties의 정의
\ref{varieties-definition-variety})\footnote{물론 대수기하학자들은 여전히
$X$가 $\mathbf{Q}$ 위에서 기하적으로 기약이어야 하는지를 두고 논쟁한다.}.
어쨌든 더 일반적인 생각은 흔히 $S$로 나타내는 주어진 {\it 바탕 스킴} 위의
스킴들을 다루는 것이다. ``$X$를 $S$ 위의 스킴이라 하자''라는 말은 단순히
$X$에 사상 $X \to S$가 주어져 있다는 뜻으로 사용한다. 도식에서는
{\it 구조 사상} $X \to S$를 $X$에서 $S$로 내려가는 화살표로 그리려고
하겠다. 흔히 $X$의 내부 기하보다 $S$에 대한 $X$의 상대적 성질에 더
관심을 둔다. 예를 들어 $X \to S$의 섬유, 밑변환 뒤의 $X$ 등에 관해 알고
싶다.

\medskip\noindent
관례적으로 쓰이는 용어 몇 가지를 도입한다. 물론 이 용어는 범주의 한 대상
위의 대상들의 범주를 생각하는 특수한 경우일 뿐이다. Categories의 예
\ref{categories-example-category-over-X}를 보라.

\begin{definition}
\label{schemes-definition-base-change}
$S$를 스킴이라 하자.
\begin{enumerate}
\item $X$가 {\it $S$ 위의 스킴}이라는 말은 $X$에 스킴의 사상 $X \to S$가
주어져 있다는 뜻이다. 사상 $X \to S$를 때로 {\it 구조 사상}이라고 한다.
\item $R$이 환이면 $X$가 {\it $R$ 위의 스킴}이라고 말하며, 이는 $X$가
$\Spec(R)$ 위의 스킴이라는 뜻이다.
\item {\it 스킴의 사상 $f : X \to Y$가 $S$ 위의 스킴의 사상이라는 것}은
합성 $X \to Y \to S$, 즉 $f$와 $Y$의 구조 사상의 합성이 $X$의 구조
사상과 같다는 뜻이다.
\item $\Mor_S(X, Y)$로 $X$에서 $Y$로 가는 $S$ 위의 모든 사상의 집합을
나타낸다.
\item $X$를 $S$ 위의 스킴이라 하고 $S' \to S$를 스킴의 사상이라 하자.
$X$의 {\it 밑변환}은 스킴 $X_{S'} = S' \times_S X$이며, 이는 $S'$ 위의
스킴이다.
\item $f : X \to Y$를 $S$ 위의 스킴의 사상이라 하고 $S' \to S$를 스킴의
사상이라 하자. $f$의 {\it 밑변환}은 유도된 사상
$f' : X_{S'} \to Y_{S'}$이다(즉, 사상
$\text{id}_{S'} \times_{\text{id}_S} f$이다).
\item $R$을 환이라 하고 $X$를 $R$ 위의 스킴이라 하며 $R \to R'$을 환
사상이라 하자. {\it 밑변환} $X_{R'}$는 스킴
$\Spec(R') \times_{\Spec(R)} X$이며, 이는 $R'$ 위의 스킴이다.
\end{enumerate}
\end{definition}

\noindent
다음은 전형적인 결과이다.

\begin{lemma}
\label{schemes-lemma-base-change-immersion}
$S$를 스킴이라 하자. $f : X \to Y$를 $S$ 위의 스킴의 몰입(각각 닫힌 몰입,
각각 열린 몰입)이라 하자. 그러면 $f$의 임의의 밑변환은 몰입(각각 닫힌
몰입, 각각 열린 몰입)이다.
\end{lemma}

\begin{proof}
$f$의 밑변환을 사상 $S' \to S$를 통해 다음 가환 도식의 왼쪽 위 세로 화살표로
생각할 수 있다.
$$
\xymatrix{
X_{S'} \ar[r] \ar[d] & X \ar[d] \ar@/^4ex/[dd] \\
Y_{S'} \ar[r] \ar[d] & Y \ar[d] \\
S' \ar[r] & S
}
$$
이 도식은 $X_{S'} \cong Y_{S'} \times_Y X$를 함의하며, 보조정리
\ref{schemes-lemma-fibre-product-immersion}에서 결과가 따른다.
\end{proof}

\noindent
사실 이런 종류의 결과는 매우 전형적이므로 이를 표현하는 용어가 따로 있다.

\begin{definition}
\label{schemes-definition-preserved-by-base-change}
성질과 밑변환.
\begin{enumerate}
\item $\mathcal{P}$를 바탕 위의 스킴의 성질이라 하자. $\mathcal{P}$가
{\it 임의의 밑변환 아래 보존된다}, 또는 간단히 $\mathcal{P}$가
{\it 밑변환 아래 보존된다}는 것은 $X/S$가 $\mathcal{P}$를 가질 때마다
임의의 밑변환 $X_{S'}/S'$도 $\mathcal{P}$를 가진다는 뜻이다.
\item $\mathcal{P}$를 바탕 위의 스킴의 사상들의 성질이라 하자.
$\mathcal{P}$가 {\it 임의의 밑변환 아래 보존된다}, 또는 간단히
{\it 밑변환 아래 보존된다}는 것은 $f : X \to Y$가 $S$
위에서 $\mathcal{P}$를 가질 때마다 임의의 밑변환
$f' : X_{S'} \to Y_{S'}$가 $S'$ 위에서 $\mathcal{P}$를 가진다는 뜻이다.
\end{enumerate}
\end{definition}

\noindent
이제 ``닫힌 몰입임''은 임의의 밑변환 아래 보존된다고 말할 수 있다.

\begin{definition}
\label{schemes-definition-fibre}
$f : X \to S$를 스킴의 사상이라 하고 $s \in S$를 점이라 하자.
{\it 스킴론적 섬유 $X_s$, 곧 $f$의 $s$ 위 스킴론적 섬유}, 또는
간단히 {\it $f$의 $s$ 위 섬유}란 다음 섬유곱 도식에 놓이는 스킴이다.
$$
\xymatrix{
X_s = \Spec(\kappa(s)) \times_S X \ar[r] \ar[d] &
X \ar[d] \\
\Spec(\kappa(s)) \ar[r] &
S
}
$$
섬유 $X_s$는 언제나 $\kappa(s)$ 위의 스킴으로 생각한다.
\end{definition}

\begin{lemma}
\label{schemes-lemma-fibre-topological}
$f : X \to S$를 스킴의 사상이라 하자. 도식들
$$
\vcenter{
\xymatrix{
X_s \ar[r] \ar[d] & X \ar[d] \\
\Spec(\kappa(s)) \ar[r] & S
}
}
\quad\text{and}\quad
\vcenter{
\xymatrix{
\Spec(\mathcal{O}_{S, s}) \times_S X \ar[r] \ar[d] & X \ar[d] \\
\Spec(\mathcal{O}_{S, s}) \ar[r] & S
}
}
$$
을 생각하자. 두 경우 모두 이 도식들은 위상공간의 섬유곱 도식을 유도한다.
특히 위쪽 가로 화살표는 그 상으로의 위상동형이다.
\end{lemma}

\begin{proof}
아핀 열린집합 $U \subset S$로서 $s$를 포함하는 것을 택하자. 아래쪽 가로 사상들은
$U$를 통해 인수분해된다. 예를 들어 보조정리
\ref{schemes-lemma-morphism-from-spec-local-ring}을 보라. 따라서 $S$가 아핀이라고
가정할 수 있다. $X$도 아핀이면 결과는 Algebra의 주의
\ref{algebra-remark-fundamental-diagram}에서 따른다. 일반적인 경우에는
$X$를 아핀 열린집합들로 덮어 결과를 얻는다.
\end{proof}

\begin{lemma}
\label{schemes-lemma-local-ring-fibre}
$f : X \to S$를 스킴의 사상이라 하고 $x \in X$의 상을 $s \in S$라 하자.
$x$를 $X_s$의 점으로 생각하면(보조정리 \ref{schemes-lemma-fibre-topological}를 보라)
$$
\mathcal{O}_{X_s, x} \cong
\mathcal{O}_{X, x}/\mathfrak m_s\mathcal{O}_{X, x} \cong
\mathcal{O}_{X, x} \otimes_{\mathcal{O}_{S, s}} \kappa(s)
$$
이다.
\end{lemma}

\begin{proof}
보조정리 \ref{schemes-lemma-fibre-topological}의 증명에서와 같이 이는 $X$와 $S$가
아핀인 경우로 환원된다. 아핀인 경우 이 명제는 Algebra의 주의
\ref{algebra-remark-local-ring-fibre}로 옮겨진다.
\end{proof}

\section{준콤팩트 사상}
\label{schemes-section-quasi-compact}

\noindent
스킴의 기저 위상공간이 준콤팩트이면 그 스킴을 {\it 준콤팩트}라고 한다.
이에 대응하는 상대적 개념은 다음과 같이 정의한다.

\begin{definition}
\label{schemes-definition-quasi-compact}
스킴의 사상이 유도하는 위상공간 사이의 사상이 준콤팩트이면 그 사상을
{\it 준콤팩트}라고 한다. Topology의 정의
\ref{topology-definition-quasi-compact}를 보라.
\end{definition}

\begin{lemma}
\label{schemes-lemma-quasi-compact-affine}
$f : X \to S$를 스킴의 사상이라 하자. 다음 조건들은 서로 동치이다.
\begin{enumerate}
\item $f : X \to S$가 준콤팩트이다.
\item 모든 아핀 열린집합의 역상이 준콤팩트이다.
\item 어떤 아핀 열린 덮개 $S = \bigcup_{i \in I} U_i$가 존재하여
$f^{-1}(U_i)$가 모든 $i$에 대해 준콤팩트이다.
\end{enumerate}
\end{lemma}

\begin{proof}
(3)에서와 같은 덮개 $S = \bigcup_{i \in I} U_i$가 주어졌다고 하자. 먼저
$U \subset S$를 임의의 아핀 열린집합이라 하자. 각 $u \in U$에 대해 지표
$i(u) \in I$를 택하여 $u \in U_{i(u)}$가 되게 할 수 있다. 표준
열린집합들은 $U_{i(u)}$의 위상의 기저를 이루므로
$W_u \subset U \cap U_{i(u)}$를 택하여 $U_{i(u)}$에서 표준 열린집합이
되게 할 수 있다. 콤팩트성에 의해
유한 개의 점 $u_1, \ldots, u_n \in U$를 택하여
$U = \bigcup_{j = 1}^n W_{u_j}$로 쓸 수 있다. 각 $j$에 대해
$f^{-1}U_{i(u_j)} = \bigcup_{k \in K_j} V_{jk}$를 아핀 열린집합들의 유한
합집합으로 쓰자. $W_{u_j} \subset U_{i(u_j)}$가 표준 열린집합이므로
$f^{-1}(W_{u_j}) \cap V_{jk}$는 $V_{jk}$의 표준 열린집합이다. Algebra의
보조정리 \ref{algebra-lemma-spec-functorial}을 보라. 따라서
$f^{-1}(W_{u_j}) \cap V_{jk}$는 아핀이고, 그러므로 $f^{-1}(W_{u_j})$는
아핀 열린집합들의 유한 합집합이다. 이로써 임의의 아핀 열린집합의 역상이
아핀 열린집합들의 유한 합집합임을 보였다.

\medskip\noindent
이제 모든 아핀 열린집합의 역상이 아핀 열린집합들의 유한 합집합이라고
가정하자. $K \subset S$를 임의의 준콤팩트 열린집합이라 하자. $S$의
위상에는 아핀 열린집합들로 이루어진 기저가 있으므로 $K$는 아핀
열린집합들의 유한 합집합이다. 따라서 $K$의 역상도 아핀 열린집합들의 유한
합집합이다. 그러므로 $f$는 준콤팩트이다.

\medskip\noindent
마지막으로 $f$가 준콤팩트라고 가정하자. 이 경우 앞 단락의 논증에 의해
임의의 아핀 열린집합의 역상은 아핀 열린집합들의 유한 합집합이다.
\end{proof}

\begin{lemma}
\label{schemes-lemma-quasi-compact-preserved-base-change}
준콤팩트임은 바탕 위의 스킴 사상의 성질이며, 임의의 밑변환 아래
보존된다.
\end{lemma}

\begin{proof}
생략한다.
\end{proof}

\begin{lemma}
\label{schemes-lemma-composition-quasi-compact}
준콤팩트 사상들의 합성은 준콤팩트이다.
\end{lemma}

\begin{proof}
정의와 Topology의 보조정리
\ref{topology-lemma-composition-quasi-compact}에서 따른다.
\end{proof}

\begin{lemma}
\label{schemes-lemma-closed-immersion-quasi-compact}
닫힌 몰입은 준콤팩트이다.
\end{lemma}

\begin{proof}
정의와 Topology의 보조정리
\ref{topology-lemma-closed-in-quasi-compact}에서 따른다.
\end{proof}

\begin{example}
\label{schemes-example-open-immersion-not-quasi-compact}
일반적으로 열린 몰입은 준콤팩트가 아니다. 표준적인 예는 열린 부분공간
$U \subset X$이다. 여기서 $X = \Spec(k[x_1, x_2, x_3, \ldots])$이고,
$U$는 $X \setminus \{0\}$이며, $0$은 $X$의 점이고 극대 아이디얼
$(x_1, x_2, x_3, \ldots)$에 대응한다.
\end{example}

\begin{lemma}
\label{schemes-lemma-image-quasi-compact-closed}
$f : X \to S$를 스킴의 준콤팩트 사상이라 하자. 다음 조건들은 서로
동치이다.
\begin{enumerate}
\item $f(X) \subset S$가 닫혀 있다.
\item $f(X) \subset S$가 특수화 아래 안정적이다.
\end{enumerate}
\end{lemma}

\begin{proof}
Topology의 보조정리 \ref{topology-lemma-open-closed-specialization}에 의해
(1) $\Rightarrow$ (2)이다. (2)를 가정하자. $U \subset S$를 아핀
열린집합이라 하자. $f(X) \cap U$가 닫혀 있음을 보이면 충분하다.
$U \cap f(X)$가 $U$ 안의 특수화 아래 안정적이므로 $S$가 아핀인 경우로
환원되었다. $f$가 준콤팩트이므로 $X = f^{-1}(S)$가 준콤팩트임을 얻는데,
이는 $S$가 아핀이기 때문이다. 따라서
$X = \bigcup_{i = 1}^n U_i$로 쓸 수 있으며, 여기서
$U_i \subset X$는 아핀 열린집합이다. $S = \Spec(R)$라 하고
$U_i = \Spec(A_i)$라 하자. 여기서 어떤 $R$-대수를 $A_i$로 나타냈다. 그러면
$f(X) = \Im(\Spec(A_1 \times \ldots \times A_n)
\to \Spec(R))$이다. 따라서 보조정리는 Algebra의 보조정리
\ref{algebra-lemma-image-stable-specialization-closed}에서 따른다.
\end{proof}

\begin{lemma}
\label{schemes-lemma-quasi-compact-closed}
$f : X \to S$를 스킴의 준콤팩트 사상이라 하자. 그러면 $f$가 닫힌
사상일 필요충분조건은 $f$를 따라 특수화를 올릴 수 있는 것이다. Topology의
정의 \ref{topology-definition-lift-specializations}를 보라.
\end{lemma}

\begin{proof}
Topology의 보조정리
\ref{topology-lemma-closed-open-map-specialization}에 따르면 $f$가 닫힌
사상이면 $f$를 따라 특수화를 올릴 수 있다. 역으로 $f$를 따라 특수화를
올릴 수 있다고 하자. $Z \subset X$를 닫힌 부분집합이라 하자. 정의
\ref{schemes-definition-reduced-induced-scheme}에 따라 $Z$를 유도된 축약 스킴
구조를 갖춘 스킴으로 생각할 수 있다. $Z \subset X$가 닫혀 있으므로
$f$를 $Z$에 제한한 사상도 여전히 준콤팩트이다. 또한 $Z \to S$를 따라서도
특수화를 올릴 수 있다. Topology의 보조정리
\ref{topology-lemma-lift-specialization-composition}을 보라. 따라서 $f(X)$가
닫혀 있음을, $f$를 따라 특수화를 올릴 수 있을 때 보이면 충분하다. 특히
$f(X)$는 특수화 아래 안정적이다. Topology의 보조정리
\ref{topology-lemma-lift-specializations-images}를 보라. 그러므로 보조정리
\ref{schemes-lemma-image-quasi-compact-closed}에 의해 $f(X)$는 닫혀 있다.
\end{proof}

\section{보편 닫힘의 값매김 판정법}
\label{schemes-section-valuative-criterion-universal-closedness}

\noindent
Topology의 절 \ref{topology-section-proper}에서는 모든 섬유가 준콤팩트인
위상공간의 닫힌 사상으로서의 고유 사상, 또는 모든 밑변환이 닫힌 사상인
사상에 관해 논의한다. 대수기하학에서 이에 대응하는 개념은 다음과 같다.

\begin{definition}
\label{schemes-definition-universally-closed}
스킴의 사상 $f : X \to S$의 모든 밑변환
$f' : X_{S'} \to S'$이 닫힌 사상이면 그 사상을 {\it 보편 닫힌 사상}이라고
한다.
\end{definition}

\noindent
실제로 ``보편적으로''라는 형용사는 흔히 이런 방식으로 쓰인다. 즉, 사상의
성질 $\mathcal{P}$가 주어졌을 때 ``$X \to S$가 보편적으로
$\mathcal{P}$이다''라는 말은 모든 밑변환 $X_{S'} \to S'$가
$\mathcal{P}$를 가진다는 것과 동치이다.

\medskip\noindent
보편 닫힌 사상의 성질에 대한 더 자세한 논의는 Morphisms의 절
\ref{morphisms-section-proper}을 보라. 이 절에서는 보편 닫힌 사상과 값매김
판정법의 존재 부분을 만족하는 사상의 관계만을 논의한다.

\begin{lemma}
\label{schemes-lemma-specializations-lift}
$f : X \to S$를 스킴의 사상이라 하자.
\begin{enumerate}
\item $f$가 보편 닫힌 사상이면 $f$의 임의의 밑변환을 따라 특수화를 올릴
수 있다. Topology의 정의 \ref{topology-definition-lift-specializations}를
보라.
\item $f$가 준콤팩트이고 $f$의 임의의 밑변환을 따라 특수화를 올릴 수
있으면 $f$는 보편 닫힌 사상이다.
\end{enumerate}
\end{lemma}

\begin{proof}
(1)은 Topology의 보조정리
\ref{topology-lemma-closed-open-map-specialization}의 직접적인 귀결이다.
(2)는 보조정리 \ref{schemes-lemma-quasi-compact-closed}와
\ref{schemes-lemma-quasi-compact-preserved-base-change}에서 따른다.
\end{proof}

\begin{definition}
\label{schemes-definition-valuative-criterion}
$f : X \to S$를 스킴의 사상이라 하자. $f$가 {\it 값매김 판정법의 존재
부분을 만족한다}는 것은 임의의 가환 실선 도식
$$
\xymatrix{
\Spec(K) \ar[r] \ar[d] & X \ar[d] \\
\Spec(A) \ar[r] \ar@{-->}[ru] & S
}
$$
이 주어졌을 때, 여기서 $A$는 분수체가 $K$인 값매김환이며 점선 화살표가
항상 존재한다는 뜻이다. 위와 같은 임의의 도식에서 점선 화살표가 많아야
하나라면(물론 존재는 요구하지 않는다) $f$가 {\it 값매김 판정법의 유일성
부분을 만족한다}고 한다.
\end{definition}

\noindent
{\it 값매김환}은 그 분수체 안의 지배 관계에 관하여 극대인 국소 정역이다.
Algebra의 정의 \ref{algebra-definition-valuation-ring}을 보라. 따라서
값매김환의 스펙트럼에는 유일한 일반점 $\eta$와 유일한 닫힌점 $0$이 있으며,
물론 특수화 $\eta \leadsto 0$이 있다. 값매김환이 중요한 이유는 임의의
스킴에서 점들의 임의의 특수화가 어떤 값매김환의 스펙트럼에서 오는 어떤
사상 아래 $\eta \leadsto 0$의 상이기 때문이다. 정확한 결과는 다음과 같다.

\begin{lemma}
\label{schemes-lemma-points-specialize}
\begin{slogan}
특수화는 값매김환으로 목격된다.
\end{slogan}
$S$를 스킴이라 하자. $s' \leadsto s$를 $S$의 점들의 특수화라 하자. 그러면
\begin{enumerate}
\item 값매김환 $A$와 사상 $f : \Spec(A) \to S$가 존재하여
일반점 $\eta$, 즉 $\Spec(A)$의 일반점은 $s'$로, 특수점은 $s$로 사상된다.
\item 체 확대 $K/\kappa(s')$가 주어지면, 확대
$\kappa(\eta)/\kappa(s')$가 주어진 확대와 동형이 되도록 정할 수 있으며,
이 확대는 $f$가 유도한다.
\end{enumerate}
\end{lemma}

\begin{proof}
$s' \leadsto s$를 $S$ 안의 특수화라 하고, $K/\kappa(s')$를 체의 확대라
하자. 보조정리 \ref{schemes-lemma-specialize-points}와 보조정리
\ref{schemes-lemma-characterize-points} 뒤의 논의에 의해 환 사상들
$\mathcal{O}_{S, s} \to \kappa(s') \to K$를 얻는다. $A \subset K$를
분수체가 $K$이고 $\mathcal{O}_{S, s} \to K$의 상을 지배하는 임의의
값매김환이라 하자. Algebra의 보조정리 \ref{algebra-lemma-dominate}를 보라.
환 사상 $\mathcal{O}_{S, s} \to A$는 사상 $f : \Spec(A) \to S$를 유도한다.
보조정리 \ref{schemes-lemma-morphism-from-spec-local-ring}을 보라. 이 사상은 구성에
의해 필요한 모든 성질을 가진다.
\end{proof}

\begin{lemma}
\label{schemes-lemma-lift-specializations-valuative}
$f : X \to S$를 스킴의 사상이라 하자. 다음 조건들은 서로 동치이다.
\begin{enumerate}
\item $f$의 임의의 밑변환을 따라 특수화를 올릴 수 있다.
\item 사상 $f$는 값매김 판정법의 존재 부분을 만족한다.
\end{enumerate}
\end{lemma}

\begin{proof}
(1)이 성립한다고 가정하자. 정의 \ref{schemes-definition-valuative-criterion}에서와
같은 실선 도식이 주어졌다고 하자. 점선 화살표를 찾기 위해 $X \to S$를
$X_{\Spec(A)} \to \Spec(A)$로 바꾸어도 된다. 가정은 밑변환 아래
안정적이기 때문이다. 따라서 $S = \Spec(A)$라고 가정할 수 있다.
$x' \in X$를 $\Spec(K) \to X$의 상이라 하면
$\kappa(x') \subset K$이다. 보조정리 \ref{schemes-lemma-characterize-points}를 보라.
가정에 의해 특수화 $x' \leadsto x$가 $X$ 안에 존재하여 $x$는
$S = \Spec(A)$의 닫힌점으로 사상된다. 국소환 사상
$A \to \mathcal{O}_{X, x}$와 환 사상
$\mathcal{O}_{X, x} \to \kappa(x')$를 얻는다. 보조정리
\ref{schemes-lemma-specialize-points}와 보조정리 \ref{schemes-lemma-characterize-points}
뒤의 논의를 보라. 합성
$A \to \mathcal{O}_{X, x} \to \kappa(x') \to K$는 주어진 단사
$A \to K$이다. $A \to \mathcal{O}_{X, x}$가 국소 사상이므로
$\mathcal{O}_{X, x} \to K$의 상은 $A$를 지배하고, 따라서 Algebra의 정의
\ref{algebra-definition-valuation-ring}에 의해 $A$와 같다. 그러므로 환 사상
$\mathcal{O}_{X, x} \to A$와 이에 따른 사상 $\Spec(A) \to X$를 얻는다
(보조정리 \ref{schemes-lemma-morphism-from-spec-local-ring}과 그 뒤의 논의를 보라).
이로써 (2)를 증명했다.

\medskip\noindent
역으로 (2)가 성립한다고 가정하자. 다음 가환 도식을 생각하면 밑변환
$X_{S'} \to S'$이 $f$의 임의의 밑변환일 때 그것도 값매김 판정법의 존재
부분을 만족한다는 것이
즉시 따른다.
$$
\xymatrix{
\Spec(K) \ar[r] \ar[d] & X_{S'} \ar[r] \ar[d] & X \ar[d] \\
\Spec(A) \ar[r] \ar@{-->}[ru] \ar@{-->}[rru] & S' \ar[r] & S
}
$$
즉, 더 수평에 가까운 점선 화살표에서 섬유곱의 정의에 의해 다른 점선
화살표를 얻는다. 따라서 $f$를 따라 특수화를 올릴 수 있음을 보이면
충분하다. $s' \leadsto s$를 $S$ 안의 특수화라 하고, $x' \in X$를 $s'$ 위의
점이라 하자. 보조정리 \ref{schemes-lemma-points-specialize}를 $s' \leadsto s$와 체
확대 $K = \kappa(x')/\kappa(s')$에 적용한다. 가환 도식
$$
\xymatrix{
\Spec(K) \ar[rr] \ar[d] & & X \ar[d] \\
\Spec(A) \ar[r] \ar@{-->}[rru] &
\Spec(\mathcal{O}_{S, s}) \ar[r] & S
}
$$
을 얻고, 조건 (2)에 의해 점선 화살표를 얻는다. $x$를 $\Spec(A)$의 닫힌점의
$X$ 안에서의 상이라 하면 이것이 문제의 해이다. 즉, $x$는 $x'$의
특수화이며 $s$로 사상된다.
\end{proof}

\begin{proposition}[보편 닫힘의 값매김 판정법]
\label{schemes-proposition-characterize-universally-closed}
$f$를 스킴의 준콤팩트 사상이라 하자. 그러면 $f$가 보편 닫힌 사상일
필요충분조건은 $f$가 값매김 판정법의 존재 부분을 만족하는 것이다.
\end{proposition}

\begin{proof}
이는 위의 보조정리 \ref{schemes-lemma-specializations-lift}와
\ref{schemes-lemma-lift-specializations-valuative}의 형식적인 귀결이다.
\end{proof}

\begin{example}
\label{schemes-example-projective-line-universally-closed}
$k$를 체라 하자. 구조 사상 $p : \mathbf{P}^1_k \to \Spec(k)$, 곧
$k$ 위의 사영직선의 구조 사상을 생각하자. 예
\ref{schemes-example-projective-line}을 보라. 위의 값매김 판정법을 사용하여 $p$가
보편 닫힌 사상임을 보이자. 구성에 의해 $\mathbf{P}^1_k$는 두 아핀
열린집합으로 덮이므로 $p$는 준콤팩트이다. 가환 도식
$$
\xymatrix{
\Spec(K) \ar[r]_\xi \ar[d] & \mathbf{P}^1_k \ar[d] \\
\Spec(A) \ar[r]^\varphi & \Spec(k)
}
$$
이 주어졌다고 하자. 여기서 $A$는 값매김환이고 $K$는 그 분수체이다.
$\mathbf{P}^1_k$는 $\Spec(k[x])$와 $\Spec(k[y])$를 붙여서 얻으며,
$D(x)$와 $D(y)$를 $x = y^{-1}$을 통해 붙인다(또는 더 대칭적으로
$xy = 1$을 통해 붙인다). 이 도식에 대각선으로 들어맞는 사상
$\Spec(A) \to \mathbf{P}^1_k$가 존재함을 보이기 위해, 대칭성에 의해
$\xi$의 상이 열린집합 $\Spec(k[x])$에 들어간다고 가정할 수 있다. 이로부터
다음 환의 가환 도식을 얻는다.
$$
\xymatrix{
K & k[x] \ar[l]^{\xi^\sharp} \\
A \ar[u] & k \ar[u] \ar[l]_{\varphi^\sharp}
}
$$
Algebra의 보조정리 \ref{algebra-lemma-valuation-ring-x-or-x-inverse}에 의해
$\xi^\sharp(x) \in A$이거나 $\xi^\sharp(x)^{-1} \in A$이다. 첫 번째
경우에는 환 사상
$$
k[x] \to A,
\ \lambda \mapsto \varphi^\sharp(\lambda),
\  x \mapsto \xi^\sharp(x)
$$
을 얻으며, 이는 위의 환의 도식에 들어맞으므로 원하는 결과를 얻는다.
두 번째 경우에는 환 사상
$$
k[y] \to A,
\ \lambda \mapsto \varphi^\sharp(\lambda),
\ y \mapsto \xi^\sharp(x)^{-1}.
$$
을 얻는다. 이는 사상
$\Spec(A) \to \Spec(k[y]) \to \mathbf{P}^1_k$를 주며, 이 예의 처음 가환
도식에 대각선으로 들어맞는다(확인은 생략한다).
\end{example}

\section{분리 공리}
\label{schemes-section-separation-axioms}

\noindent
위상공간 $X$가 하우스도르프일 필요충분조건은 대각선
$\Delta \subset X \times X$가 닫힌 부분집합인 것이다. 대수기하학에서
이에 대응하는 것은 스킴 $X$가 바탕 스킴 $S$ 위에 주어졌을 때 대각 사상
$$
\Delta_{X/S} : X \longrightarrow X \times_S X.
$$
을 생각하는 것이다. 이것은
$\text{pr}_1 \circ \Delta_{X/S} = \text{id}_X$ 및
$\text{pr}_2 \circ \Delta_{X/S} = \text{id}_X$를 만족하는 유일한 스킴
사상이다(섬유곱이 있는 임의의 범주에서 존재한다).

\begin{lemma}
\label{schemes-lemma-diagonal-affines-closed}
아핀 스킴 사이의 사상에 대한 대각 사상은 닫혀 있다.
\end{lemma}

\begin{proof}
사상 $\Spec(S) \to \Spec(R)$에 대응하는 대각 사상은 환 준동형
$S \otimes_R S \to S$, $a \otimes b \mapsto ab$에 대응하는 스펙트럼
사상이다. 이 준동형은 명백히 전사이므로 $S \cong S \otimes_R S/J$이며,
여기서 $J \subset S \otimes_R S$는 어떤 아이디얼이다. 따라서
예 \ref{schemes-example-closed-immersion-affines}에 의해 $\Delta$는 닫힌 몰입이다.
\end{proof}

\begin{lemma}
\label{schemes-lemma-diagonal-immersion}
\begin{slogan}
상대 스킴의 대각 사상은 몰입이다.
\end{slogan}
$X$를 $S$ 위의 스킴이라 하자. 대각 사상 $\Delta_{X/S}$는 몰입이다.
\end{lemma}

\begin{proof}
$V \subset X$가 아핀 열린집합이고 그 상이 아핀 열린집합 $U \subset S$에
들어가면 $V \times_U V$는 $X \times_S X$의 아핀 열린집합임을 상기하자.
보조정리 \ref{schemes-lemma-fibre-product-affines}와
\ref{schemes-lemma-open-fibre-product}를 보라. $W$를 $X \times_S X$의 열린
부분스킴으로서 이 아핀 열린집합들 $V \times_U V$의 합집합이라 하자.
보조정리 \ref{schemes-lemma-closed-local-target}에 의해 각 사상
$\Delta_{X/S}^{-1}(V \times_U V) \to V \times_U V$가 닫힌 몰입임을 보이면
충분하다. $V = \Delta_{X/S}^{-1}(V \times_U V)$이므로 확인할 것은
$\Delta_{V/U}$가 닫힌 몰입이라는 것뿐이며, 이는 보조정리
\ref{schemes-lemma-diagonal-affines-closed}이다.
\end{proof}

\begin{definition}
\label{schemes-definition-separated}
$f : X \to S$를 스킴의 사상이라 하자.
\begin{enumerate}
\item $f$의 대각 사상 $\Delta_{X/S}$가 닫힌 몰입이면 이를 {\it 분리 사상}이라고
한다.
\item $f$의 대각 사상 $\Delta_{X/S}$가 준콤팩트 사상이면 이를 {\it 준분리
사상}이라고 한다.
\item 스킴 $Y$를 사상 $Y \to \Spec(\mathbf{Z})$가 분리일 때 {\it 분리
스킴}이라고 한다.
\item 스킴 $Y$를 사상 $Y \to \Spec(\mathbf{Z})$가 준분리일 때
{\it 준분리 스킴}이라고 한다.
\end{enumerate}
\end{definition}

\noindent
보조정리 \ref{schemes-lemma-diagonal-immersion}과
\ref{schemes-lemma-immersion-when-closed}에 의해 $\Delta_{X/S}$가 닫힌 몰입일
필요충분조건은 $\Delta_{X/S}(X) \subset X \times_S X$가 닫힌 부분집합인
것이다. 또한 보조정리 \ref{schemes-lemma-closed-immersion-quasi-compact}에 의해
분리 사상은 준분리 사상이다. 준분리 사상을 도입하는 이유는 대수다양체를
연구할 때(특히 모듈라이, 대수적 스택 등을 다룰 때) 비분리 사상이 자연스럽게
나타나기 때문이다. 그러나 대개 이들은 여전히 준분리이다.

\begin{example}
\label{schemes-example-not-quasi-separated}
다음은 준분리가 아닌 사상의 예이다.
$X = X_1 \cup X_2 \to S = \Spec(k)$라 하고,
$X_1 = X_2 = \Spec(k[t_1, t_2, t_3, \ldots])$의 두 복사본을
$\{0\} = \{(t_1, t_2, t_3, \ldots)\}$의 여집합을 따라 붙였다고 하자
(예 \ref{schemes-example-affine-space-zero-doubled}에서처럼 붙인다). 이 경우 아핀
스킴 $X_1 \times_S X_2$의 대각 사상 $\Delta_{X/S}$에 의한 역상은
$\Spec(k[t_1, t_2, t_3, \ldots]) \setminus \{0\}$이며, 이는 준콤팩트가
아니다.
\end{example}

\begin{lemma}
\label{schemes-lemma-where-are-they-equal}
$X$, $Y$를 $S$ 위의 스킴이라 하자. $a, b : X \to Y$를 $S$ 위의 스킴
사상이라 하자. 가장 큰 국소 닫힌 부분스킴 $Z \subset X$가 존재하여
$a|_Z = b|_Z$를 만족한다. 실제로 $Z$는 $(a, b)$의 등화자이다. 또한 $Y$가
$S$ 위에서 분리이면 $Z$는 닫힌 부분스킴이다.
\end{lemma}

\begin{proof}
범주론적 이유에 의해 $(a, b)$의 등화자는 다음 도식의 섬유곱 $Z$이다.
$$
\xymatrix{
Z = Y \times_{(Y \times_S Y)} X \ar[r] \ar[d] &
 X \ar[d]^{(a , b)} \\
Y \ar[r]^-{\Delta_{Y/S}} & Y \times_S Y
}
$$
따라서 보조정리 \ref{schemes-lemma-base-change-immersion},
\ref{schemes-lemma-diagonal-immersion} 및 정의 \ref{schemes-definition-separated}에서 결과가
따른다.
\end{proof}

\begin{lemma}
\label{schemes-lemma-characterize-quasi-separated}
$f : X \to S$를 스킴의 사상이라 하자. 다음 조건들은 서로 동치이다.
\begin{enumerate}
\item 사상 $f$가 준분리이다.
\item 아핀 열린집합의 모든 쌍 $U, V \subset X$ 중 $S$의 공통된 어떤 아핀
열린집합으로 사상되는 것에 대해 교집합 $U \cap V$는 $X$의 아핀 열린집합들의 유한
합집합이다.
\item 아핀 열린 덮개 $S = \bigcup_{i \in I} U_i$가 존재하고, 각 $i$마다
아핀 열린 덮개 $f^{-1}U_i = \bigcup_{j \in I_i} V_j$가 존재하여, 각 $i$와
각 쌍 $j, j' \in I_i$에 대해 교집합 $V_j \cap V_{j'}$가 $X$의 아핀
열린집합들의 유한 합집합이다.
\end{enumerate}
\end{lemma}

\begin{proof}
(3)이 (1)을 함의함을 보이자. 보조정리
\ref{schemes-lemma-affine-covering-fibre-product}에 의해 덮개
$X \times_S X = \bigcup_i \bigcup_{j, j'} V_j \times_{U_i} V_{j'}$는
$X \times_S X$의 아핀 열린 덮개이다. 또한
$\Delta_{X/S}^{-1}(V_j \times_{U_i} V_{j'}) = V_j \cap V_{j'}$이다. 따라서
보조정리 \ref{schemes-lemma-quasi-compact-affine}에서 함의가 따른다.

\medskip\noindent
(1) $\Rightarrow$ (2)는 (2)의 가정 아래 섬유곱 $U \times_S V$가
$X \times_S X$의 아핀 열린집합이라는 사실에서 따른다. (2)
$\Rightarrow$ (3)은 자명하다.
\end{proof}

\begin{lemma}
\label{schemes-lemma-characterize-separated}
$f : X \to S$를 스킴의 사상이라 하자.
\begin{enumerate}
\item $f$가 분리이면 아핀 열린집합의 모든 쌍 $(U, V)$ 중 $X$에 놓이고 $S$의
공통된 어떤 아핀 열린집합으로 사상되는 것에 대해 다음이 성립한다.
\begin{enumerate}
\item 교집합 $U \cap V$는 아핀이다.
\item 환 준동형
$\mathcal{O}_X(U) \otimes_{\mathbf{Z}} \mathcal{O}_X(V)
\to \mathcal{O}_X(U \cap V)$
은 전사이다.
\end{enumerate}
\item 임의의 두 점 $x_1, x_2 \in X$가 공통된 한 점 $s \in S$ 위에 놓이고,
각각 아핀 열린집합 $x_1 \in U$, $x_2 \in V$에 들어가며 이 두 열린집합이
$S$의 공통된 어떤 아핀 열린집합으로 사상되고 (a), (b)가
성립한다면 $f$는 분리이다.
\end{enumerate}
\end{lemma}

\begin{proof}
$f$가 분리라고 가정하자. $(U, V)$가 (1)에서와 같은 쌍이라고 하자.
$W = \Spec(R)$를 $S$의 아핀 열린집합으로서 $f(U)$와 $f(V)$를 모두
포함하는 것이라 하자. $U = \Spec(A)$ 및 $V = \Spec(B)$라 쓰되,
$R$-대수 $A$와 $B$를 사용한다. 보조정리
\ref{schemes-lemma-open-fibre-product}에 의해
$U \times_S V = U \times_W V = \Spec(A \otimes_R B)$는
$X \times_S X$의 아핀 열린집합이다. 따라서 보조정리
\ref{schemes-lemma-closed-subspace-scheme}에 의해
$\Delta^{-1}(U \times_S V) \to U \times_S V$를
$\Spec((A \otimes_R B)/J) \to \Spec(A \otimes_R B)$와 동일시할 수
있으며, 여기서 $J \subset A \otimes_R B$는 어떤 아이디얼이다.
그러므로 $U \cap V = \Delta^{-1}(U \times_S V)$는 아핀이다. 또한
$A \otimes_{\mathbf{Z}} B \to (A \otimes_R B)/J$가 전사이므로 (1)(b)가
성립한다.

\medskip\noindent
(2)에 서술된 가정이 성립한다고 하자. 아핀 열린집합 $U \times_S V$들은,
(2)에서와 같은 쌍 $(U, V)$에 대해 택했을 때, 명백히 $X \times_S X$의 아핀 열린
덮개를 이룬다(예를 들어 보조정리
\ref{schemes-lemma-affine-covering-fibre-product}을 보라). 따라서 각 사상
$U \cap V = \Delta_{X/S}^{-1}(U \times_S V) \to U \times_S V$가 닫힌
몰입임을 보이면 충분하다. 보조정리 \ref{schemes-lemma-closed-local-target}을 보라.
가정 (a)에 의해 $U \cap V = \Spec(C)$이며, 여기서 $C$는 어떤 환이다.
$W = \Spec(R)$를 $S$의 아핀 열린집합으로서 $U$와 $V$가 모두 사상되는
것으로 택하고
$U = \Spec(A)$, $V = \Spec(B)$라 쓰면 가정 (b)는 합성
$$
A \otimes_{\mathbf{Z}} B \to A \otimes_R B \to C
$$
이 전사라는 뜻이다. 따라서 $A \otimes_R B \to C$는 전사이고
$\Spec(C) \to \Spec(A \otimes_R B)$가 닫힌 몰입이라고 결론 내린다.
\end{proof}

\begin{example}
\label{schemes-example-projective-line-separated}
$k$를 체라 하자. 구조 사상 $p : \mathbf{P}^1_k \to \Spec(k)$, 곧
$k$ 위의 사영직선의 구조 사상을 생각하자. 예
\ref{schemes-example-projective-line}을 보라. 위 보조정리를 사용하여 $p$가 분리임을
보이자. 구성에 의해 $\mathbf{P}^1_k$는 두 아핀 열린집합
$U = \Spec(k[x])$와 $V = \Spec(k[y])$로 덮이며 그 교집합은
$U \cap V = \Spec(k[x, y]/(xy - 1))$이다(자명한 표기법을 사용했다).
따라서 보조정리 \ref{schemes-lemma-characterize-separated}의 조건 (2)(a)와 (2)(b)가
아핀 열린집합의 쌍 $(U, U)$, $(U, V)$, $(V, U)$ 및 $(V, V)$에 대해
성립하는지만 확인하면 충분하다. 쌍 $(U, U)$와 $(V, V)$에 대해서는
자명하다. 쌍 $(U, V)$에 대해서는 $U \cap V$가 아핀임을 보이는 것과 환
준동형
$$
k[x] \otimes_{\mathbf{Z}} k[y] \longrightarrow k[x, y]/(xy - 1)
$$
이 전사임을 보이는 것으로 귀착된다. 첫 번째는 참이고, 두 번째도 우변의
임의의 원소를 $x$의 다항식과 $y$의 다항식의 합으로 쓸 수 있으므로
명백하다.
\end{example}

\begin{lemma}
\label{schemes-lemma-fibre-product-after-map}
$f : X \to T$와 $g : Y \to T$를 같은 공역을 갖는 스킴 사상이라 하고,
$h : T \to S$를 스킴의 사상이라 하자. 그러면 유도된 사상
$i : X \times_T Y \to X \times_S Y$는 몰입이다. $T \to S$가 분리이면
$i$는 닫힌 몰입이다. $T \to S$가 준분리이면 $i$는 준콤팩트 사상이다.
\end{lemma}

\begin{proof}
일반 범주론에 의해 다음 도식
$$
\xymatrix{
X \times_T Y \ar[r] \ar[d] & X \times_S Y \ar[d] \\
T \ar[r]^{\Delta_{T/S}} \ar[r] & T \times_S T
}
$$
은 섬유곱 도식이다. 보조정리 \ref{schemes-lemma-diagonal-immersion},
\ref{schemes-lemma-fibre-product-immersion} 및
\ref{schemes-lemma-quasi-compact-preserved-base-change}에서 결과가 따른다.
\end{proof}

\begin{lemma}
\label{schemes-lemma-semi-diagonal}
$g : X \to Y$를 $S$ 위의 스킴 사상이라 하자. 사상
$i : X \to X \times_S Y$는 몰입이다. $Y$가 $S$ 위에서 분리이면 이는 닫힌
몰입이다. $Y$가 $S$ 위에서 준분리이면 이는 준콤팩트이다.
\end{lemma}

\begin{proof}
이는 보조정리 \ref{schemes-lemma-fibre-product-after-map}을 사상
$X = X \times_Y Y \to X \times_S Y$에 적용한 특수한 경우이다.
\end{proof}

\begin{lemma}
\label{schemes-lemma-section-immersion}
$f : X \to S$를 스킴의 사상이라 하자. $s : S \to X$를 $f$의 단면이라
하자(식으로는 $f \circ s = \text{id}_S$). 그러면 $s$는 몰입이다. $f$가
분리이면 $s$는 닫힌 몰입이다. $f$가 준분리이면 $s$는 준콤팩트이다.
\end{lemma}

\begin{proof}
이는 보조정리 \ref{schemes-lemma-semi-diagonal}을 $g =s$에 적용한 특수한 경우이며,
그 사상은 $i = s : S \to S \times_S X$이다.
\end{proof}

\begin{lemma}
\label{schemes-lemma-separated-permanence}
영속성 성질.
\begin{enumerate}
\item 분리 사상들의 합성은 분리이다.
\item 준분리 사상들의 합성은 준분리이다.
\item 분리 사상의 밑변환은 분리이다.
\item 준분리 사상의 밑변환은 준분리이다.
\item 분리 사상들의 (섬유)곱은 분리이다.
\item 준분리 사상들의 (섬유)곱은 준분리이다.
\end{enumerate}
\end{lemma}

\begin{proof}
$X \to Y \to Z$를 사상들이라 하자. $X \to Y$와 $Y \to Z$가 분리라고
가정하자. 합성
$$
X \to X \times_Y X \to X \times_Z X
$$
은 닫혀 있다. 첫 번째 사상은 가정에 의해, 두 번째 사상은 보조정리
\ref{schemes-lemma-fibre-product-after-map}에 의해 닫혀 있기 때문이다. 같은 논증은
``준분리''에도 성립한다(같은 참조를 사용한다).

\medskip\noindent
$f : X \to Y$를 바탕 $S$ 위의 스킴 사상이라 하자. $S' \to S$를 스킴의
사상이라 하고, $f' : X_{S'} \to Y_{S'}$를 $f$의 밑변환이라 하자. 그러면
$f'$의 대각 사상은 사상
$$
\Delta_{f'} :
X_{S'} = S' \times_S X
\longrightarrow
X_{S'} \times_{Y_{S'}} X_{S'} = S' \times _S (X \times_Y X)
$$
이며, 이것은 $\Delta_f$의 밑변환임을 쉽게 알 수 있다. 따라서 (3)과 (4)는
닫힌 몰입과 준콤팩트 사상이 임의의 밑변환 아래 보존된다는 사실에서 따른다
(보조정리 \ref{schemes-lemma-fibre-product-immersion}과
\ref{schemes-lemma-quasi-compact-preserved-base-change}).

\medskip\noindent
$f : X \to Y$와 $g : U \to V$를 바탕 $S$ 위의 스킴 사상이라 하자. 그러면
$f \times g$는 $X \times_S U \to X \times_S V$($g$의 밑변환)와
$X \times_S V \to Y \times_S V$($f$의 밑변환)의 합성이다. 따라서 (5)와
(6)은 (1)--(4)에서 따른다.
\end{proof}

\begin{lemma}
\label{schemes-lemma-compose-after-separated}
\begin{slogan}
분리 사상과 준분리 사상에는 소거법칙이 성립한다.
\end{slogan}
$f : X \to Y$와 $g : Y \to Z$를 스킴의 사상이라 하자. $g \circ f$가
분리이면 $f$도 분리이다. $g \circ f$가 준분리이면 $f$도 준분리이다.
\end{lemma}

\begin{proof}
$g \circ f$가 분리라고 가정하자. 대각 사상의 분해
$X \to X \times_Y X \to X \times_Z X$를 생각하되, 이는 $g \circ f$의
대각 사상이다. 보조정리
\ref{schemes-lemma-fibre-product-after-map}에 의해 뒤의 사상은 몰입이다. 가정에 의해
$X$의 $X \times_Z X$ 안의 상은 닫혀 있다. 따라서 $X \times_Y X$ 안에서도
닫혀 있다. 그러므로 보조정리 \ref{schemes-lemma-immersion-when-closed}에 의해
$X \to X \times_Y X$가 닫힌 몰입이다.

\medskip\noindent
$g \circ f$가 준분리라고 가정하자. 아핀 열린집합 $V \subset Y$로서 그
상이 $Z$의 어떤 아핀 열린집합에 들어가는 것을 택하자. 아핀 열린집합
$U_1, U_2 \subset X$로서 $V$로 사상되는 것들을 택하자. 그러면
$U_1 \cap U_2$는 아핀 열린집합들의 유한 합집합이다. 실제로 $U_1, U_2$는
$Z$의 공통된 아핀 열린집합으로 사상된다.
$Y$를 $V$와 같은 아핀 열린집합들로 덮을 수 있으므로 보조정리
\ref{schemes-lemma-characterize-quasi-separated}에서 결과가 따른다.
\end{proof}

\begin{lemma}
\label{schemes-lemma-quasi-compact-permanence}
$f : X \to Y$와 $g : Y \to Z$를 스킴의 사상이라 하자. $g \circ f$가
준콤팩트이고 $g$가 준분리이면 $f$는 준콤팩트이다.
\end{lemma}

\begin{proof}
$f$가 합성 $(1, f) : X \to X \times_Z Y \to Y$와 같기 때문이다. 첫 번째
사상은 준분리 사상 $X \times_Z Y \to X$의 단면이므로 보조정리
\ref{schemes-lemma-section-immersion}에 의해 준콤팩트이다(이는 $g$의 밑변환이다.
보조정리 \ref{schemes-lemma-separated-permanence}을 보라). 두 번째 사상은
$g \circ f$의 밑변환이므로 준콤팩트이다. 보조정리
\ref{schemes-lemma-quasi-compact-preserved-base-change}를 보라. 준콤팩트 사상들의
합성은 준콤팩트이다. 보조정리 \ref{schemes-lemma-composition-quasi-compact}를 보라.
\end{proof}

\begin{lemma}
\label{schemes-lemma-affine-separated}
아핀 스킴은 분리이다. 아핀 스킴에서 다른 스킴으로 가는 사상은 분리이다.
\end{lemma}

\begin{proof}
$U = \Spec(A)$를 아핀 스킴이라 하자. 그러면 보조정리
\ref{schemes-lemma-diagonal-affines-closed}에 의해 $U \to \Spec(\mathbf{Z})$의 대각
사상은 닫혀 있다. 따라서 정의 \ref{schemes-definition-separated}에 의해 $U$는
분리이다. $U \to X$가 스킴의 사상이면 보조정리
\ref{schemes-lemma-compose-after-separated}을 사상들
$U \to X \to \Spec(\mathbf{Z})$에 적용하여 $U \to X$가 분리임을 얻는다.
\end{proof}

\noindent
상들이 아핀 열린집합에 포함되는 아핀 열린집합의 쌍만을 생각해야만 그
교집합이 아핀이라고 결론 내릴 수 있는지 궁금했을지도 모른다. 흔히 그럴
필요가 없다!

\begin{lemma}
\label{schemes-lemma-curiosity}
$f : X \to S$를 사상이라 하자. $f$가 분리이고 $S$가 분리 스킴이라고
가정하자. $U \subset X$와 $V \subset X$가 아핀 열린집합이면
$U \cap V$는 아핀이다(그리고 $U \times V$의 닫힌 부분스킴이다).
\end{lemma}

\begin{proof}
이 경우 보조정리 \ref{schemes-lemma-separated-permanence}에 의해 $X$는 분리이다.
따라서 $U \cap V$가 아핀임은 보조정리
\ref{schemes-lemma-characterize-separated}을 사상 $X \to \Spec(\mathbf{Z})$에
적용하여 얻는다.
\end{proof}

\noindent
한편 다음 예는 아핀 스킴의 상이 어떤 아핀 열린집합에 포함된다고 기대할
수 없음을 보여 준다.

\begin{example}
\label{schemes-example-image-affine-projective}
예 \ref{schemes-example-not-affine}의 아핀이 아닌 스킴
$U = \Spec(k[x, y]) \setminus \{(x, y)\}$를 생각하자. 한편 스킴
$$
\mathbf{GL}_{2, k} = \Spec(k[a, b, c, d, 1/ad - bc]).
$$
을 생각하자. 사상 $\mathbf{GL}_{2, k} \to U$가 있으며, 이는 환 준동형
$x \mapsto a$, $y \mapsto b$에 대응한다. 이 사상이 전사임은 쉽게 알 수 있으므로
그 상은 $U$의 어떤 아핀 열린집합에도 포함되지 않는다. 실제로 아핀 스킴
$\mathbf{GL}_{2, k}$는 $\mathbf{P}^1_k$ 위로도 전사이며,
$\mathbf{P}^1_k$는 {\it 어떠한} 아핀 스킴으로도 몰입조차 갖지 않는다.
\end{example}

\begin{remark}
\label{schemes-remark-quasi-compact-and-quasi-separated}
$P$를 다음 $4$개 성질 중 하나라 하자. 즉 ``준분리'', ``분리'',
``준콤팩트이고 준분리'', 또는 ``준콤팩트이고 분리''이다. 그러면 다음이
성립한다.
\begin{enumerate}
\item 임의의 아핀 스킴은 $P$를 가진다.
\item 스킴의 사상 $f : X \to Y$가 주어지고 $Y$가 $P$를 가진다고 하자.
그러면 $f$가 $P$를 가질 필요충분조건은 $X$가 $P$를 가지는 것이다.
\item $X \to Y$와 $Z \to Y$가 스킴의 사상이고 $X$, $Y$, $Z$가 $P$를
가지면 $X \times_Y Z$도 $P$를 가진다.
\end{enumerate}
(1)은 명백하다. $f : X \to Y$를 (2)에서와 같다고 하자. $f$가 $P$를
가지면 $X$도 그러하다. 이는 보조정리 \ref{schemes-lemma-separated-permanence}과
\ref{schemes-lemma-composition-quasi-compact}을 $X \to Y \to \Spec(\mathbf{Z})$에
적용하여 얻는다. $X$가 $P$를 가지면 $f$도 그러하다. 이는 보조정리
\ref{schemes-lemma-compose-after-separated}과 \ref{schemes-lemma-quasi-compact-permanence}을
$X \to Y \to \Spec(\mathbf{Z})$에 적용하여 얻는다.
$X \to Y$와 $Z \to Y$를 (3)에서와 같다고 하자. 사영
$X \times_Y Z \to Z$는 $P$를 가진다. 실제로 화살표 $X \to Y$가 (2)에
의해 $P$를 가지며 이 사영은 그 밑변환이다. 보조정리 \ref{schemes-lemma-separated-permanence}과
\ref{schemes-lemma-quasi-compact-preserved-base-change}를 보라. 따라서 (2)에 의해
$X \times_Y Z$가 $P$를 가진다.
\end{remark}

\section{분리성의 값매김 판정법}
\label{schemes-section-valuative-separatedness}

\begin{lemma}
\label{schemes-lemma-separated-implies-valuative}
$f : X \to S$를 스킴의 사상이라 하자. $f$가 분리이면 $f$는 값매김 판정법의
유일성 부분을 만족한다.
\end{lemma}

\begin{proof}
정의 \ref{schemes-definition-valuative-criterion}에서와 같은 도식이 주어졌다고 하자.
도식에 들어맞는 두 사상 $a, b : \Spec(A) \to X$가 있다고 하자.
$Z \subset \Spec(A)$를 $a$와 $b$의 등화자라 하자. 보조정리
\ref{schemes-lemma-where-are-they-equal}에 의해 이는 $\Spec(A)$의 닫힌 부분스킴이다.
가정에 의해 이는 $\Spec(A)$의 일반점을 포함한다. $A$가 정역이므로
$Z = \Spec(A)$이다. 따라서 원하는 대로 $a = b$이다.
\end{proof}

\begin{lemma}[분리성의 값매김 판정법]
\label{schemes-lemma-valuative-criterion-separatedness}
\begin{reference}
\cite[II Proposition 7.2.3]{EGA}
\end{reference}
$f : X \to S$를 사상이라 하자. 다음을 가정하자.
\begin{enumerate}
\item 사상 $f$가 준분리이다.
\item 사상 $f$가 값매김 판정법의 유일성 부분을 만족한다.
\end{enumerate}
그러면 $f$는 분리이다.
\end{lemma}

\begin{proof}
가정 (1), 명제 \ref{schemes-proposition-characterize-universally-closed}, 보조정리
\ref{schemes-lemma-diagonal-immersion} 및 \ref{schemes-lemma-immersion-when-closed}에 의해
사상 $\Delta_{X/S} : X \to X \times_S X$가 값매김 판정법의 존재 부분을
만족함을 보이면 충분하다. 가환 실선 도식
$$
\xymatrix{
\Spec(K) \ar[r] \ar[d] & X \ar[d] \\
\Spec(A) \ar[r] \ar@{-->}[ru] & X \times_S X
}
$$
이 주어졌다고 하자. 오른쪽 아래 화살표는 두 사상
$a, b : \Spec(A) \to X$로서 $S$ 위에 놓이는 것들에 대응한다. (2)에 의해
$a = b$이다. 따라서 $a$를
점선 화살표로 사용하면 된다.
\end{proof}

\section{모노사상}
\label{schemes-section-monomorphisms}

\begin{definition}
\label{schemes-definition-monomorphism}
스킴의 사상이 스킴의 범주에서 모노사상이면 이를 {\it 모노사상}이라고 한다.
Categories의 정의 \ref{categories-definition-mono-epi}를 보라.
\end{definition}

\begin{lemma}
\label{schemes-lemma-monomorphism}
\begin{slogan}
스킴 사상이 모노사상일 필요충분조건은 그 대각 사상이 동형인 것이다.
\end{slogan}
$j : X \to Y$를 스킴의 사상이라 하자. $j$가 모노사상일 필요충분조건은
대각 사상 $\Delta_{X/Y} : X \to X \times_Y X$가 동형인 것이다.
\end{lemma}

\begin{proof}
이는 섬유곱이 있는 임의의 범주에서 성립한다.
\end{proof}

\begin{lemma}
\label{schemes-lemma-monomorphism-separated}
스킴의 모노사상은 분리이다.
\end{lemma}

\begin{proof}
동형은 닫힌 몰입이라는 사실과 위의 보조정리 \ref{schemes-lemma-monomorphism}에서
따른다.
\end{proof}

\begin{lemma}
\label{schemes-lemma-composition-monomorphism}
모노사상들의 합성은 모노사상이다.
\end{lemma}

\begin{proof}
이는 임의의 범주에서 성립한다.
\end{proof}

\begin{lemma}
\label{schemes-lemma-base-change-monomorphism}
모노사상의 밑변환은 모노사상이다.
\end{lemma}

\begin{proof}
이는 섬유곱이 있는 임의의 범주에서 성립한다.
\end{proof}

\begin{lemma}
\label{schemes-lemma-injective-points}
$j : X \to Y$를 스킴의 사상이라 하자. $j$가 점에서 단사이면 $j$는
분리이다.
\end{lemma}

\begin{proof}
$z$를 $X \times_Y X$의 점이라 하자. 그러면
$x = \text{pr}_1(z)$와 $\text{pr}_2(z)$는 같은 점이다. $j$가 이 두 점을
같은 점 $y$, 즉 $Y$의 한 점으로 사상하기 때문이다. 아핀 열린 근방
$V \subset Y$를 $y$에 대해, 아핀 열린 근방 $U \subset X$를 $x$에 대해 택하여
$j(U) \subset V$가 되게 할 수 있다. 그러면
$z \in U \times_V U \subset X \times_Y X$이다. 따라서
$X \times_Y X$는 아핀 열린집합들 $U \times_V U$의 합집합이다.
$\Delta_{X/Y}^{-1}(U \times_V U) = U$이고
$U \to U \times_V U$가 닫힌 몰입이므로 $\Delta_{X/Y}$가 닫힌 몰입이라고
결론 내린다(보조정리 \ref{schemes-lemma-diagonal-immersion}의 증명에 있는 논증을
보라).
\end{proof}

\begin{lemma}
\label{schemes-lemma-injective-points-surjective-stalks}
$j : X \to Y$를 스킴의 사상이라 하자. 다음을 가정하자.
\begin{enumerate}
\item $j$가 점에서 단사이다.
\item 임의의 $x \in X$에 대해 환 준동형
$j^\sharp_x : \mathcal{O}_{Y, j(x)} \to \mathcal{O}_{X, x}$가 전사이다.
\end{enumerate}
그러면 $j$는 모노사상이다.
\end{lemma}

\begin{proof}
$a, b : Z \to X$를 $j \circ a  = j \circ b$를 만족하는 두 스킴 사상이라
하자. 그러면 (1)에 의해 기저 위상공간 사이의 사상으로서 $a = b$이다.
임의의 $z \in Z$에 대해
$a^\sharp_z \circ j^\sharp_{a(z)} = b^\sharp_z \circ j^\sharp_{b(z)}$가
준동형 $\mathcal{O}_{Y, j(a(z))} \to \mathcal{O}_{Z, z}$ 사이의 등식으로
성립한다. 준동형 $j^\sharp_x$들의 전사성에 의해
$a^\sharp_z = b^\sharp_z$, $\forall z \in Z$이다. 이는
$a^\sharp = b^\sharp$를 함의한다. 따라서 원하는 대로 스킴 사상으로서
$a = b$라고 결론 내린다.
\end{proof}

\begin{lemma}
\label{schemes-lemma-immersions-monomorphisms}
스킴의 몰입은 모노사상이다. 특히 모든 몰입은 분리이다.
\end{lemma}

\begin{proof}
보조정리 \ref{schemes-lemma-injective-points-surjective-stalks}의 판정 조건을
확인함으로써 이를 알 수 있다. 조금 더 우아하게는 보조정리
\ref{schemes-lemma-restrict-map-to-opens}와
\ref{schemes-lemma-characterize-closed-subspace}에 의해 열린 몰입과 닫힌 몰입이
모노사상이라는 사실을 이용할 수 있다. 따라서 이들의 합성인 임의의 몰입도
모노사상이다.
\end{proof}

\begin{lemma}
\label{schemes-lemma-subscheme-of-separated-scheme}
$f : X \to S$를 분리 사상이라 하자. 임의의 국소 닫힌 부분스킴
$Z \subset X$는 $S$ 위에서 분리이다.
\end{lemma}

\begin{proof}
보조정리 \ref{schemes-lemma-immersions-monomorphisms}와 분리 사상들의 합성이
분리라는 사실(보조정리 \ref{schemes-lemma-separated-permanence})에서 따른다.
\end{proof}

\begin{example}
\label{schemes-example-Q-over-Z}
사상 $\Spec(\mathbf{Q}) \to \Spec(\mathbf{Z})$는 모노사상이다. 이는
$\mathbf{Q} \otimes_{\mathbf{Z}} \mathbf{Q} = \mathbf{Q}$이기 때문이다.
더 일반적으로 임의의 스킴 $S$와 임의의 점 $s \in S$에 대해 표준 사상
$$
\Spec(\mathcal{O}_{S, s}) \longrightarrow S
$$
은 모노사상이다.
\end{example}

\begin{lemma}
\label{schemes-lemma-mono-towards-spec-field}
$k_1, \ldots, k_n$을 체들이라 하자. 스킴의 임의의 모노사상
$X \to \Spec(k_1 \times \ldots \times k_n)$에 대해 부분집합
$I \subset \{1, \ldots, n\}$가 존재하여
$X \cong \Spec(\prod_{i \in I} k_i)$이며, 이는
$\Spec(k_1 \times \ldots \times k_n)$ 위의 스킴으로서의 동형이다. 더 일반적으로
$X = \coprod_{i \in I} \Spec(k_i)$가 체의 스펙트럼들의 서로소 합집합이고
$Y \to X$가 모노사상이면 부분집합 $J \subset I$가 존재하여
$Y = \coprod_{i \in J} \Spec(k_i)$이다.
\end{lemma}

\begin{proof}
먼저 $n = 1$인 경우(또는 $\# I = 1$인 경우)로 환원한다. 이는 열린 닫힌
부분스킴 $\Spec(k_i)$들의 역상을 취하여 할 수 있다. 이 경우 $X$는 점을 하나만 가지므로
아핀이다. 이에 대응하는 대수 문제는 다음과 같다. 대수 준동형
$k \to R$이 $R \otimes_k R \cong R$을 만족하면 $R \cong k$이거나
$R = 0$이다. 이는 차원을 생각하면 성립한다. Algebra의 보조정리
\ref{algebra-lemma-epimorphism-over-field}도 보라.
\end{proof}

\section{준연접 가군에 대한 함자성}
\label{schemes-section-quasi-coherent}

\noindent
$X$를 스킴이라 하자. 범주 $\QCoh(\mathcal{O}_X)$로 Modules의 정의
\ref{modules-definition-quasi-coherent}에서 정의한 준연접
$\mathcal{O}_X$-가군들의 범주를 나타낸다. 절
\ref{schemes-section-quasi-coherent-affine}에서 범주 $\QCoh(\mathcal{O}_X)$가
$X$가 아핀이면 좋은 성질을 많이 가짐을 보았다. 준연접이라는
성질은 $X$ 위에서 국소적이므로, 임의의 스킴 $X$ 위의 준연접 층들의 범주도
이 성질들을 물려받는다. 이를 열거하면 다음과 같다.
\begin{enumerate}
\item $\mathcal{O}_X$-가군의 층 $\mathcal{F}$가 준연접일 필요충분조건은
$\mathcal{F}$를 각 아핀 열린집합 $U = \Spec(R)$에 제한한 것이
$\widetilde M$의 꼴인 것이다. 여기서 어떤 $R$-가군을 $M$으로 나타냈다.
\item $\mathcal{O}_X$-가군의 층 $\mathcal{F}$가 준연접일 필요충분조건은
$\mathcal{F}$를 어떤 아핀 열린 덮개의 각 성분에 제한한 것이 준연접인
것이다.
\item 준연접 층들의 임의의 직합은 준연접이다.
\item 준연접 층들의 임의의 여극한은 준연접이다.
\item
\label{schemes-item-kernel}
준연접 층 사이의 사상의 핵과 여핵은 준연접이다.
\item $\mathcal{O}_X$-가군의 짧은 완전열
$0 \to \mathcal{F}_1 \to \mathcal{F}_2 \to \mathcal{F}_3 \to 0$이 주어졌을
때 셋 중 둘이 준연접이면 나머지 하나도 준연접이다.
\item 스킴의 사상 $f : Y \to X$가 주어졌을 때 준연접
$\mathcal{O}_X$-가군의 당김은 준연접 $\mathcal{O}_Y$-가군이다. Modules의
보조정리 \ref{modules-lemma-pullback-quasi-coherent}를 보라.
\item 두 준연접 $\mathcal{O}_X$-가군의 텐서곱은 준연접이다. Modules의
보조정리 \ref{modules-lemma-tensor-product-permanence}를 보라.
\item 준연접 $\mathcal{O}_X$-가군 $\mathcal{F}$가 주어졌을 때
$\mathcal{F}$ 위의 텐서 대수, 대칭 대수 및 외대수는 준연접이다. Modules의
보조정리 \ref{modules-lemma-whole-tensor-algebra-permanence}를 보라.
\item 두 준연접 $\mathcal{O}_X$-가군 $\mathcal{F}$, $\mathcal{G}$가
주어지고 $\mathcal{F}$가 유한 표시이면 내부 Hom
$\SheafHom_{\mathcal{O}_X}(\mathcal{F}, \mathcal{G})$은 준연접이다.
Modules의 보조정리
\ref{modules-lemma-internal-hom-locally-kernel-direct-sum}과 위의
(\ref{schemes-item-kernel})을 보라.
\end{enumerate}
한편 일반적으로 준연접 가군의 밂이 준연접인 것은 아니다. 다음은 그것이
성립하는 경우이다.

\begin{lemma}
\label{schemes-lemma-push-forward-quasi-coherent}
$f : X \to S$를 스킴의 사상이라 하자. $f$가 준콤팩트이고 준분리이면
$f_*$는 준연접 $\mathcal{O}_X$-가군을 준연접 $\mathcal{O}_S$-가군으로
보낸다.
\end{lemma}

\begin{proof}
문제는 $S$ 위에서 국소적이므로 $S$가 아핀이라고 가정할 수 있다. $X$가
준콤팩트이므로 $X = \bigcup_{i = 1}^n U_i$로 쓸 수 있고, 각 $U_i$는 아핀
열린집합이다. $f$가 준분리이므로
$U_i \cap U_j = \bigcup_{k = 1}^{n_{ij}} U_{ijk}$로 쓸 수 있으며, 여기서
$U_{ijk}$들은 아핀 열린집합이다. 보조정리
\ref{schemes-lemma-characterize-quasi-separated}를 보라. 제한들을
$f_i : U_i \to S$ 및 $f_{ijk} : U_{ijk} \to S$로 나타내되, 이는 $f$의
제한들이다. 열린집합 $V$를 $S$에서 임의로 택하고 층 $\mathcal{F}$를 $X$
위에서 임의로 택하면 다음이 성립한다.
\begin{eqnarray*}
f_*\mathcal{F}(V) & = & \mathcal{F}(f^{-1}V) \\
& = &
\Ker\left(
\bigoplus\nolimits_i \mathcal{F}(f^{-1}V \cap U_i)
\to
\bigoplus\nolimits_{i, j, k} \mathcal{F}(f^{-1}V \cap U_{ijk})\right) \\
& = &
\Ker\left(
\bigoplus\nolimits_i f_{i, *}(\mathcal{F}|_{U_i})(V)
\to
\bigoplus\nolimits_{i, j, k} f_{ijk, *}(\mathcal{F}|_{U_{ijk}})(V)\right) \\
& = &
\Ker\left(
\bigoplus\nolimits_i f_{i, *}(\mathcal{F}|_{U_i})
\to
\bigoplus\nolimits_{i, j, k} f_{ijk, *}(\mathcal{F}|_{U_{ijk}})\right)(V)
\end{eqnarray*}
달리 말해 층들의 완전열
$$
0 \to f_*\mathcal{F}
\to \bigoplus f_{i, *}\mathcal{F}_i
\to \bigoplus f_{ijk, *}\mathcal{F}_{ijk}
$$
이 있다. 여기서 $\mathcal{F}_i, \mathcal{F}_{ijk}$는 $\mathcal{F}$를 해당
열린집합에 제한한 것을 나타낸다. $\mathcal{F}$가 준연접
$\mathcal{O}_X$-가군이면 $\mathcal{F}_i$는 준연접
$\mathcal{O}_{U_i}$-가군이고 $\mathcal{F}_{ijk}$는 준연접
$\mathcal{O}_{U_{ijk}}$-가군이다. 따라서 보조정리
\ref{schemes-lemma-widetilde-pullback}에 의해 완전열의 두 번째 항과 세 번째 항은
준연접 $\mathcal{O}_S$-가군이다. 그러므로 $f_*\mathcal{F}$는 준연접
$\mathcal{O}_S$-가군이다.
\end{proof}

\noindent
이를 사용하면 스킴의 (닫힌) 몰입을 다음과 같이 특징지을 수 있다.

\begin{lemma}
\label{schemes-lemma-characterize-closed-immersions}
$f : X \to Y$를 스킴의 사상이라 하자. 다음을 가정하자.
\begin{enumerate}
\item $f$는 $X$와 $Y$의 어떤 닫힌 부분집합 사이의 위상동형을 유도한다.
\item $f^\sharp : \mathcal{O}_Y \to f_*\mathcal{O}_X$는 전사이다.
\end{enumerate}
그러면 $f$는 스킴의 닫힌 몰입이다.
\end{lemma}

\begin{proof}
(1)과 (2)를 가정하자. (1)에 의해 사상 $f$는 준콤팩트이다(Topology의
보조정리 \ref{topology-lemma-closed-in-quasi-compact}를 보라). 조건 (1)과
(2)는 보조정리 \ref{schemes-lemma-injective-points-surjective-stalks}의 조건 (1)과
(2)를 함의한다. 따라서 $f : X \to Y$는 모노사상이다. 특히 $f$는 분리이다.
보조정리 \ref{schemes-lemma-monomorphism-separated}를 보라. 그러므로 위의 보조정리
\ref{schemes-lemma-push-forward-quasi-coherent}를 적용할 수 있어
$f_*\mathcal{O}_X$가 준연접 $\mathcal{O}_Y$-가군임을 얻는다. 따라서
$\mathcal{O}_Y \to f_*\mathcal{O}_X$의 핵은 보조정리
\ref{schemes-lemma-extension-quasi-coherent}에 의해 준연접이다. 준연접 층은
국소적으로 단면들에 의해 생성되므로(Modules의 정의
\ref{modules-definition-quasi-coherent}를 보라) $f$는 닫힌 몰입이다. 정의
\ref{schemes-definition-closed-immersion-locally-ringed-spaces}를 보라.
\end{proof}

\noindent
이 보조정리를 사용하여 다음 보조정리를 증명할 수 있다.

\begin{lemma}
\label{schemes-lemma-composition-immersion}
스킴의 몰입들의 합성은 몰입이고, 스킴의 닫힌 몰입들의 합성은 닫힌
몰입이며, 스킴의 열린 몰입들의 합성은 열린 몰입이다.
\end{lemma}

\begin{proof}
열린 부분공간의 열린 부분공간도 열린 부분공간이므로 열린 몰입의 경우에는
명백하다.

\medskip\noindent
$a : Z \to Y$와 $b : Y \to X$가 스킴의 닫힌 몰입이라고 하자.
$c = b \circ a$도 닫힌 몰입임을 확인하겠다. 가정에 의해 $a$와 $b$는 닫힌
부분집합 위로의 위상동형이며, 따라서 $c = b \circ a$도 닫힌 부분집합
위로의 위상동형이다. 또한 준동형 $\mathcal{O}_X \to c_*\mathcal{O}_Z$는
전사이다. 이는 전사 준동형들 $\mathcal{O}_X \to b_*\mathcal{O}_Y$와
$b_*\mathcal{O}_Y \to b_*a_*\mathcal{O}_Z$의 합성으로 인수분해되기
때문이다($b_*$가 완전하므로 전사이다. Modules의 보조정리
\ref{modules-lemma-i-star-exact}를 보라). 따라서 위의 보조정리
\ref{schemes-lemma-characterize-closed-immersions}에 의해 $c$는 닫힌 몰입이다.

\medskip\noindent
마지막으로 몰입의 경우를 다룬다. $a : Z \to Y$와 $b : Y \to X$가 스킴의
몰입이라고 하자. 이는 열린 부분스킴 $V \subset Y$와 $U \subset X$가
존재하여 $a(Z) \subset V$, $b(Y) \subset U$이고
$a : Z \to V$와 $b : Y \to U$가 닫힌 몰입이라는 뜻이다. $Y$의 위상은
$U$의 위상에서 유도되므로 열린집합 $U' \subset U$를 찾아
$V = b^{-1}(U')$가 되게 할 수 있다. 그러면
$Z \to V = b^{-1}(U') \to U'$은 닫힌 몰입들의 합성이므로 닫힌 몰입이다.
따라서 $Z \to X$는 몰입이며 원하는 결과를 얻었다.
\end{proof}
